\chapter{分割与场景理解}

\paragraph*{本章阅读锚点}
原始3DGS只负责把几何、颜色和不透明度渲染成图像，并不知道哪个高斯属于哪件物体、哪类语义、哪条语言查询或哪个检测框。分割与场景理解工作的共同增量，是在高斯基元上增加“可渲染的语义属性”，再利用SAM、CLIP、DINO、语言模型、2D检测器或点云网络把2D语义监督提升到3D。读这一章时可以抓住三条锚点路线：Feature 3DGS把任意维语义特征作为高斯属性并可微渲染；Gaussian Grouping和SAGA把SAM掩码蒸馏成3D实例亲和特征；OpenGaussian、LangSplatV2等把开放词汇语言特征压缩、离散或稀疏化，使3DGS可以被文本查询。后续方法大多是在这些锚点上改变特征压缩方式、跨视角一致性、边界处理、交互提示、前馈预测或具体任务接口。

\section{开放词汇场景理解}

开放词汇理解的基本问题是：CLIP类2D特征强在语义泛化，但跨视角不一致且维度高；3DGS强在可实时渲染，但每个高斯存储高维语言特征会拖慢训练和推理。因此，这条路线通常先从多视角图像提取语言或实例特征，再通过量化、码本、稀疏系数、图传播或跨视角聚合，把它们绑定到高斯上。



% METHOD: 2312.16084
\paragraph{LangSplat（语言高斯场锚点）（2023.12）\cite{2312_16084}}
LangSplat 是开放词汇 3DGS 的早期锚点。它面对的基线不是普通语义分割网络，而是 LERF 一类 NeRF 语言场：这类方法把 CLIP 特征蒸馏进连续辐射场，能用文本查询三维区域，但每次查询都要做体渲染；同时，CLIP 是图像/区域级特征，直接把多尺度裁剪特征分配给三维点，会出现“同一点既属于鼻子、熊头、熊”这类尺度歧义。LangSplat 的改变是把语言场从 NeRF 换成显式高斯场，并把语言特征做成可 splat 的高斯属性。

具体流程有三步。第一步仍先训练或载入一个 3DGS 几何与颜色场，然后在训练视图上用 SAM 产生层级掩码，用 CLIP/OpenCLIP 对掩码区域编码，得到边界更清楚的 2D 语言监督。第二步处理存储问题：若每个高斯直接保存 512 维 CLIP 向量，显存和渲染代价会很高；LangSplat 因此为每个场景训练一个 language autoencoder，把高维语义特征 $\mathbf f$ 压缩成低维 latent $\mathbf z=E(\mathbf f)$，查询或监督时再由 $D(\mathbf z)$ 回到 CLIP 空间。第三步把 $\mathbf z_i$ 当作每个高斯的可渲染属性，像颜色一样按 3DGS 的透明度顺序混合：
\[
  \hat{\mathbf z}(\mathbf u)=\sum_i T_i(\mathbf u)\alpha_i(\mathbf u)\mathbf z_i .
\]
训练时，渲染出的 $\hat{\mathbf z}$ 经解码后与对应视图的 CLIP/SAM 区域特征对齐；查询时，把文本嵌入和解码后的语言特征做相似度比较，即可得到开放词汇目标区域。这样，LangSplat 不再为每个查询反复体渲染多尺度 NeRF 特征，而是用 tile-based splatting 直接渲染语言特征图。

这也是后续方法需要先理解 LangSplat 的原因。OpenGaussian 把 LangSplat 的“视图语言场查询”推进到三维点级实例选择；GOI、GAGS、CAGS 等方法分别改阈值决策、粒度一致性和高斯图传播；LangSplatV2 则保留 SAM/CLIP 监督，但用稀疏字典替代 autoencoder/MLP 解码，使高维语言特征查询达到交互速度。论文报告在高分辨率设置下相对 LERF 语言场有约 199 倍速度提升。官方代码为 \url{https://github.com/minghanqin/LangSplat}，GitHub 关注度很高，约 1066 stars；OpenAlex 约 172 次引用。

\textbf{实现细节：} 方法先训练普通 3DGS 获得几何和颜色场，再从多视图图像中提取 SAM 层级掩码和 CLIP 语义特征。直接把 512 维 CLIP 特征存到每个高斯会导致显存和渲染代价过高，因此 LangSplat 为每个场景训练语言 autoencoder，将高维语言特征压缩到低维 latent；训练时把高斯语言 latent 像颜色一样 splat 到视图中，与 2D 语言特征监督对齐。查询时，文本嵌入与渲染/解码后的语言特征做相似度匹配，得到开放词汇目标区域。SAM 的层级语义用于减少边界模糊和尺度查询开销，使目标边界比直接蒸馏 CLIP 更清楚。 核心贡献是把 2D 视觉语言特征变成可实时渲染的 3D language field：用 SAM 生成层级语义区域，用 CLIP/OpenCLIP 提供语言监督，再用 scene-wise autoencoder 将高维语言特征压缩后存入高斯。 scene-wise autoencoder 带来额外训练和解码步骤，后续 LangSplatV2 进一步用稀疏字典替代 MLP 解码；点级实例选择和跨场景泛化也由 OpenGaussian、Gen-LangSplat 等后续方法继续改进。
\textbf{方法信息：} 代码已公开，仓库为 minghanqin/LangSplat (1066 stars)；GitHub 关注度很高，minghanqin/LangSplat 约 1066 stars；引用很多，约 172 次引用，5.4 次/月。入选依据：企业:否；代码:是；星标:是；引用:是。


% METHOD: 2311.18482
\paragraph{Language Embedded 3D Gaussians（2023.11）\cite{2311_18482}}
是早期把语言嵌入接入3DGS的代表。它从多视图图像提取CLIP和DINO混合语言特征，再用聚类码本量化局部冗余特征，降低每个高斯的存储开销。为处理多视角语义不一致，它为高斯学习不确定性值，并根据不确定性降低语言特征的空间频率，让不可靠区域更平滑。训练时渲染语义特征图，与2D特征监督对齐。相对Feature 3DGS，它更专门处理语言特征的内存和跨视角噪声问题。

流程上，层级随机裁剪先让CLIP看到不同尺度的局部区域，DINO特征补充边界和纹理细节，二者拼接成每个视图的密集语言特征图。随后，k-means码本把高维语义压成低维索引，高斯只需学习紧凑语义表示。多视角不一致通过不确定性来吸收：不可靠区域的语义特征被平滑，可靠区域保留更高频细节。最终，3DGS在渲染RGB的同时渲染语义特征图，和2D语言特征监督对齐，从而支持文本查询。
\textbf{实现细节：} Language Embedded 3D Gaussians：密集像素级特征。早期方法 Language Embedded 3D Gaussians（2023.11） 与 LangSplat 同期提出，走了一条不同的路线。动机。LangSplat 采用 SAM 掩码区域裁剪后提取 CLIP 特征，掩码粒度影响特征质量；该方法直接从 CLIP 和 LSeg 的密集像素级特征出发，通过量化压缩和不确定性学习的平滑嵌入来处理多视角不一致性。LangSplat 采用 SAM 掩码区域裁剪后提取 CLIP 特征，而该方法直接从 CLIP 和 LSeg 的密集像素级特征出发，通过专用量化方案将高维特征压缩到低维码本空间以缓解内存需求，并引入基于不确定性学习的平滑嵌入过程来对抗多视角特征不一致性和点状表示的高频偏置。
\textbf{方法信息：} 引用很多，Semantic Scholar 约 229 次。
\textbf{设计动机与边界：} 相对原始 3DGS 以独立高斯和图像重建误差为主的设定，Language Embedded 3D Gaussians for Open-Vocabulary Scene Understanding 把问题落在以下可实现的改动上：方法包含四个步骤。第一步，从多视图图像中提取密集语言特征：使用层级随机裁剪技术提取CLIP特征，结合DINO特征增强细节，拼接形成混合语言特征图。第二步，量化方案：利用局部语义特征的冗余性，通过k-means聚类构建紧凑的码本，将高维原始语义特征量化为低维紧凑表示，显著降低计算和存储开销。第三步，不确定性引导嵌入：针对多视图图像间的语义不一致性和点云表示的高频归纳偏置问题，学习每个高斯的不确定性值，根据不确定性降低语义特征的空间频率，使语言特征更平滑且保持高精度。第四步，训练时通过渲染语义特征图并与2D特征监督对比，联合优化3D-GS属性和语言嵌入。其设计目的不是单独增加一个模块，而是用上述表示、监督或渲染约束消除原流程在目标任务中的具体失配。提出Language Embedded 3D Gaussians表示方法，通过专用量化方案将高维语义特征压缩存储在3D高斯上，大幅降低内存需求（从数GB到可管理级别）。设计不确定性引导的低频嵌入机制，解决多视图特征不一致和点云表示的高频归纳偏置问题，实现高精度开放词汇查询。适用时还应保留论文给出的边界：语义特征精度受CLIP和DINO模型限制；量化可能损失部分语义信息；对非常小物体的查询精度有限。
% METHOD: 2404.14249
\paragraph{CLIP-GS（2024.04）\cite{2404_14249}}
面向室内语义理解，核心改动是用低维语义表示替代直接存储CLIP高维特征。它利用物体内部语义一致性学习紧凑特征，并用码本压缩提高语义渲染速度；同时在2D和3D层面做一致性约束：2D层面通过多视角渲染语义结果聚合更稳定监督，3D层面鼓励同物体高斯特征接近。相对早期语言嵌入方法，CLIP-GS的重点是效率和室内跨视角一致，而不是任意场景的开放词汇检索。

语义属性压缩不把512维CLIP特征直接附着到每个高斯，而是学习十几维量级的紧凑语义向量，再通过码本量化减少渲染负担。视角不一致则通过多视图监督和3D一致正则处理，使同一物体的高斯在语义空间聚集，减少不同视角CLIP特征各说各话。这样，CLIP-GS把语言高斯从“能查询”推进到“能快速查询”。
\textbf{方法信息：} Baidu 参与了这篇论文；仓库：\url{https://github.com/gbliao/clip-gs}，当前约 56 stars，核验时主要是README，完整代码暂未确认。
\textbf{设计动机与边界：} 相对原始 3DGS 以独立高斯和图像重建误差为主的设定，CLIP-GS 把问题落在以下可实现的改动上：CLIP-GS解决现有3DGS语义理解方法的两个核心问题。(1) 效率瓶颈：现有方法将高维CLIP特征（如512维）直接附加到每个高斯上，严重影响光栅化效率。SAC方法利用物体内语义的自然统一性，学习低维紧凑的语义表示（如16维），通过码本量化进一步压缩，使语义渲染速度达到100+FPS。(2) 语义不一致性：2D CLIP特征的视角不一致性导致3D语义分割出现不一致。3DCR在2D和3D两个层面引入一致性约束：在2D层面，利用3DGS渲染的语义结果建立跨视角一致性约束（通过特征聚合得到更稳定的2D语义标签作为监督）；在3D层面，鼓励属于同一物体的3D高斯原语具有相似的特征。两者结合显著提升了分割的准确性和一致性。其设计目的不是单独增加一个模块，而是用上述表示、监督或渲染约束消除原流程在目标任务中的具体失配。提出CLIP-GS，通过语义属性紧凑性（SAC）和3D一致性正则化（3DCR）实现高效一致的3D室内语义理解。SAC利用物体内统一语义学习紧凑的语义高斯表示，实现100+FPS语义渲染。3DCR在2D和3D域同时强制语义一致性。在ScanNet和Replica上mIoU分别提升21.20%和13.05%。适用时还应保留论文给出的边界：依赖SAM和CLIP等预训练模型的质量；对开放词汇中罕见概念的理解能力有限；在室外复杂场景上的泛化性未验证。
% METHOD: 2405.17596
\paragraph{GOI（2024.05）\cite{2405_17596}}
GOI针对的是开放词汇查询里的“阈值不可靠”问题。LERF、LangSplat一类方法通常把文本嵌入和渲染出的语义特征做相似度比较，再用经验阈值截出目标区域；当查询带有空间关系或长尾属性时，固定阈值容易把相邻物体一起选中，或者只选中目标的一部分。GOI的改变有两处：先用APE提取像素对齐的视觉语言特征，再用Trainable Feature Clustering Codebook（TFCC）把高维语义压成低维码本特征并嵌入高斯；查询时不直接按相似度截断，而是在语义空间中学习一个可优化超平面，把“目标高斯”和“非目标高斯”分开。

具体流程是：多视角图像先经过APE得到边界更清楚的像素级语义图，3DGS训练时同步学习每个高斯的低维语义向量，TFCC用码本条目的组合来表示原始高维特征，降低渲染和存储成本。用户给出文本查询后，GOI用2D referring expression segmentation模型在图像上产生目标区域提示，并用这些2D结果微调语义超平面；最终被超平面判为正侧的高斯就是Gaussians of Interest。相对LangSplat的多尺度特征匹配，GOI把“找阈值”变成了“学决策边界”，因此更适合带关系的短语查询和边界敏感的局部编辑。
\textbf{实现细节：} GOI：语义空间超平面。GOI（2024.05，Gaussian Open-vocabulary Instance） 引入可优化的语义空间超平面，通过二维参照表达式分割模型微调特征空间中的分割边界，并用可训练特征聚类码本（TFCC，Trainable Feature Clustering Codebook）压缩高维语言嵌入，实现更精确的开放词汇查询定位。动机。LangSplat 的语言特征未在特征空间中建立明确的分割边界，查询时依赖余弦相似度的经验阈值；GOI 用可优化的超平面显式定义分割边界，使开放词汇查询的定位更加精确。
\textbf{实现补充：} 实现上还要核对输入、输出、参数、特征、初始化、训练损失、正则、渲染或推理接口，以及优化中的剪枝、分裂、采样、编码或解码步骤；这些抓手决定该方法到底改了 3DGS 管线的哪一环。
\textbf{设计动机与边界：} 相对原始 3DGS 以独立高斯和图像重建误差为主的设定，GOI 把问题落在以下可实现的改动上：GOI框架包含三个核心模块。（1）语义特征提取：使用视觉-语言基础模型APE提取像素级对齐的语义特征图，而非CLIP的全局特征或多尺度裁剪方式，避免了语义边界模糊问题。（2）可训练特征聚类码本（TFCC）：基于场景先验将高维语义特征压缩为紧凑低维向量并嵌入每个3D高斯。TFCC通过聚类学习一个码本，每个高斯的语义特征用码本条目的加权组合表示，显著减少存储和渲染成本。（3）可优化语义空间超平面：不同于现有方法使用固定经验阈值选择语义特征距离近的区域，GOI将特征选择过程建模为特征空间中的超平面划分，仅保留与查询高度相关的特征。利用现成的2D指代表达分割（RES）模型来微调该超平面参数，实现目标区域与其他区域的更精确区分。整个框架端到端可优化，支持任意自然语言查询3D场景中的感兴趣区域。其设计目的不是单独增加一个模块，而是用上述表示、监督或渲染约束消除原流程在目标任务中的具体失配。提出GOI框架，首次将可优化语义空间超平面引入3DGS的开放词汇查询任务。TFCC模块利用场景先验将噪声高维语义特征压缩为紧凑低维向量嵌入3DGS，大幅降低存储和渲染开销。用RES模型微调语义空间超平面替代固定经验阈值，显著提升查询精度。在多个基准上超越LERF、LangSplat等先前方法。适用时还应保留论文给出的边界：依赖APE模型的像素级语义特征提取质量；TFCC压缩可能损失部分细粒度语义信息；RES模型微调需要额外计算开销；仅在室内外静态场景验证，未涉及动态场景。
\textbf{方法信息：} Tencent 参与了这篇论文；引用很多，Semantic Scholar 约 60 次。



% METHOD: 2406.02058
\paragraph{OpenGaussian（2024.06）\cite{2406_02058}}
OpenGaussian 是点级开放词汇 3DGS 的重要锚点。Feature 3DGS、LangSplat 等方法已经能把 2D 语言特征提升到三维并渲染成像素级语义图，但它们的查询通常仍落在“从某个视角看哪些像素相关”上；机器人抓取、点云理解或遮挡部件识别更需要直接判断“哪些三维高斯点属于这个语义目标”。OpenGaussian 的改变是先学习具有三维一致性的点级实例特征，再把这些实例特征和 2D CLIP 语义关联起来，而不是直接把高维 CLIP 向量塞进每个高斯。

第一步是实例特征学习。方法为每个高斯附加一个低维实例特征，例如 6 维向量，并把该特征像颜色一样通过 $\alpha$-blending 渲染到训练视图。监督来自每个视图的 SAM 二值掩码，但不要求先把所有视图的掩码跨帧关联起来。对一个 SAM 掩码，渲染出的特征在掩码内部应接近同一个均值，这是 intra-mask smoothing loss；不同掩码的均值特征应彼此远离，这是 inter-mask contrastive loss。由于同一个三维高斯会在不同视角被反复渲染，这两个损失把“同物体内一致、不同物体间可分”的约束落回三维点级特征。

第二步是两级码本离散化。仅靠连续特征做点击选择时，很难设一个对所有场景都稳定的相似度阈值；而 $\alpha$ 混合也会让同一物体内部出现特征噪声。OpenGaussian 因此先用粗码本结合三维位置做位置相关聚类，再用细码本细化实例特征，使同一实例的高斯更接近离散一致的编码。第三步再做实例级 3D-2D 特征关联：把三维高斯投影到 2D 掩码，利用对应掩码区域的 CLIP 特征给三维实例挂上开放词汇语义。最终查询时，文本和三维高斯的语言特征比较，可以直接选出相关高斯点并渲染验证。相对 LangSplat，OpenGaussian 的核心增量是把开放词汇理解从视图像素层推进到三维点层。
\textbf{锚点细化：} OpenGaussian：针对 LangSplat 将 CLIP 特征蒸馏到每个高斯低维向量、查询精度受限于区域级粗糙匹配（一个高斯可能横跨多个物体）的局限，用\textbf{6 维实例特征 + 由粗到精码本离散化}替代区域级 CLIP 特征压缩，核心思想是\textbf{学习三维点级实例特征并通过无跨帧关联的 SAM 掩码训练三维一致性}——通过实例级三维-二维特征关联将三维点链接到 CLIP 特征，首次实现基于 3DGS 的三维点级开放词汇理解，超越 LangSplat 的区域级查询能力。
\textbf{实现细节：} OpenGaussian：三维点级开放词汇理解。动机。LangSplat 将 CLIP 特征蒸馏到每个高斯的低维向量中，查询精度受限于区域级的粗糙匹配；OpenGaussian 学习 6 维实例特征并通过由粗到精的码本离散化，实现精确的三维点级语义查询。其核心差异在于：LangSplat 将 CLIP 特征蒸馏到每个高斯的低维向量中，而 OpenGaussian 学习 6 维实例特征 $f \in \mathbb{R}^{}6$（无需额外的维度压缩或量化步骤），通过无跨帧关联的 SAM 掩码训练具有三维一致性的实例特征，并提出由粗到精的两阶段码本进行特征离散化，最后通过实例级三维-二维特征关联将三维点链接到 CLIP 特征。
\textbf{方法信息：} Baidu VIS 参与了这篇论文；代码：\url{https://github.com/yanmin-wu/OpenGaussian}；GitHub 关注度较高，约 212 stars；引用很多，Semantic Scholar 约 149 次。



% METHOD: 2412.13654
\paragraph{GAGS（2024.12）\cite{2412_13654}}
GAGS把LangSplat式语言特征蒸馏中的核心噪声定位到“粒度不一致”。同一只碗在近景里可能被SAM切成碗、面、筷子，在远景里又被切成一个整体；如果把这些区域的CLIP特征都蒸馏到同一个3D位置，高斯语言场会学到平均化、退化的特征。GAGS没有简单增加多个语言场，而是先让2D分割本身更稳定：从已训练3DGS渲染深度图，把图像划成patch，并根据深度决定SAM提示点密度，远处区域采样更密以获得足够细的掩码，近处区域采样更疏以避免把完整物体切碎。

第二个改变发生在蒸馏损失上。每个高斯仍学习一个可渲染语义特征，训练时把特征图splat到视图中，与由SAM区域和CLIP编码得到的2D特征对齐；但GAGS额外解码一个可学习粒度因子，用它在sub-part、part、object等候选粒度中选择多视角更一致的监督。这样，它把“该用哪个粒度的CLIP特征”交给场景自身学习，而不是像LangSplat那样为多个粒度各训练一套特征场。相对GOI的查询后超平面微调，GAGS主要修正训练阶段的跨视角监督，因此查询时只需要渲染单个粒度感知语言场。
\textbf{实现细节：} GAGS：分割粒度一致性。GAGS（2024.12，Granularity-Aware Gaussian Splatting） 从分割粒度角度解决多视角不一致性，指出 LangSplat 所依赖的 SAM 掩码粒度在不同视角下的不稳定性是特征不一致的重要来源。动机。SAM 掩码的粒度因视角距离而异——近处物体被过度细分为小掩码，远处物体被合并为大掩码，导致同一物体在不同视角的特征不一致；GAGS 将提示点密度与相机距离关联以稳定分割粒度。GAGS 提出将 SAM 的提示点密度与相机距离关联以改善多视角分割一致性，并学习粒度因子以无监督方式选择多视角一致的特征用于蒸馏，从而在特征提取源头解决一致性问题。
\textbf{实现补充：} 实现上还要核对输入、输出、参数、特征、初始化、训练损失、正则、渲染或推理接口，以及优化中的剪枝、分裂、采样、编码或解码步骤；这些抓手决定该方法到底改了 3DGS 管线的哪一环。
\textbf{设计动机与边界：} 相对原始 3DGS 以独立高斯和图像重建误差为主的设定，GAGS 把问题落在以下可实现的改动上：方法分两步：(1) 粒度感知分割：观察到SAM分割粒度由提示点密度控制，且远距离物体需要更小分割掩码。因此将提示点密度与渲染深度图关联：将图像分patch，根据深度计算每个patch的采样点数量。远处视角用更高密度分割更小物体，近处用更低密度。这显著提升了跨视角分割一致性。(2) 粒度感知蒸馏：解码一个可学习粒度因子，在无监督方式下学习，仅选择多视角一致的2D CLIP特征进行蒸馏。每个Gaussian关联可训练特征向量，通过splatting渲染特征图与CLIP特征图对齐训练。无需学习三个独立特征场（如LangSplat），降低内存和时间成本。其设计目的不是单独增加一个模块，而是用上述表示、监督或渲染约束消除原流程在目标任务中的具体失配。提出两项关键策略：(1) SAM提示点密度与相机距离关联，显著提升分割多视角一致性；(2) 可学习粒度因子无监督选择多视角一致的2D特征进行蒸馏。推理速度比基线快2倍。适用时还应保留论文给出的边界：依赖SAM分割质量和CLIP特征提取器的性能。
\textbf{方法信息：} Huawei 参与了这篇论文；代码：\url{https://github.com/whu-usi3dv/gags}。



% METHOD: 2503.06271
\paragraph{SplatTalk（2025.03）\cite{2503_06271}}
把3DGS变成3D视觉问答的语言token来源。它先从LLaVA-OV提取每帧视觉语言token，并训练自编码器压缩到低维空间；再基于FreeSplat式泛化3DGS框架，同时预测RGB高斯和语言特征高斯。训练时渲染RGB和语言特征图，分别对齐图像和压缩后的LLaVA特征。推理时，3DGS中的语言特征直接作为token输入LLM回答问题。相对开放词汇分割，SplatTalk不只输出掩码，而是把可渲染3D语言场作为问答模型的上下文。
\textbf{实现细节：} SplatTalk：零样本三维视觉问答。SplatTalk（2025.03） 利用泛化三维高斯框架（基于 FreeSplat 扩展）生成可直接输入预训练 LLM 的三维标记，实现仅依赖带位姿图像的有效零样本三维视觉问答。其关键创新在于特征压缩流水线：先用 LLaVA-OV（SigLIP 视觉编码器 + Qwen2 语言模型）提取 $27 \times 27 \times 3584$ 维高维特征，再通过自编码器将每个 3584 维特征向量压缩至 256 维（压缩比约 14 倍），以 256 维低维特征作为伪真值训练 3DGS 语言表示。
\textbf{方法信息：} Google / DeepMind 参与了这篇论文。
\textbf{设计动机与边界：} 相对原始 3DGS 以独立高斯和图像重建误差为主的设定，SplatTalk 把问题落在以下可实现的改动上：SplatTalk的流程包含三个阶段：(1)特征自编码器训练——从LLaVA-OV提取每帧RGB图像的高维语言特征（视觉token），训练自编码器将其压缩到低维空间，作为伪真值；(2)自监督3D-语言GS训练——基于FreeSplat架构构建泛化3DGS框架，CNN骨干提取多尺度图像特征，通过自适应代价体积预测深度，反投影为3D高斯三元组。在渲染RGB图像的同时渲染语言特征图，与压缩后的LLaVA特征对齐训练。RGB和语言特征联合训练确保多模态一致性；(3)3D VQA推理——从训练好的3DGS中提取语言特征，直接作为token输入LLM，结合语言提示生成自由形式的回答。无需3D输入（点云/网格）。其设计目的不是单独增加一个模块，而是用上述表示、监督或渲染约束消除原流程在目标任务中的具体失配。提出SplatTalk，首个利用泛化3DGS框架生成可直接输入预训练LLM的3D token的方法，实现零样本3D视觉问答。通过自监督方式将2D语言特征蒸馏到3D高斯中，仅用带位姿的RGB图像即可实现与SOTA 3D LMM竞争的性能。适用时还应保留论文给出的边界：依赖LLaVA特征质量；对复杂空间推理问题的能力有限；特征维度压缩可能丢失部分信息。
% METHOD: 2504.11893
\paragraph{CAGS（2025.04）\cite{2504_11893}}
GAGS从2D掩码粒度入手，CAGS则把问题进一步推到3D高斯之间的上下文。它的判断是：颜色可以在同一物体内部剧烈变化，但语义通常更稀疏、更连续；如果每个高斯只独立学习语言特征，SAM跨视角粒度不一致会把一个物体切成碎片。CAGS在高斯集合上构建局部图，用最远点采样选出约10\%高斯作为锚点，每个锚点连接16个近邻，再通过两层图传播聚合邻域的语义和颜色上下文，使同一局部物体区域的语言特征更一致。

监督也从逐像素对比改成掩码级对比。CAGS在SAM掩码内部计算渲染特征的质心，再让跨视角、同一对象区域的质心靠近，而不是把每个像素都当成独立正负样本。训练时先冻结或稳定高斯位置，预计算邻域关系，避免每轮都重建大规模图；推理时仍按文本特征和高斯语言特征相似度选择对象。相对OpenGaussian的点级语义和GAGS的粒度选择，CAGS强调“邻域传播 + 掩码质心监督”，适合减少开放词汇3D实例分割中的碎片化。
\textbf{实现细节：} CAGS：局部图上下文传播。CAGS（2025.04） 通过构建局部图在高斯间传播上下文特征，减少不一致粒度导致的噪声，并采用掩码中心对比学习平滑 SAM 衍生的跨视角特征。与 Feature 3DGS 的逐像素特征蒸馏不同，CAGS 的图传播机制利用高斯的空间邻域关系实现特征平滑，从根本上解决了跨视角 SAM 掩码不一致的问题。
\textbf{实现补充：} 实现上还要核对输入、输出、参数、特征、初始化、训练损失、正则、渲染或推理接口，以及优化中的剪枝、分裂、采样、编码或解码步骤；这些抓手决定该方法到底改了 3DGS 管线的哪一环。
\textbf{方法信息：} 代码：\url{https://github.com/wistzz/cags}。
\textbf{设计动机与边界：} 相对原始 3DGS 以独立高斯和图像重建误差为主的设定，CAGS 把问题落在以下可实现的改动上：CAGS框架包含三个核心组件：(1) 上下文特征传播：在3DGS上构建局部图，使用最远点采样（FPS）选择10%高斯基元作为锚点，每个锚点连接k=16个最近邻。通过图卷积在相邻高斯基元间传播语义特征，平滑同一物体内的特征表示，减少粒度不一致的噪声。(2) 掩码感知对比学习：不同于逐像素InfoNCE损失，在SAM掩码内计算特征质心，跨视角对质心施加对比监督，减少不一致分割的影响。(3) 预计算策略：预先计算邻域关系，降低大规模场景的训练计算成本。关键洞察：语义空间比颜色空间更稀疏，同一物体内部语义一致，需要更强的空间上下文利用。其设计目的不是单独增加一个模块，而是用上述表示、监督或渲染约束消除原流程在目标任务中的具体失配。提出上下文感知高斯泼溅框架，通过局部图构建上下文特征传播和掩码感知对比学习，解决SAM跨视角分割粒度不一致问题。适用时还应保留论文给出的边界：图构建和FPS采样的超参数需要调优；对极度细碎物体的处理仍有挑战；预计算增加存储需求。
% METHOD: 2505.24746
\paragraph{LaGa（2025.05）\cite{2505_24746}}
处理语言高斯中的视角依赖语义。它先训练高斯亲和特征，把同一SAM掩码内的高斯聚到一起，不同掩码分开，再根据阈值把高斯分成3D物体实例。之后，每个3D物体从多视角2D语义特征中自适应聚类出多个代表性描述符，权重由全局语义对齐和该描述符出现频率共同决定。相对直接对每个高斯存单一语言向量，LaGa把语义提升到对象级，并保留同一物体在不同视角下可能呈现的多个描述。

场景分解阶段对每个视角渲染高斯亲和特征图，并把SAM掩码内部的平均特征作为掩码原型；对比损失让同一掩码内的高斯靠近，让不同掩码分离。训练后，阈值化亲和特征即可把高斯分配给3D物体实例，建立跨视角的对象连接。语义聚合阶段不强迫每个对象只有一个语言向量，而是从多视角2D语义特征中聚类出一组代表性描述符，并根据全局语义对齐程度和出现频率赋权。
\textbf{方法信息：} Huawei 参与了这篇论文；代码：\url{https://github.com/SJTU-DeepVisionLab/LaGa}。
\textbf{方法信息：} Huawei 参与了这篇论文；代码：\url{https://github.com/sjtu-deepvisionlab/laga}。
\textbf{设计动机与边界：} 相对原始 3DGS 以独立高斯和图像重建误差为主的设定，Tackling View-Dependent Semantics in 3D Language Gaussian Splatting 把问题落在以下可实现的改动上：LaGa的流程包含两个阶段：1)3D场景分解：采用对比学习方法训练高斯亲和特征，指示两个3D高斯是否属于同一物体。对每个视角渲染亲和特征图，为每个2D掩码分配掩码亲和原型（通过掩码平均池化）。训练目标鼓励同一掩码内的特征聚集、不同掩码的特征分离。训练后，同一3D物体的亲和特征和掩码原型收敛到紧凑聚类。通过阈值化将高斯分配到3D物体实例，建立跨视角语义连接；2)视角聚合语义表示：对每个3D物体，通过自适应聚类从多视角2D语义特征中提取代表性语义描述符集合。分配权重基于两个因素：全局对齐（描述符与全局特征的方向相似性）和视角频率（描述符在多视角中出现的频率）。加权描述符实现直接3D场景理解，无需依赖2D渲染特征图，更全面地捕获视角依赖的3D语义。其设计目的不是单独增加一个模块，而是用上述表示、监督或渲染约束消除原流程在目标任务中的具体失配。提出LaGa(Language Gaussians)，解决3D语言高斯溅射中的视角依赖语义问题。核心发现：简单将2D语义投影到3D高斯会导致不完整的语义分配（假阳性和假阴性）。LaGa通过两步解决：1)3D场景分解——将多视角2D掩码分组为一致的3D物体，建立跨视角语义连接；2)视角聚合语义表示——通过自适应聚类提取代表性语义描述符集合，基于多视角语义重加权。在LERF-OVS数据集上，相同设置下mIoU比前SOTA提升+18.7%。适用时还应保留论文给出的边界：依赖2D分割和语义先验的质量；场景分解对复杂遮挡场景可能不够精确；自适应聚类的参数选择影响性能。
% METHOD: 2506.23309
\paragraph{SurgTPGS（2025.06）\cite{2506_23309}}
SurgTPGS把文本可提示语义高斯带到手术场景。手术视频里的组织会形变，器械会遮挡，光照和反光也比普通室内场景更不稳定；如果只把CAT-Seg或SAM得到的2D语义逐帧贴到高斯上，语义区域会随形变漂移。SurgTPGS以Deform3DGS式动态表示为底座，先从训练视图提取SAM分割和CAT-Seg像素级语义特征，再用自编码器压缩语义维度，避免把高维VLM特征直接挂到每个动态高斯上。

它的增量在动态语义跟踪。语义感知变形跟踪模块用卷积网络增强局部语义一致性，再通过斜率sigmoid概率图和门控机制判断哪些语义区域应随组织变形一起移动；语义区域感知优化则把区域级语义监督接到动态高斯训练里，让几何重建、变形场和语义平滑共同优化。推理时，渲染出的语义特征与文本嵌入匹配，得到可提示的3D手术区域查询。相对静态LangSplat类方法，SurgTPGS的关键不是换一个语言特征，而是把“语言特征如何随可变形组织运动”纳入训练。
\textbf{实现细节：} SurgTPGS：手术场景语义。SurgTPGS（2025.06） 将文本提示式高斯引入手术场景，结合 SAM 和视觉-语言模型的三维语义特征学习策略以及语义感知变形跟踪，实现手术工具和解剖结构的实时文本提示三维查询。与通用 LangSplat 使用静态高斯不同，SurgTPGS 采用灵活变形建模方案（FDM），通过傅里叶和多项式基函数 $\Phi(\tau; \omega, \mu, \sigma)$ 参数化形变，将语义特征与变形高斯绑定，使文本提示能在动态手术过程中持续追踪解剖结构。
\textbf{实现补充：} 实现上还要核对输入、输出、参数、特征、初始化、训练损失、正则、渲染或推理接口，以及优化中的剪枝、分裂、采样、编码或解码步骤；这些抓手决定该方法到底改了 3DGS 管线的哪一环。
\textbf{方法信息：} 代码：\url{https://github.com/lastbasket/surgtpgs}。
\textbf{设计动机与边界：} 相对原始 3DGS 以独立高斯和图像重建误差为主的设定，SurgTPGS 把问题落在以下可实现的改动上：SurgTPGS首先结合SAM（分割）和CAT-Seg（视觉语言模型）提取像素级语义特征，通过自编码器压缩特征维度。语义感知变形跟踪模块使用基于卷积的网络增强语义特征的空间一致性，通过斜率sigmoid概率图和门控机制跟踪语义区域的变形。语义区域感知优化利用区域级语义信息监督训练，特别促进重建质量和语义平滑性。渲染的语义特征通过映射到文本嵌入实现文本可提示的3D查询。整体框架基于Deform3DGS，支持可变形组织的动态重建。其设计目的不是单独增加一个模块，而是用上述表示、监督或渲染约束消除原流程在目标任务中的具体失配。提出SurgTPGS，首个文本可提示的手术场景Gaussian Splatting方法，结合SAM和CAT-Seg实现3D语义特征学习，提出语义感知变形跟踪和语义区域感知优化，实现手术工具和解剖结构的实时文本查询3D理解。适用时还应保留论文给出的边界：依赖SAM和CAT-Seg的分割质量；语义特征训练增加计算开销；对出血和严重遮挡的处理仍有局限。
% METHOD: 2507.07136
\paragraph{LangSplatV2（2025.07）\cite{2507_07136}}
LangSplatV2 解决的是语言高斯的实时性问题。LangSplat 的关键设计是用 SAM 的多尺度掩码提取较干净的 CLIP 语义特征，再把 512 维 CLIP 特征压缩成低维潜变量存到高斯上；查询时先渲染低维特征图，再用 MLP 解码回 CLIP 空间。这个设计节省训练内存，但推理时的 MLP 解码成为瓶颈，高分辨率开放词汇查询仍难实时。LangSplatV2 的改变是完全取消这个重解码器。

它把每个高斯的高维语言特征表示为全局字典上的稀疏系数：场景共享一个语义字典，每个高斯只存少量非零系数，高维 CLIP 特征由这些系数和字典项线性组合得到。训练时仍继承 LangSplat 的 SAM 三层语义尺度和 OpenCLIP 特征监督，但优化目标从“学习低维潜变量并让 MLP 解码”变成“学习稀疏系数场和全局字典”。渲染时也不再 splat 完整 512 维特征，而是只 splat 少量非零系数；之后用一次轻量矩阵乘法恢复高维特征，CUDA 实现进一步减少系数混合开销。这样，渲染成本接近低维特征泼溅，语义表达又保留高维 CLIP 的辨别能力。

相对 OpenGaussian，LangSplatV2 不是为了把查询推进到点云任务，而是为了让高维开放词汇语言场达到交互速度；相对 LangSplat，它用稀疏字典替代 autoencoder/MLP。论文报告高维特征 splatting 达到 476.2 FPS，三维开放词汇文本查询达到 384.6 FPS，分别比 LangSplat 快 42 倍和 47 倍，同时保持或提升查询精度。
\textbf{实现细节：} LangSplatV2（2025.07） 针对 LangSplat 的速度瓶颈进行优化。动机。LangSplat 的场景专属自编码器是推理速度的主要障碍，每次查询都需解码高维特征；LangSplatV2 用稀疏编码替代自编码器，假设每个高斯是全局词典中的稀疏编码，完全消除解码器需求，实现 476.2 FPS 的高维特征渲染。该工作发现重量级解码器是推理速度的主要障碍，转而假设每个高斯是全局词典中的稀疏编码，学习三维稀疏系数场从而完全消除解码器需求。通过 CUDA 优化的稀疏系数 splatting，实现 476.2 FPS 的高维特征渲染和 384.6 FPS 的三维开放词汇查询，分别比 LangSplat 快 42 倍和 47 倍。
\textbf{设计动机与边界：} 相对原始 3DGS 以独立高斯和图像重建误差为主的设定，LangSplatV2 把问题落在以下可实现的改动上：LangSplatV2首先对LangSplat进行详细的性能分析，发现重量级MLP解码器是主要速度瓶颈。为解决此问题，提出将每个高斯视为全局字典中的稀疏码，学习一个3D稀疏系数场，使得高维CLIP特征可以通过稀疏系数和全局字典的乘积来近似表示，从而完全消除解码器。在渲染阶段，提出高效的稀疏系数泼溅方法：由于稀疏系数中大部分值为零，只需泼溅非零系数对应的字典项，将高维特征渲染的时间复杂度降低到接近超低维特征泼溅的水平。同时利用CUDA进行深度优化。该方法保持与LangSplat相同的SAM语义层级结构，使用三层语义尺度获取精确的3D语言场。为了消除原始 3DGS 在该目标任务中的具体失配，方法把上述表示、监督或渲染约束作为可学习和可执行的改动，而不是单独堆叠模块。提出LangSplatV2，通过假设每个高斯作为全局字典中的稀疏码，学习3D稀疏系数场完全消除重量级MLP解码器。结合CUDA优化的稀疏系数泼溅方法，实现高维特征泼溅476.2 FPS和三维开放词汇文本查询384.6 FPS，分别比LangSplat加速42倍和47倍，同时保持或提升查询精度。适用时还应保留论文给出的边界：依赖预训练CLIP特征的质量和语义表达能力；稀疏编码假设可能在某些极端场景下引入近似误差；全局字典的大小需要在速度和精度之间权衡。
\textbf{方法信息：} 代码：\url{https://github.com/zhaoyujie2002/langsplatv2}；GitHub 关注度较高，约 237 stars；引用很多，Semantic Scholar 约 33 次。



% METHOD: 2507.19830
\paragraph{Language Embedded 3DGS in the Wild（2025.07）\cite{2507_19830}}
Language Embedded 3DGS in the Wild处理的是无约束照片集合里的语言高斯。WE-GS、WildGaussians一类方法已经能从互联网照片中重建同一地标或建筑，但这些图像有不同光照、后期风格、游客遮挡和季节变化；如果直接从原图提取CLIP特征并蒸馏，语言场会把“瞬态外观”和“稳定结构”混在一起。该方法先利用in-the-wild 3DGS在同一视角渲染多种外观版本，再从这些版本提取多外观CLIP特征，同时估计瞬态和外观不确定性。

语义训练时，不确定性图决定哪些视图和外观更可信：高不确定区域对语言特征蒸馏贡献降低，稳定结构区域则被反复聚合。语言特征进入瞬态不确定性感知自编码器压缩，再形成多外观语言场；最后用后集成策略融合不同外观下的语义响应。相对LangSplat默认输入是较干净、配准良好的训练图像，这个方法把WildGaussians处理外观扰动的思想扩展到开放词汇语义，使建筑构件和风格查询不被游客、阴影或曝光变化主导。
\textbf{实现细节：} Taking Language Embedded 3DGS into the Wild：互联网照片集。Taking Language Embedded 3DGS into the Wild（2025.07） 将语言嵌入三维高斯扩展到无约束互联网照片集。动机。LangSplat 假设受控拍摄条件（固定光照、无瞬态物体），但互联网照片集存在光照剧变、瞬态物体（行人、车辆）和不同外观模式；该方法通过多外观 CLIP 特征和瞬态/外观不确定性图，用不确定性感知的自编码器将这些挑战统一处理。提出多外观 CLIP 特征和瞬态/外观不确定性图引导优化，引入瞬态不确定性感知自编码器和后集成策略。与 LangSplat 假设受控拍摄条件不同，互联网照片集存在光照剧变、瞬态物体（行人、车辆）和不同外观模式，该方法通过不确定性感知的自编码器将这些挑战统一处理。
\textbf{实现补充：} 实现上还要核对输入、输出、参数、特征、初始化、训练损失、正则、渲染或推理接口，以及优化中的剪枝、分裂、采样、编码或解码步骤；这些抓手决定该方法到底改了 3DGS 管线的哪一环。
\textbf{方法信息：} 代码：\url{https://github.com/yuzewang1998/takinglangsplatw}。
\textbf{设计动机与边界：} 相对原始 3DGS 以独立高斯和图像重建误差为主的设定，Taking Language Embedded 3D Gaussian Splatting into the Wild 把问题落在以下可实现的改动上：基于in-the-wild 3DGS方法(如WE-GS)构建的辐射场，扩展开放词汇语义。核心创新：(1)从同一视角渲染多外观图像，提取多外观CLIP特征和瞬态/外观不确定性图引导后续优化；(2)瞬态不确定性感知自编码器压缩语言特征；(3)多外观语言场3DGS表示学习；(4)后集成策略融合多外观语言特征。提出PT-OVS新基准数据集用于评估无约束照片集的开放词汇分割性能。其设计目的不是单独增加一个模块，而是用上述表示、监督或渲染约束消除原流程在目标任务中的具体失配。扩展语言嵌入3DGS到无约束照片集，提出多外观CLIP特征和瞬态/外观不确定性图引导优化，设计瞬态不确定性感知自编码器和后集成策略，并构建PT-OVS基准数据集。适用时还应保留论文给出的边界：依赖in-the-wild 3DGS重建质量；网络照片的极端光照变化影响语言特征提取；PT-OVS数据集规模有限。
% METHOD: 2508.14278
\paragraph{GALA（2025.08）\cite{2508_14278}}
基于Scaffold-GS训练场景特定实例特征，再用双可学习码本和交叉注意力映射到通用语言空间。第一阶段，SAM掩码通过对比学习蒸馏实例特征；第二阶段，实例特征查询两个语义码本，注意力权重被熵约束推向清晰的一对一映射。这样，每个锚点不需要存储完整高维语言向量，而是通过码本获得视角无关语义。相对CAGS的局部图传播，GALA更像在实例特征和通用语言类别之间学习一个可解释的码本桥。

Scaffold-GS锚点在这里同时携带几何特征和分割特征。SAM掩码给出自监督实例划分，同掩码内的锚点特征被拉近，不同掩码被推远；双码本模块再把这些场景特定实例特征映射到更通用的语言嵌入。熵损失让注意力权重更接近一对一选择，减少同一实例在多个语义条目之间摇摆。这个设计避免了每个场景都学出一套不可迁移的实例特征。
\textbf{实现细节：} GALA：自监督实例特征场。在实例特征的学习策略上，GALA（2025.08，Group-Align-Language Attention） 与上述方法形成鲜明对比：它不依赖二维掩码的跨视角对齐，而是通过自监督对比学习蒸馏场景特定的三维实例特征场。动机。ObjectGS 和 OP2GS 等方法依赖二维掩码的跨视角对齐，但 SAM 掩码本身在不同视角间不一致；GALA 通过自监督对比学习直接在三维空间中学习实例特征，完全绕过二维掩码对齐的难题。GALA 基于 Scaffold-GS 架构，采用两阶段训练——第一阶段自监督重建场景几何和实例特征场，第二阶段渲染至二维并通过双可学习码本的交叉注意力模块将实例特征映射为泛化语言特征。线性注意力设计使仅用二维训练即可实现二维和三维无缝分割，码本与注意力权重熵损失强制一对一映射，确保实例内特征一致性并支持开放词汇查询。
\textbf{实现补充：} 实现上还要核对输入、输出、参数、特征、初始化、训练损失、正则、渲染或推理接口，以及优化中的剪枝、分裂、采样、编码或解码步骤；这些抓手决定该方法到底改了 3DGS 管线的哪一环。
\textbf{方法信息：} 仓库：\url{https://github.com/vdigpku/gala3d}，当前约 306 stars；核验时主要是项目页/README，完整代码暂未确认。
\textbf{设计动机与边界：} 相对原始 3DGS 以独立高斯和图像重建误差为主的设定，GALA 把问题落在以下可实现的改动上：GALA是一个开放词汇3DGS场景理解框架。基于Scaffold-GS，两阶段训练：(1) 场景重建和实例特征蒸馏：使用SAM生成实例掩码，通过对比学习（同掩码内特征拉近、跨掩码推远）以自监督方式蒸馏3D实例特征场。每个锚点分配几何特征和分割特征。(2) 语义码本映射：核心创新——双可学习码本交叉注意力模块。场景特定实例特征通过交叉注意力查询双码本，映射到通用化语言嵌入。码本编码视角无关的语义嵌入，每个码本条目对应一个语义类别。注意力权重的熵损失强制一对一映射，增强实例内特征一致性。线性注意力设计使得仅通过2D训练即可支持3D查询，避免了为每个高斯存储高维语言特征的内存开销。其设计目的不是单独增加一个模块，而是用上述表示、监督或渲染约束消除原流程在目标任务中的具体失配。提出双码本交叉注意力机制，将场景特定实例特征映射到通用化语言特征空间，以线性注意力设计实现仅2D训练即可支持2D和3D开放词汇查询，显著降低内存消耗。适用时还应保留论文给出的边界：码本大小需预设；对SAM掩码质量有依赖；仅验证了室内场景。
% METHOD: 2510.23930
\paragraph{PlanarGS（2025.10）\cite{2510_23930}}
PlanarGS不是用语言特征做查询，而是把视觉语言模型变成室内几何先验生成器。普通3DGS依赖光度误差优化，在白墙、地板、门板这类低纹理大平面上容易几何漂移；只看RGB误差时，很多错误深度也能解释图像。PlanarGS先用GroundedSAM按“wall”“door”“floor”等文本提示找出候选平面区域，再通过跨视角融合、深度一致性和法线验证生成更可信的平面先验掩码。

几何优化分两条线约束高斯。DUSt3R提供多视角深度和法线先验，引导高斯整体落在合理表面附近；LP3得到的平面先验则在初始化和训练中把对应高斯展平、对齐到同一平面，减少墙面和地面的起伏。训练完成后，方法从PlanarGS渲染深度图做TSDF融合，得到更适合网格重建的室内表面。相对开放词汇语义高斯，PlanarGS说明视觉语言模型也可以服务于几何：文本提示不是查询目标，而是帮助系统找到“应该是平面”的区域。
\textbf{方法信息：} 代码：\url{https://github.com/SJTU-ViSYS-team/PlanarGS}。
\textbf{设计动机与边界：} 相对原始 3DGS 以独立高斯和图像重建误差为主的设定，PlanarGS 把问题落在以下可实现的改动上：PlanarGS框架核心包含：(1) LP3管线——使用预训练视觉-语言分割模型GroundedSAM，通过文本提示（如'wall'、'door'、'floor'）检测平面区域，通过跨视角融合和利用深度/法线先验的几何验证精炼区域建议，生成高质量平面先验掩码；(2) 几何先验监督——利用DUSt3R多视角3D基础模型获取深度和法线先验，引导3DGS优化；(3) 平面先验监督——将高斯展平并利用平面先验引导初始化和优化过程，强制同平面约束。最终通过TSDF融合从PlanarGS渲染的深度图恢复高保真3D网格。关键优势是文本提示可灵活适应新场景，远优于传统启发式或专用平面检测器。其设计目的不是单独增加一个模块，而是用上述表示、监督或渲染约束消除原流程在目标任务中的具体失配。提出PlanarGS，利用预训练视觉-语言分割模型检测平面区域，设计LP3管线进行跨视角融合和几何验证，引入平面先验和几何先验监督项，在室内基准上大幅超越SOTA方法实现高保真3D表面重建。适用时还应保留论文给出的边界：依赖预训练基础模型质量；文本提示需要人工选择；对非平面复杂几何的处理能力有限。
% METHOD: 2601.04754
\paragraph{ProFuse（2026.01）\cite{2601_04754}}
ProFuse把开放词汇语义注册做成3DGS之外的上下文融合流程。OpenGaussian、LangSplat一类方法通常要训练或蒸馏语义特征场；ProFuse则先用标准重建得到几何，再用密集跨视图对应把2D语义可靠地聚到3D。它用DINOv2全局描述子检索可能重叠的图像对，再用密集匹配网络建立warp场和置信度图；SAM提供非重叠掩码，CLIP给每个掩码语言特征。

这些2D掩码不会立即写入高斯，而是先通过跨视图聚类形成3D context proposals，记录哪些视图中的掩码指向同一场景内容。每个proposal的语言特征由成员掩码加权聚合，再根据3DGS渲染时的透射率和不透明度可见性分配给相关高斯。查询时，文本特征直接与高斯语义特征比较。相对训练语言场，ProFuse的增量是把“跨视图一致的语义分组”放在注册前处理，使系统无需渲染监督语义训练，也能降低单视角CLIP特征噪声。
\textbf{实现细节：} ProFuse：渲染免监督直接注册。前三类方法在特征表示和训练范式上寻求泛化，但仍依赖可微分渲染进行语义监督。ProFuse（2026.01，Proposal-Fused Gaussian Segmentation） 则提出完全绕过渲染监督的替代方案。动机。所有基于可微分渲染的方法存在渲染-比较过程中的语义失真——$\alpha$-混合会将不同对象的特征混合在一起；ProFuse 通过直接注册（而非渲染-比较）将跨视角掩码分组为三维上下文提案，避免渲染造成的语义模糊。该方法通过密集对应引导的预注册阶段初始化高斯几何并链接跨视角分割掩码形成三维上下文提案（Context Proposals），每个提案记录推断为同一场景内容的跨视角掩码分组。随后通过上下文引导的注册阶段，利用基于透射率和不透明度的可见性加权为每个提案计算全局语言特征并分配给对应高斯。这种"直接注册"策略避免了渲染-比较的语义失真问题，每场景约 5 分钟即可完成，兼具语义保真度和效率优势。
\textbf{实现补充：} 实现上还要核对输入、输出、参数、特征、初始化、训练损失、正则、渲染或推理接口，以及优化中的剪枝、分裂、采样、编码或解码步骤；这些抓手决定该方法到底改了 3DGS 管线的哪一环。
\textbf{方法信息：} 代码：\url{https://github.com/chiou1203/ProFuse}。
\textbf{设计动机与边界：} 相对原始 3DGS 以独立高斯和图像重建误差为主的设定，ProFuse 把问题落在以下可实现的改动上：方法流程：(1)密集对应预注册——使用预训练密集匹配网络（DINOv2+金字塔卷积）计算图像对之间的密集对应，通过DINOv2全局描述子进行kNN检索候选对，利用密集warp场和置信度图建立跨视图匹配。使用SAM生成非重叠区域掩码，CLIP提取语言对齐特征。通过跨视图聚类将相关掩码分组为3D上下文提案，记录哪些视图掩码指向同一场景内容；(2)上下文引导的注册——为每个3D上下文提案计算全局语言特征（通过成员嵌入的加权聚合），使用基于透射率和不透明度的可见性权重将特征分配给对应Gaussian；(3)查询阶段——直接在3D空间中通过文本查询与per-Gaussian特征比较实现开放词汇检索。整个流程无需渲染监督的语义训练，语义融合在标准重建优化之外完成。其设计目的不是单独增加一个模块，而是用上述表示、监督或渲染约束消除原流程在目标任务中的具体失配。提出ProFuse框架，通过密集对应引导的预注册阶段初始化精确几何并构建3D上下文提案，实现无需渲染监督微调的开放词汇3DGS语义理解，语义附着约5分钟/场景，速度是SOTA的两倍。适用时还应保留论文给出的边界：依赖SAM和CLIP等预训练模型质量；密集对应匹配在大视角变化场景下可能不够鲁棒。
% METHOD: 2604.01388
\paragraph{LESV（2026.04）\cite{2604_01388}}
指出3DGS本身的重叠和非结构化会造成语义注册模糊，因此改用稀疏体素光栅化作为几何表示。稀疏体素提供不重叠、确定性的像素到3D基元映射，便于把多视角2D特征置信度感知地融合到3D；几何通过深度和法线先验优化，语义则利用AM-RADIO密集对齐特征避免层级训练。相对本章其他3DGS方法，LESV是一个路线提醒：当语义精确性优先于高斯表达自由度时，结构化体素可能比重叠高斯更好注册语言特征。
实现上，LESV用SVRaster的分层稀疏体素替换重叠高斯，并按视角相关Morton码排序体素，从渲染顺序上减少弹出和混合伪影。体素网格给每个像素贡献一个更确定的3D承载单元，因此多视角2D特征可以按可见性和高斯核权重融合到同一体素，而不是在多个重叠高斯之间分摊。几何侧再用patch-wise深度一致性约束修正底层表面；语义侧直接利用AM-RADIO的密集对齐特征，避免为不同语义层级分别训练特征场。它的关键差异不是给3DGS再加一个语义头，而是通过结构化表示减少语义渗色和射线溢出。
\textbf{方法信息：} Huawei 参与了这篇论文。
\textbf{设计动机与边界：} 相对原始 3DGS 以独立高斯和图像重建误差为主的设定，LESV 把问题落在以下可实现的改动上：核心创新是将3DGS替换为稀疏体素光栅化(SVRaster)作为场景表示。SVRaster使用分层网格中的稀疏体素，通过视角相关Morton码排序避免3DGS中的弹出伪影。基于体素的显式结构实现了2D像素到3D基元的确定性映射，通过加权体素融合从多视角2D特征图聚合鲁棒的3D特征。引入patch-wise深度一致性约束来优化底层几何，使用高斯核权重控制融合灵敏度。在语义编码方面，利用AM-RADIO的密集对齐属性实现多级语义理解，无需多层级特征场训练。整个框架消除了3DGS中的语义渗色和射线溢出伪影。其设计目的不是单独增加一个模块，而是用上述表示、监督或渲染约束消除原流程在目标任务中的具体失配。提出基于稀疏体素的3D语言场框架，解决了3DGS场景理解中的空间歧义和语义渗色问题，在开放词汇3D物体检索和点云理解基准上达到SOTA，尤其在细粒度查询上表现突出。适用时还应保留论文给出的边界：依赖单目深度和法线先验的质量；SVRaster优化仍以光度保真为主，可能导致几何断裂；对遮挡区域的处理有待提升。
% METHOD: 2606.00447
\paragraph{GeoSAM-3D（2026.05）\cite{2606_00447}}
GeoSAM-3D面向更轻量的输入条件：用户只有一段短单目视频，在某一帧点选或命名目标，希望得到绑定到3DGS场景的对象掩码。它先从单目视频重建高斯场景，再以高斯质心为节点构建图；2D提示通过SAM或视频基础模型提升到可见高斯节点，成为图上的种子。与直接按欧氏距离扩散不同，GeoSAM-3D在高斯图上使用heat-kernel测地距离传播提示，使掩码更倾向沿同一表面和连通结构扩散，而不是跨空气跳到相邻物体。

这种设计把学习压力放在冻结的图像/视频基础模型和3D几何传播上，而不是为每个场景重新训练一个语言特征场。实现上，图边权来自高斯中心距离、外观和局部连接关系；传播核把种子mask转成每个节点的软响应，再投影回视图或汇总为3D掩码。相对Click-Gaussian和SAGA，GeoSAM-3D牺牲了一部分可学习特征表达，换来单目视频场景中更直接的提示传播流程；它的主要风险也很清楚：如果单目3DGS把两层表面粘在一起，测地距离也会继承这个拓扑错误。
\textbf{实现细节：} GeoSAM-3D。GeoSAM-3D（2026.05） 延续 3DGS 在"开放词汇分割"上的基本设定，主要改变是：GeoSAM-3D：冻结图像/视频基础模型，基于单目3DGS重建和高斯质心图的测地热核传播提示。3DGS 的几何-渲染统一表示保持为底层地图。其设计动机是：提出轻量开放词汇3D分割设置：用户在一帧点击或命名物体，系统在Gaussian scene graph上按测地热核传播提示，减少相邻但不连通物体之间的泄漏。具体而言，方法先从单目视频重建3DGS场景，再以高斯质心构图，把2D基础模型特征和提示映射到图节点，通过heat-kernel geodesic距离传播mask。适用边界是：依赖单目3DGS重建质量和2D基础模型；开源/仓库状态描述性质较强，系统成熟度仍需验证。
\textbf{技术实现补充：} GeoSAM-3D。GeoSAM-3D（2026.05） 延续 3DGS 在"开放词汇分割"上的基本设定，主要改变是：GeoSAM-3D：冻结图像/视频基础模型，基于单目3DGS重建和高斯质心图的测地热核传播提示。3DGS 的几何-渲染统一表示保持为底层地图。其设计动机是：提出轻量开放词汇3D分割设置：用户在一帧点击或命名物体，系统在Gaussian scene graph上按测地热核传播提示，减少相邻但不连通物体之间的泄漏。具体而言，方法先从单目视频重建3DGS场景，再以高斯质心构图，把2D基础模型特征和提示映射到图节点，通过heat-kernel geodesic距离传播mask。适用边界是：依赖单目3DGS重建质量和2D基础模型；开源/仓库状态描述性质较强，系统成熟度仍需验证。
\textbf{方法信息：} 代码：\url{https://github.com/arunshar/geosam-3d}。
\textbf{设计动机与边界：} 相对原始 3DGS 以独立高斯和图像重建误差为主的设定，GeoSAM-3D 把问题落在以下可实现的改动上：方法先从单目视频重建3DGS场景，再以高斯质心构图，把2D基础模型特征和提示映射到图节点，通过heat-kernel geodesic距离传播mask。其设计目的不是单独增加一个模块，而是用上述表示、监督或渲染约束消除原流程在目标任务中的具体失配。提出轻量开放词汇3D分割设置：用户在一帧点击或命名物体，系统在Gaussian scene graph上按测地热核传播提示，减少相邻但不连通物体之间的泄漏。适用时还应保留论文给出的边界：依赖单目3DGS重建质量和2D基础模型；开源/仓库状态描述性质较强，系统成熟度仍需验证。

\section{3D分割}

3D分割路线通常有两种范式。第一种是“先重建再分割”：在已有3DGS上附加身份、亲和或语义标签，通过2D掩码蒸馏到3D。第二种是“先分割再重建”：在优化高斯之前就把点或掩码分配给对象，使重建过程天然带实例结构。前者灵活，后者边界更干净；两者都要处理2D掩码跨视角不一致和高斯跨边界混合。
% METHOD: 2312.00732
\paragraph{Gaussian Grouping（2023.12）\cite{2312_00732}}
Gaussian Grouping 是 3DGS 实例分割和编辑的核心锚点。原始 3DGS 的每个高斯只服务于外观和几何重建，同一物体没有显式成员关系；因此即使某个视角上能用 SAM 得到对象掩码，也很难知道三维场景中哪些高斯应该作为同一个对象一起删除、移动或修复。Gaussian Grouping 的改变是在每个高斯上增加紧凑 Identity Encoding，把“对象身份”变成可渲染、可优化的高斯属性。

训练前，方法先用 SAM 为多视角图像生成 2D 掩码，并用零样本跟踪或关联策略为跨视角的同一对象建立一致 ID。训练时，Identity Encoding 和 RGB 一样通过 3DGS rasterizer splat 到图像平面，再接一个线性分类层预测像素身份；2D 监督来自这些已关联的 SAM 掩码。仅靠 2D 掩码会漏掉遮挡区域和不可见高斯，所以方法还加入 3D spatial consistency regularization：让空间中 top-$K$ 近邻高斯的身份编码保持接近，补偿视角监督不足。实现细节上，身份编码维度很低，例如 16 维，因此对渲染速度和存储开销影响较小。

推理时，系统根据 Identity Encoding 把高斯聚成对象或 stuff 组。这个组不是后处理的 2D mask，而是三维高斯集合，因此可以直接驱动局部编辑：删除一个对象就是移除对应高斯组；对象修复是在删除后结合 2D inpainting 生成背景并投影回 3D；着色、风格迁移和场景重组则是修改或交换这些对象组的外观和位置。相对 Feature 3DGS，Gaussian Grouping 不追求通用高维语义特征，而是学习面向实例分解的紧凑身份编码；相对后续 FlashSplat，它需要训练身份场，但分组结果可以自然复用于编辑。
\textbf{实现细节：} Gaussian Grouping（2023.12） 提出了一个简洁而有效的框架，奠定了 2D-3D 特征蒸馏分割的基本范式。其核心思想是为每个高斯基元附加一个紧凑的 Identity Encoding（身份编码），通过可微渲染和多视角 SAM 掩码的对比监督学习实例归属。动机。3DGS 缺乏语义结构，一个横跨前景和背景的大高斯会让分割边界模糊不清；通过为每个高斯学习身份编码并用 2D 掩码的对比监督驱动优化，可以将语义信息注入离散高斯基元。具体地，给定 $N$ 个高斯基元，每个高斯附加一个 $D$ 维身份编码 $\mathbf{e}_i \in \mathbb{R}^D$（$D$ 通常为 16 或 32）。渲染时，身份编码通过标准的 $\alpha$-混合方式合成： 其中 $\alpha_i$ 和 $T_i$ 分别为第 $i$ 个高斯的不透明度和透射率，与颜色渲染共享相同的排序和混合逻辑。等式 相应公式 的关键洞察是：身份编码的渲染与颜色渲染完全同构，因此可以复用 3DGS 的高效光栅化管线，无需设计专门的语义渲染器。监督信号来自 SAM 在每个视角生成的 2D 掩码：同一掩码内的高斯应具有相似的身份编码，不同掩码的高斯应相互区分。这一目标通过对比损失实现。此外，Gaussian Grouping 引入 3D 空间一致性正则化，鼓励空间上邻近的高斯共享身份编码，从而弥补 2D 监督的视角不一致性。基于 Gaussian Grouping 的分组结果，该方法还支持多种下游编辑应用，包括 3D 对象移除、修复、上色和场景重组。
\textbf{锚点细化：} Gaussian Grouping：针对 3DGS 缺乏语义结构、一个横跨前景和背景的大高斯会让分割边界模糊不清的局限，用\textbf{紧凑 Identity Encoding 配合多视角 SAM 掩码对比监督}替代无语义的高斯基元，核心思想是\textbf{让身份编码与颜色共享相同的 $\alpha$-混合 渲染管线}——直觉上，既然一组高斯在 RGB 渲染中共同构成某物体的外观，用相同权重渲染其身份编码即可保证语义与外观在像素级对齐，同一 SAM 掩码内的高斯被推向相似编码、不同掩码的相互区分，配合 3D 空间一致性正则化弥补 2D 监督的视角不一致性。
\textbf{方法信息：} 代码：\url{https://github.com/lkeab/gaussian-grouping}；GitHub 关注度高，约 1019 stars；引用很多，Semantic Scholar 约 460 次。



% METHOD: 2312.00860
\paragraph{SAGA（2023.12）\cite{2312_00860}}
SAGA（Segment Any 3D Gaussians）是本章交互式3D分割的重要锚点。Gaussian Grouping已经把2D掩码提升到3D实例，但它更像离线聚类；SAGA要解决的是“用户点一下或框一下，如何在3DGS中实时返回对应对象”。它从一个已经训练好的3DGS出发，不改变颜色渲染主干，而是给每个高斯额外学习一个亲和特征。这个特征不直接表示类别，而表示“哪些高斯应属于同一个可分割目标”，因此可以服务于未知类别、部件级对象和开放式交互。

SAGA的关键是尺度门控。SAM会在2D图像中产生多粒度掩码：同一位置可能是整张桌子，也可能是桌腿。方法先用3DGS深度把2D掩码投影到3D，并沿用GARField式尺度估计，把每个掩码对应为一个物理尺度；随后用软尺度门调制高斯亲和特征的各个通道，使同一组特征在不同尺度下进入不同子空间。训练时，亲和特征被渲染成2D特征图，SAM掩码内的像素在相应尺度下被拉近，不同掩码或不同尺度的像素被推开。推理时，系统先渲染当前视角的亲和特征，用户给出2D点提示后，以该点特征为查询在3D高斯集合中做相似度匹配，得到3D目标掩码；其官方代码还提供GUI、聚类和CLIP开放词汇相关流程。相对Gaussian Grouping，SAGA的增量不是多一个分割损失，而是把“交互提示 + 物理尺度 + 3D高斯亲和场”连成实时系统。

从后续工作看，SAGA的重要性在于它把SAM的2D“分割任何物体”能力转成了一个可被3DGS实时渲染和查询的亲和特征场。DCSEG沿用这类亲和特征来生成类无关3D段，再把命名交给OpenSeg；Click-Gaussian和Seg-Wild则继续围绕多粒度掩码、全局特征候选和in-the-wild视角不一致做修正。也就是说，SAGA提供的不是类别识别器，而是一种可复用接口：先把高斯组织成可提示、可聚类的对象候选，再由文本特征、交互点或下游任务决定如何命名和使用这些候选。
\textbf{实现细节：} 与 Gaussian Grouping 同期，SAGA（2023.12，Scale-Aware Gaussian Attachment） 提出了另一种身份编码学习策略。动机。Gaussian Grouping 的身份编码粒度固定，无法同时处理从部件级到整体级的分割需求；SAGA 引入尺度感知的 soft scale gate 机制，使分割结果在不同粒度层级上保持一致性。SAGA 为每个高斯附加一个 scale-gated affinity feature，引入尺度感知的对比训练，通过一个 soft scale gate 机制调整特征通道的响应幅度，使分割结果在不同粒度层级（从部件到整体）上保持一致性。SAGA 可在 4ms 内完成一次 3D 提示分割，是首批实现 3DGS 实时交互分割的方法之一。
\textbf{锚点细化：} SAGA：针对 Gaussian Grouping 身份编码粒度固定（要么全物体、要么全无）、无法同时响应"选整张椅子"和"选椅腿"两种粒度查询的局限，用\textbf{scale-aware soft scale gate}替代固定粒度身份编码，核心思想是\textbf{为每个高斯附加尺度门控的亲和特征，通过 soft scale gate 调整特征通道响应幅度}——同一组高斯在不同尺度门控下既能聚合为整体也能细分为部件，4ms 完成一次 3D 提示分割，成为首批实现 3DGS 实时交互分割的方法之一。
\textbf{方法信息：} Huawei 参与了这篇论文；代码：\url{https://github.com/Jumpat/SegAnyGAussians}；GitHub 关注度很高，约 984 stars；引用很多，Semantic Scholar 约 204 次。



% METHOD: 2407.11793
\paragraph{Click-Gaussian（2024.07）\cite{2407_11793}}
继承SAM掩码蒸馏思路，但专门优化点击式交互。它把高斯特征分成粗粒度和细粒度两部分，分别由大物体掩码和小物体掩码对比训练；全局特征引导学习从跨视角噪声2D分割中构建候选特征簇，作为软标签平滑冲突。推理时用户点击任意视角像素，系统找到对应高斯特征并做最近邻匹配得到3D目标。相对SAGA，Click-Gaussian更强调低延迟点击响应和噪声掩码的全局特征校正。
Click-Gaussian在预训练3DGS上附加分割特征，并把特征拆成粗、细两个子空间。SAM生成的训练视角掩码按面积分组：大物体监督粗粒度特征，小物体和部件监督细粒度特征；渲染出的像素特征用对比损失约束，同一掩码内靠近、不同掩码拉开。全局特征引导学习是它相对普通掩码蒸馏的关键：方法从所有训练视角收集同一3D位置的候选2D分割特征，聚成更稳定的全局簇中心，再把簇中心作为软标签监督高斯特征。点击时，用户选中的像素会落到一批高斯上，系统用训练好的特征在3D中扩展出同一对象。
\textbf{方法信息：} 引用很多，Semantic Scholar 约 64 次。
\textbf{设计动机与边界：} 相对原始 3DGS 以独立高斯和图像重建误差为主的设定，Click-Gaussian 把问题落在以下可实现的改动上：Click-Gaussian为预训练的3D高斯附加分割特征。方法流程：(1)SAM自动生成训练视角的2D掩码，按面积排序分为粗粒度（大物体）和细粒度（小物体/部件）两级；(2)每个3D高斯附加D维特征向量f，分为f\^{}c（粗粒度部分）和f\^{}f（细粒度部分），分别对应两级掩码的特征空间；(3)特征训练：渲染像素级特征图，使用对比损失（同掩码内特征拉近，不同掩码特征推远）训练特征；(4)全局特征引导学习(GFL)：核心创新。从所有训练视角的2D分割中收集每个3D位置的全局特征候选，构建特征簇中心作为软标签，解决不同视角2D分割不一致导致的3D特征学习冲突。推理时用户点击2D像素，查询对应3D高斯的特征，通过最近邻匹配分割目标物体。整个过程约10ms。其设计目的不是单独增加一个模块，而是用上述表示、监督或渲染约束消除原流程在目标任务中的具体失配。实现10毫秒每次点击的实时3D高斯交互式分割，比之前方法快15-130倍。提出两级粒度特征场（粗+细）和全局特征引导学习(GFL)解决跨视角2D分割不一致问题，无需耗时的后处理。适用时还应保留论文给出的边界：依赖SAM的2D分割质量；非常相似物体可能混淆；特征维度有限影响细粒度区分
% METHOD: 2409.08270
\paragraph{FlashSplat（2024.09）\cite{2409_08270}}
FlashSplat 把“从 2D 掩码提升到 3D 高斯标签”写成一个可求全局最优的线性问题。Gaussian Grouping 和 SA-GA 这类方法通常要通过梯度下降学习每个高斯的身份或亲和特征，训练慢且可能陷入局部最优；FlashSplat 观察到，如果 3DGS 场景已经重建好，那么每个高斯在每个像素上的 $\alpha_i T_i$ 贡献都是常数，未知量只剩下这个高斯的标签 $P_i$。

把标签当作要渲染的属性后，某视图的 2D mask 渲染值就是 $\sum_i P_i \alpha_i T_i$。因此，给定多视图 2D 掩码，目标可以写成整数线性规划：选择每个高斯属于前景还是背景，使渲染出的 mask 尽量接近输入 mask。由于权重非负且来自固定的 $\alpha$ 混合，FlashSplat 可以把求解化简为对每个高斯累积其支持前景和背景的贡献，再做闭式或近闭式的最优分配。为了抵抗 SAM 掩码噪声，方法加入 background bias 和 softened optimal assignment，让边界附近不可靠的前景证据更难把背景高斯误吸进去。

实现上，FlashSplat 用 CUDA rasterization 先统计每个高斯对不同 mask 标签的贡献；二值分割可以快速得到标签，多类场景分割再结合迭代和动态规划处理共享高斯。得到 3D 分割后，它还用深度引导把结果渲染回新视角 mask，减少半透明高斯导致的孔洞。相对 Gaussian Grouping，FlashSplat 不学习长期可复用的身份场，但在已有掩码明确时能更快给出全局一致的高斯标签，论文报告优化约 30 秒完成，约快 50 倍。
\textbf{锚点细化：} FlashSplat：针对 Gaussian Grouping 迭代优化需数千步收敛且可能陷入局部最优的局限，用\textbf{闭合形式整数线性规划求解}替代梯度下降迭代优化，核心思想是\textbf{利用 $\alpha$-混合 权重均为已知常数的性质将分割建模为加权多数投票}——具体地，2D 掩码渲染值恰好是所有前景高斯 $\alpha$-混合贡献之和，对每个高斯累计其属于前景与背景的加权贡献，取 $\arg\max$ 即得全局最优标签，约 30 秒内完成（$50\times$ 加速）。
\textbf{实现细节：} FlashSplat：闭合形式标签分配。分割边界模糊是 3DGS 分割的顽疾，大尺寸高斯横跨对象边缘导致锯齿状分割。FlashSplat（2024.09） 从数学上重新审视了这个问题，将分割标签分配从迭代优化转化为闭合形式求解。动机。Gaussian Grouping 的迭代优化需要数千步收敛且可能陷入局部最优；FlashSplat 的核心洞察是：在已重建的 3DGS 中，$\alpha$-混合 权重均为已知常数，因此分割标签分配可转化为线性整数规划问题，通过加权多数投票直接获得全局最优解。在 FlashSplat 之前，Gaussian Grouping 等方法通过梯度下降迭代优化每个高斯的标签，需要数千步才能收敛，且可能陷入局部最优。FlashSplat 的核心洞察是：给定已重建的 3DGS 场景，$\alpha$ 值 $\alpha_i$ 和透射率 $T_i$ 均为已知常数，因此 2D 掩码渲染可写成高斯标签 $P_i \in \{0,1\}$ 的线性函数。直觉上，像素被标记为前景的渲染值恰好是所有前景高斯在该像素处的 $\alpha$-混合贡献之和。据此，FlashSplat 将分割建模为整数线性规划（ILP，Integer Linear Programming）：以平均绝对误差为目标函数 $\min_{P_i} F = \sum_v \sum_{jk} | \sum_i P_i \alpha_i T_i - M_{jk}^v |$。由于 $\alpha_i \geq 0$ 且 $\sum_i \alpha_i T_i \leq 1$，该 ILP 可通过加权多数投票直接求解。对每个高斯 $G_i$，统计其在所有掩码中属于前景的加权贡献 $A_1^i = \sum_{v,j,k} \alpha_i T_i \cdot \mathbb{I}(M_{jk}^v = 1)$ 和背景的贡献 $A_0^i$，取 $\arg\max_n A_n^i$ 即为全局最优分配。此外，FlashSplat 引入可调的背景偏置 $\gamma \in [-1,1]$（L1 归一化后执行 $P_i = \arg\max\{\hat{A}_0, \hat{A}_1\}$，其中 $\hat{A}_0 = \hat{A}_0 + \gamma$），以缓解 SAM 掩码噪声导致的边界伪影。整个优化在约 30 秒内完成（约 $50\times$ 加速），且解是全局最优的。
\textbf{方法信息：} 代码：\url{https://github.com/florinshen/flashsplat}；GitHub 关注度较高，约 262 stars；引用很多，Semantic Scholar 约 62 次。



% METHOD: 2410.17505
\paragraph{PLGS（2024.10）\cite{2410_17505}}
PLGS处理的是2D全景分割结果向3DGS提升时的鲁棒性。2D全景网络能给每张图预测语义和实例，但多视角之间会出现标签冲突、遮挡缺失和漂浮误检；直接把这些标签监督到每个高斯上，容易得到闪烁或不连续的3D语义。PLGS的路线不是重新设计一个端到端全景网络，而是把3DGS变成可承载语义场和实例场的显式地图，再用结构化约束把2D噪声压到多视角一致的3D表示中。

具体做法分语义和实例两层。语义层先构建更可靠的语义锚点，把高置信2D标签提升到3D，用这些锚点作为训练期间的平滑正则，弥补3DGS没有MLP天然连续性的缺点；随后通过自训练合并渲染出的语义掩码和噪声2D掩码，逐步扩大可靠监督。实例层则把2D实例掩码投影到3D，与有向包围盒和跨视角几何关系匹配，形成一致的实例监督。相对PanopticSplatting把查询和实例分割放进端到端模型，PLGS更像一套鲁棒lifting流程：保留现有2D全景模型的能力，同时用3DGS显式几何约束修正跨视角不一致。
\textbf{实现细节：} 3DGS 全景分割基础。PLLS（2024.10） 首次将全景分割引入 3DGS 框架。动机。在 NeRF 时代，MLP 的参数化结构天然提供平滑映射可隐式平滑 2D 掩码噪声；但 3DGS 每个高斯独立编码语义属性，噪声被忠实地记忆而非平滑，导致渲染掩码不一致。PLGS 的核心动机是解决 3DGS 不满足平滑性假设这一根本问题。在 NeRF 时代，Panoptic Lifting 和 Contrastive Lift 利用 MLP 的参数化结构天然提供特征到掩码的平滑映射，即使 2D 掩码存在跨视角噪声，MLP 也能隐式地将其平滑掉。然而 3DGS 没有这样的参数化结构。每个高斯独立编码语义属性，不存在全局平滑约束，因此即使视角仅有微小变化，2D 分割模型（如 Mask2Former）产生的噪声也会被高斯忠实地记忆而非平滑，导致渲染掩码出现明显的不一致性和伪影。直接沿用 NeRF 时代的噪声处理策略（测试时增强、分割一致性损失等）也无法获得满意结果。为此，PLGS 从结构化和降噪两个层面系统性地弥补 3DGS 的平滑性缺陷。在结构化层面，PLGS 基于 Scaffold-GS 构建全景感知的高斯模型：通过体素化操作 $\mathbf{V} = \{[\mathbf{P}/\epsilon]\} \cdot \epsilon$（其中 $\epsilon$ 为体素尺寸）将点云离散化为锚点集，每个锚点携带可学习特征 $\hat{f}_v$，再由轻量 MLP 解码器（如语义解码器 $\mathbf{F}_{\mathrm{sem}}$）为周围的 $k$ 个高斯预测语义概率 $\hat{y}$ 和不透明度等属性，从而引入锚点级别的局部平滑。
\textbf{方法信息：} 代码：\url{https://github.com/wangyuyy/plgs}。



% METHOD: 2412.10972
\paragraph{DCSEG（2024.12）\cite{2412_10972}}
DCSEG把开放集3D分割拆成“先找段、再命名”两步。第一步不关心类别，只学习高斯亲和特征：SAM提供多视角2D掩码，同一掩码内投影到3D的高斯被拉近，不同掩码或不同尺度的高斯被推开；训练完成后，HDBSCAN在亲和特征空间中把高斯聚成实例级或部件级3D段。这样得到的段更像SAGA式几何语义提案，而不是固定类别输出。

第二步才引入开放词汇语义。DCSEG把每个3D聚类重新渲染成2D二值掩码，再与OpenSeg等2D开放词汇模型的预测匹配，为3D段分配文本类别或语义标签。这个解耦让系统可以单独替换2D语义模型、SAM版本或3D聚类策略；如果开放词汇分类器更新，3D段不必重新训练。相对COS3D的协作场，DCSEG牺牲了一些端到端紧耦合性，换来更清楚的工程边界和更容易维护的开放集分割流程。
\textbf{方法信息：} 代码：\url{https://github.com/lusxvr/dcseg}。
\textbf{设计动机与边界：} 相对原始 3DGS 以独立高斯和图像重建误差为主的设定，DCSEG 把问题落在以下可实现的改动上：DCSEG采用两阶段解耦设计：阶段一（类无关掩码提案）：首先用3DGS重建场景，然后受SAGA启发使用尺度感知对比学习为每个高斯附加亲和特征。具体地，从SAM的2D实例掩码生成跨视图对应关系，通过对比损失将同一SAM掩码内的像素对应的高斯拉近、不同掩码的推远。使用HDBScan密度聚类算法将高斯特征聚类为实例级或零件级3D段。尺度超参数控制SAM粒度级别。阶段二（类感知分类）：将3D聚类光栅化回2D获得二值掩码，与2D开放词汇分割模型（如OpenSeg）的预测建立对应关系，为每个3D聚类分配语义标签。两阶段完全解耦，可独立替换2D分割模型或3D表示。其设计目的不是单独增加一个模块，而是用上述表示、监督或渲染约束消除原流程在目标任务中的具体失配。提出DCSEG，将3D分割解耦为类无关掩码提案和类感知分类两个独立阶段。设计提供即插即用接口替换不同2D/3D模块，支持零额外成本的多物体实例分割，利用丰富3D几何实现鲁棒场景理解。在mIoU和mAcc上优于可比NeRF流水线，尤其对长尾类别。适用时还应保留论文给出的边界：聚类质量受HDBScan参数影响；对SAM分割错误的传播；两阶段流程可能累积误差。
% METHOD: 2503.18073
\paragraph{PanopticSplatting（2025.03）\cite{2503_18073}}
尝试端到端训练全景3DGS。它为高斯加入语义和实例特征，并引入可学习实例查询；距离感知交叉注意力把查询与视锥体内的高斯局部关联，减少全局注意力内存。标签混合和标签变形用于提升多视角一致性、减少浮游噪声。相对多阶段“2D掩码提取、3D聚类、语义赋值”的流程，PanopticSplatting把实例查询、语义渲染和高斯优化放进一个统一模型。
PanopticSplatting的训练对象仍是可渲染高斯，但每个高斯同时带语义特征和实例特征。可学习实例查询不在全场景上做昂贵全局注意力，而是在相机视锥内通过距离感知交叉注意力读取局部高斯特征，得到实例分配；语义和实例结果再通过光栅化投到2D，与全景伪标签监督对齐。标签混合负责缓和多视角伪标签冲突，标签变形则增强2D预测在视角变化下的一致性。这样，方法把“3D实例在哪里”和“渲染出来的2D标签是否一致”放到同一个优化循环里，减少多阶段流程中前端错误逐步累积的问题。
\textbf{实现细节：} 动机。多阶段流水线中语义场和实例场的训练无法受益于彼此的监督信号，且训练时间线性叠加；PanopticSplatting 将语义和实例预测统一到端到端框架中，使两者相互促进。PanopticSplatting（2025.03） 进一步提出端到端全景重建系统：为每个高斯附加可学习的语义特征 $f_\mathrm{sem}$ 和实例特征 $f_\mathrm{ins}$，这些特征与视角无关（因为语义和实例在不同视角下应保持一致），通过 $\alpha$-混合光栅化直接渲染 2D 语义图和实例图。核心架构引入 query-guided Gaussian Segmentation 和局部交叉注意力——learnable queries 通过交叉注意力与高斯特征交互，端到端预测实例掩码和语义标签，无需后处理关联。局部交叉注意力限制在视锥体内（只关注当前视角附近的高斯），有效减少训练内存使模型可处理包含更多高斯和对象的大场景。同时设计标签混合（label blending）策略——将多视角预测的标签以加权方式混合为一致的全局标签——促进一致的 3D 分割并减少噪声浮游体，以及标签扭曲（label warping）利用深度信息将相邻视角的标签投影到当前视角作为额外监督——增强多视角一致性和分割精度。
\textbf{方法信息：} Huawei 参与了这篇论文。
\textbf{设计动机与边界：} 相对原始 3DGS 以独立高斯和图像重建误差为主的设定，PanopticSplatting 把问题落在以下可实现的改动上：PanopticSplatting的核心创新包括：(1)查询引导的高斯分割：引入可学习查询显式建模3D实例，通过距离感知交叉注意力将实例查询与高斯基元的实例特征关联，在视锥体内进行局部交叉注意力以减少训练内存；(2)标签混合：促进一致的3D分割，减少噪声浮游体；(3)标签变形：在2D预测上增强多视角一致性和分割精度。方法为每个高斯基元引入语义特征和实例特征，通过光栅化渲染2D语义和实例信息。解决了现有方法因多阶段处理导致的累积误差和手工设计组件依赖问题。其设计目的不是单独增加一个模块，而是用上述表示、监督或渲染约束消除原流程在目标任务中的具体失配。提出首个端到端的高斯溅射全景重建系统，通过查询引导的高斯分割和局部交叉注意力实现无需跨帧关联的2D实例掩码提升，解决现有方法的多阶段累积误差问题。适用时还应保留论文给出的边界：局部交叉注意力的视锥体范围需要手动设定；对极其拥挤场景的实例分割可能受限；开放词汇能力依赖于2D基础模型质量。
% METHOD: 2503.19443
\paragraph{COB-GS（2025.03）\cite{2503_19443}}
COB-GS关注的是对象边界处的“一个高斯跨两个物体”问题。普通语义高斯方法会为每个高斯学习标签或特征，但边界高斯的椭球支撑区域可能同时覆盖前景和背景；如果只在特征层面做分类，边界会糊，如果直接删掉这些高斯，渲染质量又会掉。COB-GS先给每个高斯学习连续掩码值，通过3DGS的$\alpha$合成渲染出2D掩码，并用SAM2产生的对象掩码监督这些连续标签。

然后它用掩码梯度统计找出真正跨越对象边界的高斯：若一个高斯在投影区域内同时贡献给前景和背景，且掩码变化集中在其覆盖范围内，就被视为边界模糊基元。方法不是删除它，而是把它分裂成更小的高斯，让几何支撑区域更贴近边界；分裂后再做视觉优化，修复局部纹理和透明度变化。相对Segment then Splat的对象级重建约束，COB-GS是边界级精修工具，适合接在已有3DGS或分割系统后面提升轮廓清晰度。
\textbf{实现细节：} COB-GS：边界自适应高斯分裂。FlashSplat 证明了标签分配可以一步到位，但并未解决边界高斯本身的几何歧义：横跨对象边缘的大尺寸高斯即使被正确标记，渲染出的掩码边界仍然模糊。COB-GS（2025.03，Clean Object Boundaries via Gaussian Splitting） 正是从这一局限出发，将优化焦点从"标签分配"推进到"高斯几何改造"。动机。FlashSplat 虽能一步完成标签分配，但横跨对象边缘的大尺寸高斯即使被正确标记，渲染出的掩码边界仍然模糊；COB-GS 将语义梯度异常高的"边界歧义高斯"沿语义边界分裂为两个子高斯，使高斯分布贴合对象轮廓。COB-GS 为每个高斯引入连续掩码标签 $m_i \in (0,1)$，通过 $\alpha$-混合渲染后与 SAM2 生成的 2D 掩码做监督训练。其核心操作是 boundary-adaptive Gaussian splitting（边界自适应高斯分裂）：在训练过程中统计每个高斯的掩码标签梯度 $\nabla m_i$，梯度异常高的高斯即为语义上不确定的"边界歧义高斯"，它们同时被前景和背景掩码施加了矛盾信号。COB-GS 将这些歧义高斯沿语义边界分裂为两个子高斯，分别独立学习前景和背景的掩码标签，使高斯分布贴合对象轮廓。
\textbf{方法信息：} 代码：\url{https://github.com/zestfuljx/cob-gs}。
\textbf{设计动机与边界：} 相对原始 3DGS 以独立高斯和图像重建误差为主的设定，COB-GS 把问题落在以下可实现的改动上：COB-GS是一种3DGS精化和分割方法，其核心创新包括两个部分。（1）语义引导的边界自适应高斯分裂：引入连续掩码标签m\_i属于(0,1)用于每个高斯基元，通过alpha合成渲染得到2D掩码结果。利用SAM2生成的语义掩码作为监督，通过掩码梯度统计识别跨物体边界的模糊高斯，然后将这些高斯分裂为更小的基元，使其更好地贴合物体边界。（2）视觉优化：修复因分裂操作导致的边界区域纹理退化问题，通过精炼边界结构的纹理来保持高视觉质量。与现有方法直接删除模糊边界高斯不同，COB-GS通过联合优化语义和视觉信息，使两个层级有效协作，既能清晰分割边界又保留高质量渲染。两阶段的SAM2掩码生成策略确保了监督信号的准确性。其设计目的不是单独增加一个模块，而是用上述表示、监督或渲染约束消除原流程在目标任务中的具体失配。提出了COB-GS方法，通过边界自适应高斯分裂技术和语义-视觉联合优化，解决了3DGS分割中物体边界模糊的问题，在不牺牲视觉质量的前提下显著提升了分割精度和对不准确掩码的鲁棒性。适用时还应保留论文给出的边界：依赖SAM2等预训练分割模型生成掩码监督；边界自适应分裂可能增加高斯基元数量；对极度复杂拓扑边界的处理能力有待验证。
% METHOD: 2503.22204
\paragraph{Segment then Splat（2025.03）\cite{2503_22204}}
Segment then Splat把常见流程反过来：不是先训练完整3DGS、再用SAGA或聚类方法分割，而是在建模之前就把对象ID写进高斯初始化。它先用SAM在参考帧中获得多粒度物体掩码，再用SAM2沿视角序列追踪这些掩码；COLMAP得到的初始点根据多视角掩码被分配到对象集合中，每个高斯从一开始就只服务于自己的对象。这样，后续颜色和几何优化不再是单一全局场的无差别拟合，而是在2D掩码监督下训练多个对象约束的高斯子场。

这种设计的收益是，实例边界不必等到重建结束后再从语义特征里切出来。训练完成后，每个对象集合天然对应一个3D对象，可以再附加CLIP描述符做开放词汇查询；对动态输入，方法把对象ID与4DGS式变形场结合，让对象集合随时间移动。相对Gaussian Grouping和SAGA的“重建后分割”，Segment then Splat把实例一致性提前到重建阶段，适合对象级检索、编辑和动态对象跟踪，但也更依赖前端SAM/SAM2追踪的稳定性。
\textbf{方法信息：} 代码：\url{https://github.com/luyr/segment-then-splat}。
\textbf{设计动机与边界：} 相对原始 3DGS 以独立高斯和图像重建误差为主的设定，Segment then Splat 把问题落在以下可实现的改动上：该方法包含三个核心步骤：(1)利用SAM进行网格点提示获取初始帧的物体掩码，再用SAM2进行多视角物体追踪，生成多粒度（大/中/小）的多视角掩码；(2)将COLMAP初始化的每个高斯点根据掩码分配物体ID，将场景划分为多个物体特定的高斯集合，每个高斯只对其所属物体有贡献，基于2D多视角掩码监督进行引导；(3)重建完成后，为每组高斯分配CLIP嵌入，实现开放词汇查询。对于动态场景，结合4D高斯溅射的变形场来建模场景动态。方法引入了三级粒度机制（大/中/小）以支持不同尺度的物体检索。其设计目的不是单独增加一个模块，而是用上述表示、监督或渲染约束消除原流程在目标任务中的具体失配。提出了Segment then Splat框架，颠覆了3DGS领域长期以来的「先重建后分割」范式。通过在重建前将高斯点预分配给不同物体，消除了几何和语义模糊性以及动态场景中的高斯-物体错位问题，首次在静态和动态场景中统一实现了真正的3D开放词汇分割，且无需学习额外的语言场，加速了优化过程。适用时还应保留论文给出的边界：依赖SAM/SAM2的质量，对于遮挡严重的物体可能无法准确追踪；粒度划分需要预设参数；主要针对物体级别分割，对部件级别分割支持有限
% METHOD: 2504.11024
\paragraph{Easy3D（2025.04）\cite{2504_11024}}
是通用3D交互分割基线，可用于高斯或点云环境。它用体素稀疏编码器提取彩色点云特征，把用户正负点击编码为提示向量，再由轻量Transformer解码器隐式融合点击和场景特征。负掩码组件被引入隐式融合，增强背景排除能力。相对Click-Gaussian直接在高斯特征上匹配，Easy3D更像一个独立交互分割网络，高斯可作为其三维输入来源。
Easy3D先把彩色点云体素化，并用稀疏卷积编码局部几何和颜色；用户的正负点击被编码成提示向量后，进入轻量Transformer解码器，与场景特征做隐式融合。与显式把每个点和点击距离手工拼接不同，Transformer让点击提示在全局上下文中选择相关区域。负掩码组件让模型不仅学习“点击附近是什么”，也学习“哪些区域应被排除”，因此在复杂背景中更稳。对3DGS场景而言，它可作为把高斯采样或转换为点云后的交互分割模块，而不是直接修改3DGS渲染器。
\textbf{实现细节：} 面向应用的分割。CoRe-GS 面向机器人应用提出粗到细的高斯分割：先生成分割就绪的 GS 表示，再仅优化与感兴趣区域（POI）相关的高斯，大幅减少训练时间，引入基于颜色的过滤机制去除分割引起的浮游体。LISA-3D 将语言-图像分割提升到 3D：用几何感知 LoRA 层改装 LISA 模型，通过可微重投影损失强制跨视角一致性（仅需 11.6M 可训练参数），支持零样本部署。G2P 从相反方向出发——将 3DGS 的外观感知属性（不透明度解决几何歧义、尺度实现精确边界定位）迁移到点云以增强分割判别性，无需 2D 或语言监督。Easy3D 提出高效的 3D 交互分割方法，将体素稀疏编码器与轻量 Transformer 解码器结合实现隐式点击融合，在包括 3DGS 重建场景在内的多种几何分布上泛化良好。GaussianBeV 将 3DGS 引入 BEV 感知框架：用 3D 高斯精细表示场景并 splatting 到 BEV 特征图，在线集成到单阶段模型中，在 nuScenes BEV 语义分割上取得新的 SOTA。小麦穗实例分割 展示了 3DGS 在农业领域的应用：通过 3D-to-2D 多视角投影和 Grounded SAM 实现零样本 2D 分割，再融合为 3D 实例方案训练有监督分割网络。Robust Prior-Guided Segmentation 利用 SAM-HQ 生成精确 2D 掩码，引入先验引导的标签重分配方法通过学习先验强制多视角一致性。3DGS 分类去歧 发现 GS 点云的尺度、旋转和不透明度属性有助于区分线状与平面物体以及透明/反射物体，构建了首个真实世界 GS 点云分类数据集（20 类、每类 200 对象）。%=====================================================================。
\textbf{方法信息：} Meta 参与了这篇论文。
\textbf{设计动机与边界：} 相对原始 3DGS 以独立高斯和图像重建误差为主的设定，Easy3D 把问题落在以下可实现的改动上：Easy3D采用体素编码器+Transformer解码器架构：(1) 场景预处理与编码：使用体素化稀疏编码器处理彩色点云，通过稀疏卷积高效提取特征。(2) 点击编码：将用户的正负点击编码为特征向量。(3) 隐式点击融合：使用轻量Transformer解码器实现隐式点击融合，而非显式逐点处理。创新性地引入负掩码组件（传统仅用于显式融合），增强背景理解能力。(4) 输出后处理：生成分割掩码。方法设计简洁但效果显著，首次展示了在隐式点击融合中引入负掩码组件的正面效果。其设计目的不是单独增加一个模块，而是用上述表示、监督或渲染约束消除原流程在目标任务中的具体失配。提出结合体素稀疏编码器和轻量Transformer解码器的3D交互式分割方法，在域内和域外数据集上一致超越SOTA，包括高斯泼溅生成的场景。适用时还应保留论文给出的边界：依赖点云质量；对极度稀疏或噪声严重的点云效果下降；尚未充分利用GS的颜色和透明度属性。
% METHOD: 2507.07395
\paragraph{Seg-Wild（2025.07）\cite{2507_07395}}
Seg-Wild把SAGA式交互分割放到无约束照片集合中。旅游照片和互联网图像会有光照变化、后期滤镜、行人遮挡和距离差异；直接把SAM特征蒸馏到高斯，常会把瞬态遮挡当成物体的一部分。Seg-Wild基于GS-W重建框架，为每个高斯维护内在外观、动态外观和亲和特征，把稳定场景外观与瞬态变化分开处理。

语义上，它从SAM图像编码器提取2D特征，经PCA压缩后嵌入高斯，并用CUDA splatting渲染亲和特征图。尺度自适应分割模块根据视图深度和高斯密度调整SAM分割尺度：远处或稠密区域需要更细粒度，近处完整物体则避免过切。尖刺高斯切割器再把高斯投影到2D掩码边缘，统计应被切割或平滑的异常比例，减少边界尖刺。相对SAGA，Seg-Wild的增量是把in-the-wild外观分解和分割尺度选择一起纳入交互分割系统。
\textbf{方法信息：} 代码：\url{https://github.com/sugar0725/seg-wild}。
\textbf{设计动机与边界：} 相对原始 3DGS 以独立高斯和图像重建误差为主的设定，Seg-Wild 把问题落在以下可实现的改动上：Seg-Wild基于GS-W框架扩展，将SAM特征嵌入3D高斯中构建3D特征场。每个高斯新增内在外观特征sf\_i、动态外观特征df\_i和亲和力特征af\_i。通过SAM图像编码器提取2D特征嵌入（经PCA从256维压缩至64维），利用CUDA加速的泼溅操作将3D亲和力特征投影到2D平面进行特征渲染。尺度自适应分割模块(SASM)根据每个视图的深度和密度信息自适应调整SAM分割尺度（在4到8之间），解决非约束图像集中分割不一致问题。尖刺3D高斯切割器(SGC)通过将3D高斯投影到2D平面，计算需切割的高斯比例，使用SAM掩码进行平滑处理，消除分割边缘的异常尖刺。其设计目的不是单独增加一个模块，而是用上述表示、监督或渲染约束消除原流程在目标任务中的具体失配。提出Seg-Wild，首个针对非约束图像集的基于3DGS的交互式分割方法。引入尺度自适应分割模块(SASM)根据高斯深度和密度优化SAM分割尺度，以及尖刺3D高斯切割器(SGC)平滑分割结果中的异常高斯。同时设计了针对野外场景的分割质量评估基准。适用时还应保留论文给出的边界：依赖SAM分割的质量和一致性；对极端光照变化的处理能力有限；分割粒度受SAM掩码精度制约；计算开销较大。
% METHOD: 2508.08252
\paragraph{ReferSplat（2025.08）\cite{2508_08252}}
ReferSplat把开放词汇查询推进到“指代表达分割”。普通LangSplat式查询更像类别检索，适合“chair”或“red cup”；ReferSplat处理的是“左边的椅子”“桌上的物体”这类带空间关系、属性和遮挡条件的描述。它为每个高斯增加指代特征，文本编码器同时输出词级特征和句子嵌入；分割掩码不是简单由CLIP相似度阈值得到，而是由高斯指代特征与文本词特征调制后渲染，再用伪掩码监督训练。

关键模块是位置感知跨模态交互。它让语言中的方位词和关系词参与高斯特征更新，使同一类别的多个物体可以被不同描述区分；对比学习则把满足同一指代表达的高斯拉近，把不同表达的目标推开。ReferSplat还提出R3DGS任务和Ref-LERF数据集，把评价从“开放类别是否选中”扩展到“自然语言描述是否能在3D高斯场中定位目标”。相对SAGA的点提示和GOI的语义超平面，ReferSplat更关注语言-空间关系本身。
\textbf{实现细节：} ReferSplat：指代 3DGS 新任务。动机。开放词汇分割返回所有匹配类别的实例，而用户往往只需要定位特定描述的那个对象（如"被杯子挡住的水果"）；ReferSplat 定义了指代分割新任务，通过空间感知范式显式建模 3D 位置与自然语言表达式的关系。
\textbf{实现补充：} 为避免只把它读成名称或数据点，还应把实现接口明确拆开：输入是什么，输出是高斯参数、图像、轨迹、标注还是评测指标，哪些参数、网络模块、训练损失或渲染接口真正接到原始 3DGS / Scaffold-GS / 4D-GS 管线上，以及剪枝、分裂、采样、解码、优化或评测协议在哪一步发生。若该条目是survey、benchmark 或 dataset，这里的“实现”对应论文集合、任务划分、相机/轨迹/标注组织、评价指标和基线配置，而不是虚构新的 rasterizer；若是算法条目，则对应高斯位置、尺度、旋转、不透明度、特征或 MLP 解码如何进入训练和渲染流程。
\textbf{设计动机与边界：} 相对原始 3DGS 以独立高斯和图像重建误差为主的设定，ReferSplat 把问题落在以下可实现的改动上：ReferSplat框架包含三个核心组件：(1) 3D高斯指代场构建：使用文本编码器提取词特征fw和句子嵌入fe，为3D高斯引入新属性"指代特征"。分割掩码通过指代特征与词特征的调制获得，再经渲染和伪掩码监督完成训练。(2) 位置感知跨模态交互：增强指代特征与词特征之间的交互，对齐空间和语言线索。该模块解决了现有开放词汇方法中文本查询与高斯表示在训练时缺乏交互、无法理解空间关系的缺陷。(3) 对比学习：通过强制正样本高斯之间的语义一致性来区分不同文本查询。框架支持在未见视角中分割被遮挡或不可见的目标物体。其设计目的不是单独增加一个模块，而是用上述表示、监督或渲染约束消除原流程在目标任务中的具体失配。提出新的任务指代表达3DGS分割(R3DGS)，构建了首个R3DGS数据集Ref-LERF。提出ReferSplat框架，通过3D高斯指代场和位置感知跨模态交互显式建模3D高斯点与自然语言表达的空间关系，在R3DGS和3D开放词汇分割基准上均取得SOTA性能。适用时还应保留论文给出的边界：依赖2D视觉语言模型的质量；对极其复杂的空间关系描述可能处理不佳；需要多视角2D训练图像但无需3D掩码标注。
\textbf{方法信息：} 代码：\url{https://github.com/heshuting555/refersplat}；GitHub 关注度较高，约 147 stars。



% METHOD: 2508.19699
\paragraph{LabelGS（2025.08）\cite{2508_19699}}
LabelGS选择把离散标签直接写成高斯属性。SAGA学习的是亲和特征，后续还要根据提示或聚类得到对象；LabelGS则希望每个高斯在训练过程中就获得语义标签，从而支持对象级解耦和分割。它先用DEVA生成跨视角一致掩码，并在基础3DGS上渲染伪视角来增加可见性；这些额外视角能缓解真实训练图像稀疏时标签覆盖不足的问题。

遮挡是LabelGS的核心处理对象。方法用DepthAnythingV2估计深度，在相邻分割边界处比较平均深度，判断谁遮挡谁，再生成无遮挡掩码，避免后方物体标签错误监督前景高斯。主标记阶段把2D标签提升到3D高斯，并通过随机区域采样降低训练成本；Gaussian Projection Filter处理一个高斯被多个2D标签竞争的问题，只保留投影和遮挡关系更一致的标签。相对亲和特征路线，LabelGS更强调“标签、遮挡、投影”三者在优化中的显式一致。
\textbf{实现细节：} 场景编辑驱动的分割。许多分割方法的最终目标是支持场景编辑。OMEGAS 利用 2DGS 进行 3D 目标分割，生成多视角一致的目标掩码和初步模型，再用扩散先验补充未见的遮挡部分，实现从大场景中精确重建特定对象。LabelGS 引入跨视角一致语义掩码和遮挡分析模型（Occlusion Analysis Model），避免优化中的遮挡过拟合，配合随机区域采样策略将训练效率提升 22 倍。开放词汇分割。开放词汇分割允许用户用自然语言描述任意对象，在 3D 场景中定位并分割该对象。与固定类别集的传统分割不同，开放词汇分割需要理解任意自然语言描述并将其与 3D 空间中的高斯基元关联。
\textbf{方法信息：} 代码：\url{https://github.com/garrisonz/labelgs}。
\textbf{设计动机与边界：} 相对原始 3DGS 以独立高斯和图像重建误差为主的设定，LabelGS 把问题落在以下可实现的改动上：LabelGS的核心方法包含四个关键组件：(1) 交叉视角掩码生成：使用DEVA模型生成多视角一致的语义掩码，并通过训练基础3DGS后渲染伪视角来增加视角密度；(2) 遮挡分析模型（Occlusion Analysis Model）：利用DepthAnythingV2估计深度，通过比较相邻分割掩码边界段的平均深度构建遮挡列表，生成无遮挡掩码以避免3DGS在遮挡区域过拟合；(3) 主高斯标记模型（Main Gaussian Labeling）：将2D标签提升到3D高斯，通过随机区域采样策略提高效率并实现高斯表示的有效解耦；(4) 高斯投影滤波器（Gaussian Projection Filter）：解决多个标签分配给同一3D高斯的冲突问题。整体框架将3DGS优化过程与语义分割联合进行。其设计目的不是单独增加一个模块，而是用上述表示、监督或渲染约束消除原流程在目标任务中的具体失配。提出LabelGS方法，通过交叉视角一致语义掩码、遮挡分析模型（避免遮挡过拟合）、主高斯标记模型（将2D语义先验提升到3D）和高斯投影滤波器（解决标签冲突），实现了高效的3DGS场景分割。相比Feature-3DGS在1440x1080分辨率下训练速度提升22倍，同时分割性能更优。还提出了新的3D场景分割基准。适用时还应保留论文给出的边界：依赖2D语义先验（DEVA模型），可能受限于2D分割质量；对于极端遮挡场景的处理仍有改进空间
% METHOD: 2509.23947
\paragraph{CrashSplat（2025.09）\cite{2509_23947}}
把车辆损伤2D分割提升到3DGS。通用3DGS分割往往依赖多视图一致掩码或额外训练特征，但车辆划痕、凹陷、反光区域经常只在少数视角中清楚可见，强行要求跨视角一致反而会漏掉损伤。CrashSplat因此把问题拆成“先可靠地得到2D损伤掩码，再把单视图掩码提升到3D高斯集合”：YOLO11分割网络在CarDD和VehiDE等车辆损伤数据上提供2D实例掩码，COLMAP估计相机参数并初始化车辆3DGS，随后用相机矩阵把高斯中心投影回掩码视图。
单视图提升的核心是一个学习无关的可见性筛选。方法先保留投影落入损伤掩码内部的高斯，再按深度顺序遍历这些高斯并累计不透明度贡献，逻辑接近3DGS光栅化中的Z-buffer：前景高斯的累积权重会逐步遮住更深处的背景高斯。由于车辆重建中仍可能有漂浮点或反光造成的低质量高斯，CrashSplat再用深度和不透明度的正态分布模型做统计滤波，把离群高斯从损伤集合中剔除。相对SAM式多视图对齐，这条路线牺牲了一部分通用交互能力，但换来了对“只在一个角度看得见”的细小损伤更直接的适配。
\textbf{方法信息：} 代码：\url{https://github.com/dragoschileban/crashsplat}。
\textbf{设计动机与边界：} 相对原始 3DGS 以独立高斯和图像重建误差为主的设定，CrashSplat 把问题落在以下可实现的改动上：CrashSplat包含三个阶段：(1) 2D损伤实例分割：使用YOLO11分割网络（在CarDD和VehiDE数据集上预训练）对多帧车辆图像进行推理，生成2D损伤掩码。(2) 3DGS重建：通过COLMAP计算相机参数和稀疏点云，然后优化3DGS场景表示。(3) 单视图3DGS分割：将高斯基元投影到图像平面，按深度升序遍历并累积权重（类似Z缓冲），使前景高斯逐步覆盖背景。剩余离群点通过深度和不透明度的正态分布模型统计滤波去除。该无学习方法不需要额外的场景特定训练，也不依赖多视图一致性掩码（如SAM），特别适合损伤区域可能仅在一个视图中可见的实际场景。其设计目的不是单独增加一个模块，而是用上述表示、监督或渲染约束消除原流程在目标任务中的具体失配。提出端到端自动车辆损伤检测管线，将2D损伤掩码提升到3DGS空间。创新性地提出基于Z缓冲和深度-不透明度正态分布模型的无学习单视图3DGS分割方法，特别适用于损伤仅在单一视图中清晰可见的挑战性场景。适用时还应保留论文给出的边界：依赖2D分割模型的检测精度。对于大面积损伤可能需要多视图融合。深度和不透明度的正态分布假设在极端场景中可能不成立。
% METHOD: 2510.20238
\paragraph{COS3D（2025.10）\cite{2510_20238}}
COS3D把实例分割和语言查询拆成两个互相协作的场。相对LangSplat这类单语言场基线，CLIP特征有开放词汇能力，但边界往往不如实例掩码清楚；相对OpenGaussian这类先做类无关3D段、再做语言匹配的流程，实例边界清楚，却容易把语言匹配误差累积到后处理。COS3D的输入是已重建的3DGS、多视角SAM掩码和CLIP语言特征；输出不是单一语义向量，而是每个高斯同时携带低维实例特征和高维语言特征的协作场。

训练流程分两段。第一段训练实例场：把实例特征通过3DGS渲染到训练视图，用SAM掩码构造对比损失，让同一掩码内的像素/高斯特征靠近，不同掩码分开。第二段训练实例到语言的映射：实例特征作为区分性中介，把场景中的对象结构映射到CLIP语言空间，避免直接逐像素蒸馏高维语言特征造成边界松散。推理时，文本查询先在语言场中得到3D相关性图，再自适应转成实例提示，由实例场把候选区域细化成边界更完整的分割结果。相对DCSEG的“先聚类再分类”，COS3D把实例模块、语言模块和提示细化放进同一3DGS骨干中协作训练；相对LangSplat，它用实例场补上语言相似度边界松散的问题。
\textbf{实现细节：} 开放词汇实例分割。OpenSplat3D 结合 SAM 实例掩码和对比损失训练实例特征，利用视觉-语言模型的语言嵌入实现文本驱动的实例识别，在 LERF-mask、LERF-OVS 和 ScanNet++ 验证集上展示有效性。OpenInsGaussian 针对两个核心问题：上下文感知特征提取为每个掩码补充丰富语义上下文（解决预处理中个体掩码语义不足），注意力驱动特征聚合选择性融合多视角特征以缓解对齐误差和不完整性。COS3D 提出协作场（collaborative field），包含实例场和语言场两个子场：训练时通过实例到语言的特征映射和高效两阶段策略捕获两个场的内在关系；推理时通过自适应语言到实例的提示精炼弥合两个场的特性差异。
\textbf{实现补充：} 实现上还要核对输入、输出、参数、特征、初始化、训练损失、正则、渲染或推理接口，以及优化中的剪枝、分裂、采样、编码或解码步骤；这些抓手决定该方法到底改了 3DGS 管线的哪一环。
\textbf{方法信息：} 代码：\url{https://github.com/runsong123/cos3d}。



% METHOD: 2511.18386
\paragraph{SegSplat（2025.11）\cite{2511_18386}}
把开放集分割做成前馈3DGS能力。它用SAM提取多视角掩码，并为每个掩码用CLIP计算嵌入；所有视角的CLIP特征聚类成语义记忆库，每个像素只需预测离散语义索引。冻结DepthSplat式前馈模型生成3DGS时，额外为每个高斯预测语义索引，渲染时再通过文本特征与记忆库相似度查询语义。相对逐场景优化的LangSplat，SegSplat把高维语义压缩为离散索引，并一次前向完成几何和语义预测。
SegSplat的语义不是直接把512维CLIP向量绑到高斯上。它先在输入视图中用SAM得到物体掩码，并用CLIP为掩码提特征；跨视角特征经K-Means形成一个语义记忆库，训练图像中的像素只标注为记忆库索引。DepthSplat前馈生成高斯时，除了几何和外观参数，还预测每个高斯属于哪个语义索引。渲染阶段输出的是低成本索引或索引分布，文本查询再与记忆库中的CLIP原型比较。这个设计把开放词汇能力留在记忆库中，把高斯侧负担降为离散分类，因此适合前馈重建场景。
\textbf{方法信息：} Microsoft 参与了这篇论文。
\textbf{设计动机与边界：} 相对原始 3DGS 以独立高斯和图像重建误差为主的设定，SegSplat 把问题落在以下可实现的改动上：SegSplat分三步实现：(1)语义参数预测——对每张输入图像用SAM提取物体掩码，对每个掩码用CLIP计算图像嵌入，所有视角的CLIP特征经K-Means聚类形成M个语义簇（语义记忆库），每个像素分配对应簇索引生成语义索引图；(2)基于冻结的DepthSplat模型进行前向3DGS生成，额外预测每个高斯的语义索引（M类softmax分类）；(3)渲染时同时输出外观和语义特征图，通过计算查询文本CLIP特征与记忆库的相似度实现开放词汇查询。语义索引将高维CLIP特征压缩为离散索引，存储和计算开销极小。其设计目的不是单独增加一个模块，而是用上述表示、监督或渲染约束消除原流程在目标任务中的具体失配。提出SegSplat框架，在单次前向传播中同时生成3DGS几何表示和开放词汇语义特征。通过构建紧凑语义记忆库（SAM+CLIP特征聚类）并为每个高斯分配离散语义索引，无需逐场景优化即可实现开放集语义分割，显著提升3D场景语义理解的效率。适用时还应保留论文给出的边界：语义粒度受限于SAM分割质量；聚类数量M需要预设；对遮挡区域的语义预测可能不准确
% METHOD: 2512.24763
\paragraph{UniC-Lift（2025.12）\cite{2512_24763}}
改造对比学习实例分割的后处理。传统方法学习嵌入后还要用HDBSCAN等聚类把嵌入转标签，慢且不稳定。UniC-Lift把嵌入约束在超立方体中，用对比聚类损失让同实例嵌入收敛到角点，推理时简单阈值化即可解码成离散实例标签。为改善边界，它在光栅化嵌入上线性映射后计算triplet loss，避免直接对高维嵌入硬挖掘的不稳定。相对Gaussian Grouping，它把“嵌入到标签”的过程一并纳入训练目标。
UniC-Lift把嵌入空间直接设计成可解码标签空间。每个高斯嵌入经过sigmoid限制在$[0,1]^d$中，对比聚类损失把同一实例推向同一个超立方体角点；推理时只需按0.5阈值把连续嵌入二值化，就能得到实例标签，不再依赖HDBSCAN这类后处理聚类。边界处的训练另设一条低维线性映射路径，在光栅化嵌入上做triplet hard mining，避免直接在高维空间挖难例导致不稳定。相对Contrastive-Lift式“先学嵌入、再聚类”，UniC-Lift把实例编码规则前置到训练目标中。
\textbf{方法信息：} Samsung 参与了这篇论文；代码仓库：\url{https://github.com/val-iisc/unic-lift}，核验时完整实现暂未确认。
\textbf{设计动机与边界：} 相对原始 3DGS 以独立高斯和图像重建误差为主的设定，UniC-Lift 把问题落在以下可实现的改动上：方法包含两个核心创新：(1)Embedding-to-Label过程，传统两阶段方法（如Contrastive-Lift）需要昂贵的后处理聚类（HDBSCAN）将嵌入转为标签。UniC-Lift在统一框架中，将嵌入约束在[0,1]\^{}d空间（sigmoid激活），通过对比聚类损失使同类嵌入收敛到超立方体角点，然后简单阈值化(0.5)即可将二值向量直接解码为离散实例标签。将推理复杂度从O(n log c)降到O(n)，消除后处理开销；(2)边界硬挖掘，观察到物体边界出现伪影，直接对高维嵌入应用硬挖掘不稳定。解决方案是在光栅化嵌入上应用线性层后再计算triplet loss，稳定训练并显著改善边界精度。整体在3DGS上进行端到端优化。其设计目的不是单独增加一个模块，而是用上述表示、监督或渲染约束消除原流程在目标任务中的具体失配。提出UniC-Lift统一框架，将对比学习和标签分配合并为单阶段，通过Embedding-to-Label过程直接从学习嵌入解码实例标签，结合边界硬挖掘策略，在ScanNet、Replica3D和Messy-Rooms上超越基线。适用时还应保留论文给出的边界：嵌入维度限制了可区分实例数量；对2D分割质量敏感；嵌入空间的饱和程度影响性能。
% METHOD: 2601.03510
\paragraph{G2P（2026.01）\cite{2601_03510}}
不是把点云转成高斯，而是把高斯属性反过来增强点云语义分割。准备阶段通过距离近邻把高斯与点建立对应，把高斯不透明度和尺度转移到点云；尺度分布在边界处的变化用于产生边界伪标签。主训练阶段再用预训练外观编码器蒸馏高斯外观特征到点云网络。相对3DGS内部完成分割，G2P承认高斯优化可能偏离原始点位置，因此把高斯的外观和尺度作为点云分割的辅助属性。
G2P先在高斯和点之间做近邻匹配，把不透明度和尺度这两类高斯属性转移回点云：不透明度提供外观可见性线索，尺度提供局部几何范围线索。由于3DGS在密度控制中会移动和增删高斯，方法不直接把高斯中心当作点标签，而是聚合邻域属性来增强原始点。边界伪标签来自高斯尺度分布的突变区域：物体边界处相邻高斯的尺度和支撑范围更容易发生变化。主训练阶段再蒸馏预训练外观编码器特征，使点云网络同时获得图像外观、高斯几何和边界监督。
\textbf{方法信息：} Adobe 参与了这篇论文。
\textbf{设计动机与边界：} 相对原始 3DGS 以独立高斯和图像重建误差为主的设定，G2P 把问题落在以下可实现的改动上：G2P包含两个阶段：准备阶段和主训练阶段。准备阶段：(a)Gaussian-to-Point特征增强模块——通过基于距离的最近邻搜索建立Gaussian与点之间的逐点对应关系，将Gaussian的不透明度属性（外观信息）和尺度属性（几何信息）转移到点云，解决Gaussian在自适应密度优化中偏离初始位置导致的结构错位问题；(b)基于尺度的边界提取——从聚合的Gaussian尺度分布中推导边界伪标签，利用Gaussian尺度在边界处的显著变化实现精确边界定位。主训练阶段：(c)GS外观蒸馏——利用预训练的外观编码器将Gaussian外观特征转移到点云分割网络。方法在统一3D表示中集成互补的几何和外观信息，同时保持原始点几何。其设计目的不是单独增加一个模块，而是用上述表示、监督或渲染约束消除原流程在目标任务中的具体失配。提出G2P方法，将3DGS的外观感知属性（不透明度和尺度）转移到点云以增强分割判别性，通过建立点级对应关系解决Gaussian与点几何的错位问题，利用Gaussian尺度属性实现精确边界定位，无需2D或语言监督。适用时还应保留论文给出的边界：依赖3DGS优化的质量；Gaussian属性转移的准确性受Gaussian与点对应关系质量的影响。
% METHOD: 2602.03809
\paragraph{Split and Splat（2026.02）\cite{2602_03809}}
Split and Splat与Segment then Splat都把实例结构提前到重建流程中，但它更强调零样本全景分割和独立对象重建。Split阶段先用SAM2生成2D实例掩码，并通过粗到细策略合并碎片；同时利用SfM点云和单目深度估计构建粗3DGS。这个粗几何不是最终渲染目标，而是用来做深度引导的掩码传播，使不同视角中的同一对象掩码在3D上对齐。

Splat阶段再按实例独立重建3DGS：每个对象有自己的高斯集合和边界精炼，之后按顺序合并回全局场景。方法还为每个实例计算视觉语义描述符，通常来自对象裁剪视图的特征聚合，因此重建结果不仅能渲染，也能检索、删除、复制或编辑对象。相对在单一全局高斯场上学习标签，Split and Splat用“对象先分开建，再合并”的方式降低实例边界互相污染，代价是前端掩码传播和多对象重建流程更重。
\textbf{实现细节：} 实例级交互分割。Split\&Splat 从另一个角度出发：先利用深度信息跨视角传播实例掩码生成视角一致的 2D 掩码，为每个对象独立重建再合并回场景，天然支持全景分割、对象检索和 3D 编辑等下游任务，在 ScanNetv2 上取得 SOTA。PointGauss 提出点云引导的多对象分割框架，通过基于点云分割的高斯原语解码器在 1 分钟内生成 3D 实例掩码，并贡献了 DesktopObjects-360 多对象分割数据集（27000+ 2D 掩码、全局一致标注、360\textdegree{} 覆盖）。
\textbf{实现补充：} 实现上还要核对输入、输出、参数、特征、初始化、训练损失、正则、渲染或推理接口，以及优化中的剪枝、分裂、采样、编码或解码步骤；这些抓手决定该方法到底改了 3DGS 管线的哪一环。
\textbf{方法信息：} 代码：\url{https://github.com/lttm/split_and_splat}。
\textbf{设计动机与边界：} 相对原始 3DGS 以独立高斯和图像重建误差为主的设定，Split\&amp;Splat 把问题落在以下可实现的改动上：Split阶段：(1)使用SAM2生成2D实例掩码，通过粗到细策略合并；(2)利用SfM点云和单目深度估计构建粗3DGS表示；(3)通过深度引导的掩码传播实现多视图和几何一致性精炼。Splat阶段：(1)基于精炼掩码识别实例，对每个实例独立进行3DGS重建；(2)将独立对象顺序合并到全局场景中，同时精炼边界；(3)在每个实例上计算视觉语义描述符（基于裁剪的掩码视图），嵌入到3D表示中。最终输出是带有语义描述符的高斯对象集合，支持全景分割、对象检索和3D编辑等下游任务。核心创新在于显式实例建模而非在已有3DGS上拟合语义特征。其设计目的不是单独增加一个模块，而是用上述表示、监督或渲染约束消除原流程在目标任务中的具体失配。提出Split\&Splat框架，通过先分割后重建的两阶段策略实现全景场景重建。Split阶段利用深度信息跨视图传播和精炼实例掩码确保多视图一致性；Splat阶段对每个实例独立重建后合并并嵌入语义描述符，在ScanNetv2分割基准上达到SOTA性能。适用时还应保留论文给出的边界：依赖SAM2等基础模型质量；独立重建每个实例的计算开销较大；对遮挡严重的实例分割仍有限制。
% METHOD: 2602.10173
\paragraph{ArtisanGS（2026.02）\cite{2602_10173}}
是面向实际编辑的交互工具集。它把2D分割网络的单次点击或2D掩码传播到3DGS选择，并允许用户在2D和3D模式之间切换手动修正；选中的高斯区域可作为局部编辑约束，交给视频扩散模型修改。相对自动分割方法，ArtisanGS的重点不是提出新的语义场，而是把AI分割、人工纠错和局部生成编辑放进同一套高斯选择界面。
ArtisanGS的分割不是一次性自动输出，而是一个可交互的选择循环：用户在任意视图点击或给出2D掩码后，系统借助预训练2D分割模型把选择传播到对应的3D高斯集合；如果传播错误，用户可以在2D视图和3D视图之间切换，手动增删区域并重新约束选择。被选中的高斯集合随后成为局部编辑区域，视频扩散模型只在该区域内改写外观或内容。它的贡献更接近工具链：把AI分割、人工修正和生成式局部编辑统一到3DGS资产操作中，而不是训练一个新的通用语义场。
\textbf{实现细节：} 实时交互分割。Click-Gaussian 学习两级粒度（粗粒度和细粒度）的可区分特征场，通过 Global Feature-guided Learning（GFL）从跨视角的噪声 2D 分段中构建全局特征候选簇，平滑训练中的噪声，实现 10ms/次的点击分割速度（比之前方法快 15--130 倍）。GaussianCut 将场景建模为图结构，利用图割算法最小化能量函数来划分前景和背景高斯，支持点击、涂画和文本等多种交互方式，无需额外的分割感知训练。ArtisanGS 提供围绕高斯选择和分割的交互式工具套件：快速 AI 驱动的 2D 选择到 3DGS 选择的传播（允许用户干预修正错误），辅以灵活的手动选择和分割工具，可实现任何非结构化 3DGS 场景的任意二元分割。
\textbf{实现补充：} 实现上还要核对输入数据、输出指标、基线配置、评测协议、系统接口和适用边界；若调用 3DGS 渲染或重建结果，也应说明相机参数、训练设置、优化流程和指标计算如何对齐。
\textbf{方法信息：} NVIDIA 参与了这篇论文。
\textbf{设计动机与边界：} 相对原始 3DGS 以独立高斯和图像重建误差为主的设定，ArtisanGS 把问题落在以下可实现的改动上：ArtisanGS工具套件包含：1)AI驱动的3D分割：利用预训练2D分割网络将用户单次点击或2D掩码从任意视图传播到3DGS选择，允许用户干预和纠错；2)灵活的手动选择工具：支持2D和3D选择模式的自由切换，实现精确的区域选择和二值分割；3)局部编辑应用：基于自定义视频扩散模型的用户引导局部编辑，用户通过选择工具直接控制AI可修改区域。所有工具无需额外优化即可用于任何野外捕获的3DGS场景。其设计目的不是单独增加一个模块，而是用上述表示、监督或渲染约束消除原流程在目标任务中的具体失配。提出快速AI驱动的2D选择掩码到3DGS选择的传播方法，结合灵活的手动选择和分割工具，实现非结构化3DGS场景的几乎任意二值分割。开发基于自定义视频扩散模型的用户引导局部编辑方法。适用时还应保留论文给出的边界：AI分割在高度遮挡区域可能失败；视频扩散模型编辑质量受限于模型能力；大规模场景的交互性能可能受限。
% METHOD: 2602.17134
\paragraph{B3-Seg（2026.02）\cite{2602_17134}}
给FlashSplat的最优分割提供贝叶斯解释并扩展到主动视角选择。它为每个高斯的前景概率设Beta-Bernoulli后验，给定一个视角掩码后，根据该高斯对像素的责任累积成功和失败计数并更新后验。对称先验下，最优标签退化为FlashSplat式决策规则；同时，方法可解析计算候选视角的期望信息增益，贪婪选择最能减少标签不确定性的视角。相对FlashSplat一次性求标签，B3-Seg把掩码提升写成可连续更新的不确定性估计。
B$^3$-Seg把每个高斯是否属于目标写成 Bernoulli-Beta 模型：$y_i|p_i \sim \mathrm{Bernoulli}(p_i)$，$p_i \sim \mathrm{Beta}(a_i,b_i)$。当一个视角给出目标掩码后，方法按该高斯对掩码像素的渲染责任累计成功计数 $e_{i,1}$ 和失败计数 $e_{i,0}$，并用共轭更新把后验变为 $\mathrm{Beta}(a_i+e_{i,1},b_i+e_{i,0})$。如果先验对称，贝叶斯最优标签会退化为FlashSplat的决策规则；如果继续采集新视角，则可用期望信息增益选择最能降低后验不确定性的下一视角。这样，FlashSplat的经验规则被提升为可解释、可主动观测的分割过程。
\textbf{实现细节：} B$^{}3$-Seg：贝叶斯最优视角采样。B$^{}3$-Seg（2026.02） 从贝叶斯视角重新定义 3DGS 分割：将分割建模为序列 Beta-Bernoulli 贝叶斯更新，通过解析式 Expected Information Gain（EIG，期望信息增益）主动选择下一个最优视角，保证自适应单调性和次模性，产生 $(1-1/e)$ 近似的最优视角采样策略。动机。现有分割方法假设视角已给定，但并非所有视角都对分割同等有用——某些视角包含更多辨别性信息；B$^{}3$-Seg 主动选择信息增益最大的视角进行标注，以最少的视角获得最准确的分割。与 Click-Gaussian 需要预先训练特征场不同，B$^{}3$-Seg 完全不需要相机参数和训练过程，但以多视角采集为代价换取了无需训练的优势。
\textbf{方法信息：} Sony 参与了这篇论文。
\textbf{设计动机与边界：} 相对原始 3DGS 以独立高斯和图像重建误差为主的设定，B$^3$-Seg 把问题落在以下可实现的改动上：方法将3DGS分割重新表述为概率贝叶斯更新。对每个高斯g\_i，设置Bernoulli-Beta先验/后验：y\_i|p\_i ~ Bernoulli(p\_i), p\_i ~ Beta(a\_i, b\_i)。给定渲染图像和物体掩码，计算每个高斯的成功计数e\_i,1和失败计数e\_i,0（基于像素责任）。通过Beta-Bernoulli共轭性进行后验更新：Beta(a\_i, b\_i) -> Beta(a\_i + e\_i,1, b\_i + e\_i,0)。证明了对称先验下贝叶斯最优标签退化为FlashSplat的决策规则。主动相机采样：通过解析计算期望信息增益IG(v) = sum\_i [H(Beta(a\_i,b\_i)) - H(Beta(a\_i+e\_i,1, b\_i+e\_i,0))]，选择信息增益最大的候选视角。证明了EIG的非负性和边际收益递减，确认自适应单调性和子模性，提供贪婪(1-1/e)近似保证。每个选择视角使用Grounding DINO + SAM2 + CLIP重排序获取掩码。其设计目的不是单独增加一个模块，而是用上述表示、监督或渲染约束消除原流程在目标任务中的具体失配。提出B3-Seg，在无相机视角、无训练条件下实现几秒内完成开放词汇3DGS分割。将分割重新表述为序列Beta-Bernoulli贝叶斯更新，通过解析期望信息增益(EIG)主动选择下一视角。证明了EIG的自适应单调性和子模性，提供贪婪(1-1/e)近似保证。适用时还应保留论文给出的边界：依赖Grounding DINO + SAM2的分割质量；对极端细小或透明物体的分割可能困难；理论保证依赖于Beta-Bernoulli模型的假设。
% METHOD: 2604.14706
\paragraph{NG-GS（2026.04）\cite{2604_14706}}
NG-GS和COB-GS都盯住边界问题，但处理方式不同。COB-GS把跨边界高斯分裂成更小基元；NG-GS认为离散高斯在语义边界处天然会产生aliasing，因此引入NeRF式连续场来修正。它先统计多视角掩码对每个高斯的标签方差：如果同一高斯在不同视图或相邻区域中反复被前景/背景争夺，就被识别为边界高斯。

这些边界高斯被用来构建连续细化场。方法通过RBF插值把离散边界信号扩展为空间连续特征，再用多分辨率哈希编码表达多尺度变化；轻量NeRF模块接收这些特征并输出更平滑的语义结果。训练时，语义对齐损失保证NeRF细化结果与3DGS语义一致，空间连续性损失则抑制跨视角边界抖动。相对直接改变高斯几何，NG-GS把3DGS作为主表示，把NeRF只用于边界连续性补偿。
\textbf{实现细节：} NG-GS：NeRF-高斯混合弥补语义空洞。上述方法均依赖 3DGS 的离散点基表示，当高斯分布稀疏或边界区域覆盖不足时，语义特征场会出现不连续的空洞。NG-GS（2026.04，NeRF-Gaussian Splatting） 引入轻量 NeRF 模块弥补这一缺陷。动机。纯 3DGS 的离散点基表示在高斯分布稀疏或边界覆盖不足时，语义特征场会出现不连续的空洞；NG-GS 引入轻量 NeRF 模块通过 RBF 插值构建空间连续的特征场，弥补离散高斯之间的语义间隙。首先通过掩码方差分析自动识别边界歧义高斯，统计每个高斯在多视角掩码中被标记为前景和背景的频率差异，差异大的即为边界高斯。然后以这些边界高斯为锚点，用径向基函数（RBF，Radial Basis Function）插值在其周围构建空间连续的特征场，使离散高斯之间获得平滑的语义过渡。同时通过多分辨率哈希编码（MRHE，Multi-Resolution Hash Encoding）为 RBF 提供多尺度空间输入，增强表示容量。最后通过联合优化策略对齐 3DGS 和 NeRF 的输出，利用对齐损失和空间连续性损失确保分割兼顾高频细节和跨视角平滑。
\textbf{方法信息：} 代码：\url{https://github.com/bjtu-kd3d/ng-gs}。



% METHOD: 2604.22439
\paragraph{NRGS（2026.04）\cite{2604_22439}}
把训练无关的语义提升和后续神经正则分开。相对LangSplat的逐场景语义蒸馏，OccamLGS一类方法可以把多视角DINO或CLIP特征按渲染贡献加权平均到高斯上，速度快，但同一物体在不同视角、遮挡和边界处的2D特征并不总是一致，直接平均会得到带噪语义场。NRGS继承OccamLGS的高效提升作为基线，同时为每个高斯计算跨视角特征方差：低方差说明多视角观察较一致，高方差通常出现在边界、遮挡或几何不准的位置。
在这个初始语义场之上，NRGS训练一个轻量条件MLP做3D空间内的语义去噪。MLP的输入不是图像，而是高斯自身的位置、尺度、旋转、不透明度和球谐外观等属性，并可用语义粒度作为条件；输出是被修正后的语义特征。训练损失按方差加权，让低方差高斯作为更可靠教师，高方差高斯少主导参数更新。这样，它没有把成本放到更复杂的2D预处理或逐视图高维特征渲染上，而是利用3DGS已经学到的几何和外观一致性去正则化语义。公开实现的使用流程也对应这三步：先训练RGB高斯模型，再运行特征提取脚本生成不同粒度的语义checkpoint，最后训练和测试MLP正则器。
\textbf{实现细节：} NRGS：方差感知特征修正。NRGS（2026.04，Noise-Reduced Gaussian Splatting） 从正则化角度解决同一问题：跨视角 2D 语义特征的不一致性导致直接提升到 3D 后语义场噪声严重。动机。LangSplat 的迭代训练蒸馏耗时且可能过拟合，OccamLGS 的加权平均虽快但无法区分可靠与不可靠视角的信号；NRGS 采用训练无关的特征聚合加轻量 MLP 修正，以方差加权优先从多视角信号可靠的区域学习。不同于 LangSplat 需要迭代训练蒸馏特征、OccamLGS 仅做加权平均，NRGS 采用训练无关的特征聚合方式先将 2D 特征直接映射到 3D 高斯，然后引入轻量 variance-aware conditional MLP 利用高斯的固有属性（位置、协方差、颜色、不透明度）作为条件输入来修正 3D 空间中的语义误差。该 MLP 以方差加权为目标函数，优先从多视角信号可靠（方差低）的区域学习，抑制来自不一致视角的噪声信号，无需额外预处理且计算开销极小。
\textbf{方法信息：} 代码：\url{https://github.com/Gintoki-jpg/NRGS_Code}。



% METHOD: 2605.23287
\paragraph{LangFlash（2026.05）\cite{2605_23287}}
从稀疏无位姿图像前馈生成语言高斯。相对NoPoSplat这个几何前馈锚点，LangFlash新增的是语言语义分支；相对SegSplat和LSM这类前馈语义3DGS路线，它进一步取消了稳定相机位姿或重型语义预处理的默认条件。输入是少量未定姿图像，系统要一次前向给出几何、外观和语言特征。它冻结NoPoSplat几何骨干，把编码器的多视图特征送入语义分组模块，再用语言特征聚合模块为分组生成字典式语义表示。
方法的关键不是直接回归每个高斯的高维语言向量，而是先学习一组场景语义组。Siamese ViT解码器把多视图特征解成多尺度密集语义图；DETR式查询在这些密集特征上做自注意力和交叉注意力，判断哪些查询对应真实语义组，并为每个空间位置输出组概率。落到高斯上时，一个高斯不再存完整CLIP/LSeg向量，而是存它对若干语义组的局部权重；这些组共享一个全局语义字典，字典原子由预训练VLM特征加权聚合，再经Transformer块细化。这样，语义信息随高斯一起渲染，但存储和训练压力由“高维逐高斯特征”变成“字典加稀疏权重”。代码仓库的配置也体现了这个设计：几何侧沿用pose-free CroCo/NoPoSplat式编码和DPT高斯头，语义侧另设分组、字典和语言聚合模块。为支撑监督，论文还给RealEstate10K补充了密集且时间连贯的语义标注。
\textbf{实现细节：} 语言特征场的构建。LangSplat 是这一方向的先驱，为每个高斯附加 CLIP 特征并通过 3D 语言高斯光栅化实现开放词汇查询。后续工作从效率和精度两个方向改进。Ilov3Splat 引入实例级开放词汇理解：利用多分辨率哈希嵌入高效编码语言对齐的 CLIP 特征，通过 SAM 掩码上的对比损失训练实例特征场，在推理时通过两阶段 3D 聚类检索相关高斯组。OpenGaFF 提出 Gaussian Feature Field，将语义建模为高斯几何和外观的连续函数——通过以几何结构为条件增强语义预测，提升相似结构间的空间一致性。该方法还引入结构化码本作为共享语义原语集，配合码本引导注意力机制检索语言特征，降低对象内特征方差。SCOUP 从效率角度出发，提出稀疏码本提升策略：将语言表示学习与 3DGS 优化解耦——在 2D 图像区域学习码本系数后通过加权稀疏聚合提升到 3D 高斯，Top-$K$ 过滤提取每个高斯的主导多视角系数，实现 $400\times$ 训练加速和 $3\times$ 训练内存效率提升。LangFlash 提出前馈式 3D 语言高斯光栅化，从稀疏无位姿图像单次前向预测几何和语义，引入稀疏语义编码方案结合全局语义字典和逐原语局部权重，支持低延迟 3D 重建和语言一致性场景理解。Lee et al. 重新定义开放词汇分割为 3D 体分割问题：直接在 3D 点上监督语言嵌入场训练（而非 2D 像素），再将学到的语言场迁移到 3DGS，首次在不牺牲训练时间或精度的前提下实现实时渲染速度，并提出 3D 查询和评估协议。
\textbf{技术实现补充：} 语言特征场的构建。LangSplat 是这一方向的先驱，为每个高斯附加 CLIP 特征并通过 3D 语言高斯光栅化实现开放词汇查询。后续工作从效率和精度两个方向改进。Ilov3Splat 引入实例级开放词汇理解：利用多分辨率哈希嵌入高效编码语言对齐的 CLIP 特征，通过 SAM 掩码上的对比损失训练实例特征场，在推理时通过两阶段 3D 聚类检索相关高斯组。OpenGaFF 提出 Gaussian Feature Field，将语义建模为高斯几何和外观的连续函数——通过以几何结构为条件增强语义预测，提升相似结构间的空间一致性。该方法还引入结构化码本作为共享语义原语集，配合码本引导注意力机制检索语言特征，降低对象内特征方差。SCOUP 从效率角度出发，提出稀疏码本提升策略：将语言表示学习与 3DGS 优化解耦——在 2D 图像区域学习码本系数后通过加权稀疏聚合提升到 3D 高斯，Top-$K$ 过滤提取每个高斯的主导多视角系数，实现 $400\times$ 训练加速和 $3\times$ 训练内存效率提升。LangFlash 提出前馈式 3D 语言高斯光栅化，从稀疏无位姿图像单次前向预测几何和语义，引入稀疏语义编码方案结合全局语义字典和逐原语局部权重，支持低延迟 3D 重建和语言一致性场景理解。Lee et al. 重新定义开放词汇分割为 3D 体分割问题：直接在 3D 点上监督语言嵌入场训练（而非 2D 像素），再将学到的语言场迁移到 3DGS，首次在不牺牲训练时间或精度的前提下实现实时渲染速度，并提出 3D 查询和评估协议。
\textbf{方法信息：} 代码：\url{https://github.com/liylo/LangFlash}。
\textbf{设计动机与边界：} 相对原始 3DGS 以独立高斯和图像重建误差为主的设定，LangFlash 把问题落在以下可实现的改动上：LangFlash基于冻结的NoPosplat网络构建：(1)语义分组：Siamese ViT解码器产生多视图多尺度密集语义图，DETR式可学习查询通过交叉注意力获取3D场景知识，预测每个空间位置的分组概率，过滤后得到K个有效语义组及归一化权重。(2)稀疏语义编码：每个高斯基元通过语义权重获取所属组的语言特征，使用全局语义字典+逐基元局部权重的方式压缩表示。(3)语言特征聚合：利用预训练VLM提取语言特征图，按语义权重聚合为每组的初始特征，再通过transformer块更新。为支持大规模训练，对RealEstate10k数据集补充了连贯密集的语义标注。其设计目的不是单独增加一个模块，而是用上述表示、监督或渲染约束消除原流程在目标任务中的具体失配。首次实现无位姿、语言感知的前向传播3DGS重建，通过稀疏语义编码（全局语义字典+逐基元权重）大幅降低表示复杂度，同时保持高质量语义重建。适用时还应保留论文给出的边界：语义分组的数量K需要预设上限；依赖预训练骨干的质量；语义分辨率受分组粒度限制

\section{特征蒸馏与表示}

如果说前两节解决“怎么查询和怎么分割”，本节更底层：语义、语言、视频和跨模态特征到底怎样作为高斯属性被学习、渲染、反投影或预训练。这里的锚点是Feature 3DGS，它把3DGS从RGB渲染器扩展为通用特征渲染器。
% METHOD: 2312.03203
\paragraph{Feature 3DGS（2023.12）\cite{2312_03203}}
是语义特征高斯的基础锚点。它为每个高斯新增任意维语义特征向量，并实现并行N维高斯光栅化器，把这些特征像颜色一样渲染成2D特征图。训练时，高斯位置和协方差与RGB共享，但RGB和特征图解耦监督，避免高维特征破坏外观重建；2D监督来自SAM、CLIP-LSeg等基础模型。这样，点选、框选、语言分割和3D编辑都可以通过渲染出的特征图实现。后续几乎所有语言和分割高斯，都可以看作在这个“可渲染特征属性”上做压缩、一致性或任务化改造。
\textbf{实现细节：} 将 CLIP 等视觉-语言模型的二维语义特征提升到三维空间，是实现开放词汇理解的基本策略。Feature 3DGS（2023.12） 奠定了这一方向的基础。动机。SAM、CLIP 等二维基础模型能提取强大的语义特征，但仅存在于二维像素空间中，缺乏跨视角一致性；Feature 3DGS 通过将这些二维特征蒸馏到三维高斯基元中，使每个高斯拥有与视角无关的语义属性。本方法修改了 3DGS 的高斯基元属性——为每个高斯基元附加一个任意维度的语义特征向量 $f \in \mathbb{R}^N$，取代了仅渲染颜色的传统设定——而 3DGS 的几何表示（位置、协方差、不透明度）保持不变。标准 3DGS 中，每个高斯基元的可优化属性仅包括位置 $\mu \in \mathbb{R}^{}3$、旋转四元数 $q \in \mathbb{R}^{}4$、缩放因子 $s \in \mathbb{R}^{}3$、不透明度 $\alpha \in \mathbb{R}$ 以及球谐系数（最高 3 阶，用于建模视角相关颜色 $c \in \mathbb{R}^{}3$）；Feature 3DGS 在此基础上增加语义特征 $f \in \mathbb{R}^N$，其中 $N$ 可以是任意维度，从而将属性集扩展为 $\Theta_i = \{\mu_i, q_i, s_i, \alpha_i, c_i, f_i\}$。特征渲染的核心思想是复用 3DGS 原有的 $\alpha$-混合 机制：正如 RGB 图像通过逐高斯加权求和得到像素颜色，语义特征图也用完全相同的权重渲染。直觉上，如果一组高斯基元在 RGB 渲染中共同构成了某个物体的外观，那么用同一组权重渲染语义特征，就能保证语义信息与外观信息在空间上对齐。具体而言，像素 $p$ 处的渲染颜色和语义特征分别由以下公式计算： 其中 $T_i = \prod_{j=1}^{i-1}(1 - \alpha_j)$ 是透射率，$N$ 是沿射线按深度排序后与该像素重叠的高斯基元集合。两个公式共享完全相同的 $\alpha$-混合 权重 $\alpha_i T_i$，区别仅在于前者混合 RGB 颜色 $c_i$，后者混合语义特征 $f_i$——这意味着语义特征与 RGB 图像具有完全一致的像素级空间对齐关系。
\textbf{方法信息：} 引用很多，Semantic Scholar 约 451 次。
\textbf{设计动机与边界：} 相对原始 3DGS 以独立高斯和图像重建误差为主的设定，Feature 3DGS 把问题落在以下可实现的改动上：Feature 3DGS将3DGS从纯辐射场渲染扩展到语义特征场渲染。核心方法包含：（1）高维语义特征渲染：为每个3D高斯体增加语义特征向量（除颜色外）。提出并行N维高斯光栅化器，可渲染任意维度的语义特征图。关键挑战是RGB图像与特征图的空间分辨率和通道数不一致：采用共享的3D高斯属性（位置、协方差），但分别渲染RGB和特征图，通过解耦的训练策略避免互相干扰；（2）特征场蒸馏：利用2D基础模型（SAM用于可提示分割，CLIP-LSeg用于语言驱动分割）提取2D语义特征作为监督信号，训练3D高斯体的语义特征使其在渲染后与2D特征匹配；（3）应用支持：基于渲染的特征图实现多种3D语义任务——新视角语义分割（将渲染特征图通过分割头得到分割结果）、语言引导编辑（通过CLIP特征匹配实现文本驱动的3D场景编辑）、SAM点选/框选提示（首次在辐射场中实现点和边界框提示的3D分割）。初始化仍使用SfM点云。其设计目的不是单独增加一个模块，而是用上述表示、监督或渲染约束消除原流程在目标任务中的具体失配。提出Feature 3DGS，首个基于3DGS框架的特征场蒸馏技术。通过并行N维高斯光栅化器渲染任意维度语义特征图，解决RGB图像与特征图间空间分辨率和通道一致性差异的挑战。实现新视角语义分割、语言引导编辑和SAM点选/框选提示。适用时还应保留论文给出的边界：特征场质量依赖于2D基础模型的性能；高维特征渲染增加内存消耗；对遮挡区域的3D一致性可能不足；不支持动态场景
% METHOD: 2410.10719
\paragraph{4-LEGS（2024.10）\cite{2410_10719}}
把语言嵌入扩展到4D动态高斯。它在Dynamic 3DGS式表示中为高斯增加语义特征属性，使用ViCLIP从视频帧序列提取像素对齐时空特征，并通过多尺度图像金字塔捕获时间差异。每个时间步的3DGS独立优化中心和旋转，但其余高斯属性保持时序一致；训练后可用文本查询渲染时空相关图，定位包含查询事件的时间段和区域。相对静态语言高斯，4-LEGS让语言特征随动态场景的时间演化对齐。
4-LEGS在Dynamic 3DGS的高斯属性中加入语义特征，并用ViCLIP从视频序列提取带时间上下文的像素对齐特征作为监督。训练时，每个时间步允许高斯中心和旋转随动态场景变化，但外观、尺度等属性保持一致，以免语言特征被逐帧漂移打散；多尺度图像金字塔帮助模型同时看到局部动作和较长时序变化。查询时，文本嵌入与4D语言高斯渲染出的时空特征图匹配，得到目标事件出现的时间段和空间区域。它相当于把静态开放词汇3DGS扩展成可问“什么时候发生”的4D表示。
\textbf{方法信息：} Google / DeepMind 参与了这篇论文。
\textbf{设计动机与边界：} 相对原始 3DGS 以独立高斯和图像重建误差为主的设定，4-LEGS 把问题落在以下可实现的改动上：方法扩展Dynamic 3DGS表示，为每个高斯添加语义特征属性。核心流程：（1）时空特征提取：使用ViCLIP（基于CLIP的视频编码器）处理输入视频的时间序列帧，通过图像金字塔提取多尺度像素对齐特征，捕获时序差异。（2）4D语言场优化：将提取的时空特征作为监督信号，优化每个高斯的语义特征属性。每个时间步独立优化3DGS，除高斯中心和旋转外所有参数强制时序一致性。（3）开放词汇查询：通过文本查询训练的4D特征场，渲染时空相关图，预测包含查询描述的时间段和像素级概率。其设计目的不是单独增加一个模块，而是用上述表示、监督或渲染约束消除原流程在目标任务中的具体失配。提出4-LEGS，首个将语言嵌入扩展到4D动态高斯splatting的方法，利用ViCLIP视频-文本模型实现基于文本的时空事件定位，显著优于2D时空定位和静态3D语言嵌入技术。适用时还应保留论文给出的边界：依赖预训练ViCLIP模型的语义理解能力；对长视频序列的计算成本较高；密集标注评估数据有限
% METHOD: 2411.15193
\paragraph{Gradient-Weighted Feature Back-Projection（2024.11）\cite{2411_15193}}
提供不训练特征场的快速替代方案。3DGS渲染中，某个高斯颜色对像素颜色的梯度等于该高斯对像素的alpha和透射率权重；因此，可以用这些渲染贡献权重把2D特征图加权反投影到高斯上。每个高斯的特征是所有视角中它贡献过的像素特征的归一化平均。相对Feature 3DGS，它取消语义特征优化，把2D基础模型特征一次性聚合到3D，适合快速分割、affordance迁移和身份编码。
它的实现可以直接从渲染方程看懂：像素颜色$C(x,y)$是沿射线排序后的高斯颜色加权和，第$k$个高斯的权重就是不透明度和前向透射率的乘积$\alpha_k T_k$。对颜色$c_k$求梯度时，这个权重也正是该高斯对像素的渲染贡献。方法把SAM、CLIP或其他2D基础模型得到的特征$F_{2D}$按$\alpha_k T_k$累加到对应高斯上，得到$f_k=\sum F_{2D}(x,y,n)\alpha_kT_k / \sum \alpha_kT_k$。因此，高斯只从自己真正影响过的像素接收特征，遮挡和半透明贡献由渲染权重自然处理。它不是训练一个新语言场，而是利用现有3DGS的可微渲染权重做一次特征聚合。
\textbf{方法信息：} 代码：\url{https://github.com/jojijoseph/3dgs-gradient-backprojection}。
\textbf{设计动机与边界：} 相对原始 3DGS 以独立高斯和图像重建误差为主的设定，Gradient-Weighted Feature Back-Projection 把问题落在以下可实现的改动上：核心方法基于对3DGS渲染方程的梯度分析：渲染像素颜色C(x,y)是各高斯体颜色的加权和，权重为opacity*transmittance。对第k个高斯体颜色c\_k的梯度dC/dc\_k恰好等于该高斯体在渲染中的权重alpha\_k*T\_k。利用这一性质，定义特征反投影方程：f\_k = sum(F\_2D(x,y,n)*alpha\_k*T\_k) / sum(alpha\_k*T\_k)，其中F\_2D是2D特征图（来自SAM、CLIP等），求和覆盖所有视角和像素。这相当于按高斯体的渲染影响力加权聚合2D特征。支持expected和accumulated两种反投影模式，经L2归一化后等价。应用包括：3D分割（余弦相似度+阈值）、affordance迁移、身份编码等，均在3D空间直接操作。其设计目的不是单独增加一个模块，而是用上述表示、监督或渲染约束消除原流程在目标任务中的具体失配。提出Gradient-Weighted Feature Back-Projection，一种无需训练的3DGS特征场构建方法。核心洞察：3DGS渲染颜色是各高斯体颜色的加权和，其对某高斯体颜色的梯度恰好等于该高斯体的渲染权重。利用这一性质，将2D特征（如SAM、CLIP特征）按渲染权重反投影到3D高斯体上，实现高质量的2D和3D分割。方法快速、可扩展、无需后处理。适用时还应保留论文给出的边界：依赖预训练3DGS模型的质量；2D特征提取器的质量直接影响3D特征质量；对于被遮挡区域的高斯体，特征可能不够准确。
% METHOD: 2503.18052
\paragraph{SceneSplat（2025.03）\cite{2503_18052}}
SceneSplat 把 3DGS 场景理解从“每个场景单独蒸馏特征”推进到“大规模 3DGS 预训练”。前面的 LangSplat、OpenGaussian 或 Feature 3DGS 通常针对一个重建场景优化语义特征；换一个场景，就要重新提取 2D 特征、重新提升或重新训练。SceneSplat 的改变是把高斯基元本身当作模型输入，训练一个能从 3DGS 参数直接预测逐基元语义特征的通用编码器。

它首先构建 SceneSplat-7K：从 ScanNet、Matterport3D、ScanNet++ 等七个室内数据源生成 7916 个 3DGS 场景。为了得到监督标签，方法用 SAMv2 产生物体级 2D 掩码，用 SigLIP2 提取视觉语言特征，再借 Occam's LGS 这类语言高斯提升方法把 2D 特征转成每个高斯的 3D 语义标签。这样得到的是成规模的“高斯参数-语义特征”配对数据，而不是单场景优化结果。

预训练模型使用 Transformer 编码器-解码器处理高斯基元参数和局部上下文，输入包括位置、颜色、不透明度、尺度、旋转等 3DGS 属性，输出每个 primitive 的高维语言特征。监督损失由三部分组成：余弦相似度让预测特征方向对齐 VLM 标签；L2 损失约束欧氏空间误差；聚合对比损失把同类高斯的 pooled feature 拉近、不同语义类推远。除此之外，SceneSplat 还设计 GaussSSL 自监督预训练，用 masked Gaussian modeling、自蒸馏和可选的 language-Gaussian alignment 正则潜在空间。相对逐场景特征蒸馏，SceneSplat 的意义在于：它把 3DGS 表示本身变成可被基础模型处理的数据形态，为可泛化的室内开放词汇理解提供了基准和预训练路线。
\textbf{设计动机与边界：} 相对原始 3DGS 以独立高斯和图像重建误差为主的设定，SceneSplat 把问题落在以下可实现的改动上：SceneSplat方法包含三个核心部分：(1)3DGS语言标签收集：使用SAMv2进行物体级分割，SigLIP2提取特征，Occam's LGS将2D特征提升至3D高斯特征场，生成3DGS-特征配对数据；(2)视觉-语言3DGS预训练：适配Transformer编码器-解码器骨干网络，从高斯参数预测高维逐基元特征，使用余弦相似度损失、L2损失和聚合对比损失进行监督；(3)自监督预训练：通过重建和自蒸馏对齐的多目标框架正则化3DGS参数编码期间的潜在空间。SceneSplat-7K数据集来源于ScanNet、Matterport3D等七个数据集，计算资源相当于150 GPU天。其设计目的不是单独增加一个模块，而是用上述表示、监督或渲染约束消除原流程在目标任务中的具体失配。提出首个原生基于3DGS的大规模室内场景理解方法，构建了包含7916个场景的SceneSplat-7K数据集，并设计自监督学习方案从无标注场景中学习丰富的3D特征。适用时还应保留论文给出的边界：训练需要大量GPU资源（150 GPU天用于数据集生成）；目前仅限于室内场景；模型对输入高斯基元数量敏感。
\textbf{方法信息：} 代码：\url{https://github.com/unique1i/scenesplat}；GitHub 关注度较高，约 350 stars。



% METHOD: 2508.02172
\paragraph{GaussianCross（2025.08）\cite{2508_02172}}
GaussianCross把3DGS从单场景语义蒸馏工具扩展为跨模态自监督预训练表示。点云预训练常见问题是场景尺度不一致、点间判别难度不足，以及只看几何会缺少外观和语义线索。GaussianCross先把原始点云归一化到统一长方体空间，并把点转成带尺度、位置编码和体素锚点坐标的高斯表示，使不同场景在同一规范空间中接受预训练。

训练时，高斯不是为了照片级渲染，而是作为2D图像和3D点云之间的可微中间层。三属性自适应蒸馏同时使用RGB外观、深度/法线几何和2D基础模型产生的语义伪标签；每一路都有自适应权重，避免某一类噪声监督主导训练。模型通过splatting把3D高斯特征投到2D视图，与图像特征对齐，再把跨模态信号反过来塑造3D骨干。相对Feature 3DGS服务单个场景的可渲染特征，GaussianCross把“高斯可投影性”用在大规模无标注3D表征学习上。
\textbf{实现细节：} GaussianCross：跨模态自监督。GaussianCross（2025.08） 从跨模态输入出发，提出跨模态自监督三维表示学习架构，将尺度不一致的三维点云转换为统一的立方体归一化高斯表示，通过三属性自适应蒸馏 splatting 模块构建三维特征场，协同捕获外观、几何和语义线索。其立方体归一化将点云 $(x, y, z)$ 映射到 $[-1, 1]^{}3$，消除不同传感器尺度差异；三属性蒸馏 splatting 分别处理颜色、法线和语义三种属性。
\textbf{方法信息：} 代码：\url{https://github.com/RayYoh/GaussianCross}。
\textbf{设计动机与边界：} 相对原始 3DGS 以独立高斯和图像重建误差为主的设定，GaussianCross 把问题落在以下可实现的改动上：GaussianCross架构包含两个核心设计：(1)归一化长方体高斯初始化：将原始尺度不一致的3D点云转换为统一的归一化高斯表示。对场景点云进行长方体归一化，应用傅里叶位置编码，附加体素锚点坐标作为结构化先验。这种归一化策略解决了场景间尺度不一致问题，实现可泛化的预训练；(2)三属性自适应蒸馏泼溅：构建3D特征场，通过自适应蒸馏策略协同捕获外观（RGB颜色）、几何（法线/深度）和语义（分割伪标签）三类线索。设计三路自适应权重平衡不同属性的蒸馏贡献，维持跨模态一致性。预训练信号来自2D RGB图像与3D点云的跨模态对应关系，无需人工标注。整个框架通过前馈3DGS渲染进行端到端训练。其设计目的不是单独增加一个模块，而是用上述表示、监督或渲染约束消除原流程在目标任务中的具体失配。提出GaussianCross，整合前馈3DGS的跨模态自监督3D表征学习架构。将尺度不一致的点云转换为统一的归一化长方体高斯表示，设计三属性自适应蒸馏泼溅模块构建3D特征场。在ScanNet200上全微调分别提升9.3% mIoU和6.1% AP50，参数和数据效率显著。适用时还应保留论文给出的边界：依赖RGB-D数据的多模态输入；归一化可能损失场景的绝对尺度信息；大规模预训练的计算成本较高。
% METHOD: 2510.23521
\paragraph{Online 3DGS Explicit Memory for Video Segmentation（2025.10）\cite{2510_23521}}
把3DGS用作类别无关视频分割的显式记忆。首帧中像素被反投影成3D高斯并赋予分割ID特征；后续帧对未覆盖区域或深度变化区域新增高斯，并用FastSAM或SAM2输出更新语义特征。当前帧分割时，系统渲染3DGS记忆中的分割ID特征图，与基础模型预测融合，修正短时遮挡和视角变化造成的漂移。相对纯2D视频分割，它把跨帧记忆从隐式隐藏状态改成可渲染3D高斯地图。
实现上，这套记忆为每个高斯额外保存分割ID特征，而不是只保存颜色和不透明度。首帧像素根据深度反投影并初始化高斯，同时写入对应对象的ID特征；后续帧中，系统对新出现区域或深度差异大的区域创建新高斯，并用FastSAM或SAM2的输出更新已有语义特征。当前帧预测时，3DGS先把记忆中的ID特征渲染回图像平面，再与基础分割模型的2D预测融合。这样，遮挡前见过的对象可以通过3D记忆重新投影回来，减少纯2D视频分割的漂移。
\textbf{方法信息：} Amazon 参与了这篇论文。
\textbf{设计动机与边界：} 相对原始 3DGS 以独立高斯和图像重建误差为主的设定，Explicit Memory through Online 3D Gaussian Splatting Improves Class-Agnostic Video Segment 把问题落在以下可实现的改动上：方法核心是在线构建3DGS语义地图作为显式3D记忆：(1) 3DGS表示扩展——每个高斯基元额外关联分割ID特征向量，编码连续特征空间中的身份信息；(2) 在线构建——首帧将像素反投影为3D点初始化高斯，分配分割ID特征；后续帧对未覆盖区域或深度差异大区域创建新高斯，使用分割模型输出更新语义特征；(3) 记忆融合——FastSAM-Splat和SAM2-Splat分别将3DGS渲染的分割ID特征图与各自基础模型的预测融合，利用3D记忆纠正和优化当前帧分割结果。关键创新在于将3DGS从纯几何/外观表示扩展到包含语义分割信息的统一表示。其设计目的不是单独增加一个模块，而是用上述表示、监督或渲染约束消除原流程在目标任务中的具体失配。提出基于在线3DGS的显式3D记忆增强视频分割方法，开发FastSAM-Splat和SAM2-Splat，通过增量构建3DGS语义地图并融合到基础模型预测中，显著提升分割准确性和时序一致性。适用时还应保留论文给出的边界：依赖深度图和相机位姿输入；对动态场景的处理能力有限；3DGS增量构建存在累积误差。
% METHOD: 2511.16301
\paragraph{Upsample Anything（2025.11）\cite{2511_16301}}
虽然不是3DGS专用方法，但对本章很多2D基础特征提升有直接价值。它在测试时为每个高分辨率像素学习各向异性Gaussian核，通过重建下采样图像获得空间和光度混合权重；随后把同一核迁移到低分辨率基础模型特征，把特征splat到高分辨率。相对直接双线性上采样，它不生成新语义内容，而是学习局部混合方式，因此能让SAM、CLIP或DINO特征更细地对齐边界，再被蒸馏到3DGS。
Upsample Anything把上采样写成逐像素Gaussian核的测试时优化。给定高分辨率图像，它先为每个像素学习各向异性核参数，包括两个空间尺度、方向和颜色范围权重，使高分辨率图像经这些局部核混合后能重建下采样结果。学到的核并不依赖某个语义类别，而是描述“这个像素附近哪些邻居应该一起混合”。随后同一组核被迁移到低分辨率基础模型特征上，把特征按各向异性权重splat回高分辨率。对3DGS语义蒸馏而言，这类边界对齐特征可减少把模糊2D特征提升到3D时的物体串色。
\textbf{方法信息：} Microsoft 参与了这篇论文。
\textbf{设计动机与边界：} 相对原始 3DGS 以独立高斯和图像重建误差为主的设定，Upsample Anything 把问题落在以下可实现的改动上：Upsample Anything包含两个阶段：(1)测试时优化(TTO)：学习逐像素的各向异性Gaussian参数(σx, σy, θ, σr)，通过从高分辨率图像重建其下采样版本来优化。每个像素学习空间和光度相似邻居应如何混合，发现可超越图像域泛化的局部混合权重；(2)特征渲染：优化后的Gaussian核直接迁移到基础特征空间，低分辨率特征通过学习的各向异性加权机制splat到高分辨率。空间权重由各向异性Gaussian协方差决定，范围权重由图像引导。设计灵感来自联合双边上采样(JBU)，但通过逐像素各向异性核提升表达力。纯混合操作（无内容生成），保证跨骨干网络和任务的迁移能力。其设计目的不是单独增加一个模块，而是用上述表示、监督或渲染约束消除原流程在目标任务中的具体失配。提出Upsample Anything，轻量级测试时优化(TTO)框架，无需训练即可将低分辨率特征恢复为高分辨率。通过学习各向异性Gaussian核结合空间和范围线索，桥接Gaussian Splatting和联合双边上采样。学习到的核作为通用边缘感知算子，在架构和模态间无缝迁移，每张224x224图像仅需0.419秒，在语义分割、深度估计等任务上达到SOTA。适用时还应保留论文给出的边界：依赖高分辨率引导图像；测试时优化增加推理延迟；对极端下采样倍率(>16x)效果可能下降。
% METHOD: 2512.05060
\paragraph{4DLangVGGT（2025.12）\cite{2512_05060}}
4DLangVGGT构建的是可泛化的4D语言场，而不是每个动态场景单独优化一个语言3DGS。4-LEGS这类方法从动态高斯视频中逐场景蒸馏ViCLIP特征，能查询时间段和区域，但扩展到大量真实视频会受训练成本限制。4DLangVGGT冻结StreamVGGT几何编码器，利用交替空间注意力和因果时间注意力得到同时包含3D几何和时间动态的token，把“动态几何理解”作为预训练底座。

Semantic Bridging Decoder再把这些几何token转成语言对齐表示：DPT结构先细化时空上下文，随后双头分别输出语言语义和RGB重建。语言头保证token能和开放词汇语义对齐，RGB头则把语义学习约束在真实外观和几何结构上，避免纯语义监督漂离场景。训练时联合语义对齐和外观重建损失。相对4-LEGS的逐场景语言特征优化，4DLangVGGT的改变是前馈化：输入动态视频后直接生成语言感知4D场，更接近可部署的动态场景基础模型。
\textbf{方法信息：} 代码：\url{https://github.com/hustvl/4dlangvggt}。
\textbf{设计动机与边界：} 相对原始 3DGS 以独立高斯和图像重建误差为主的设定，4DLangVGGT 把问题落在以下可实现的改动上：4DLangVGGT框架包含三个核心组件：（1）StreamVGGT几何编码器：冻结的StreamVGGT交替使用空间注意力和因果时间注意力，生成编码精细3D几何结构和时间动态的几何Token。冻结策略基于两方面考虑——StreamVGGT已在大规模视频数据上预训练几何重建，提供强泛化的时空表示；冻结避免冗余优化、降低计算成本，让训练专注于语义对齐。（2）Semantic Bridging Decoder (SBD)：首先通过Dense Prediction Transformer (DPT)将几何Token精炼为上下文丰富的表示，然后双头解码——g\_Lang映射到D维语义空间（UNet结构），g\_RGB重建RGB帧确保感知一致性。双头设计实现互引导：语义嵌入保持结构保真度，RGB重建被语义监督正则化。（3）多目标训练策略：联合优化语义对齐和外观重建损失。其设计目的不是单独增加一个模块，而是用上述表示、监督或渲染约束消除原流程在目标任务中的具体失配。提出首个无需逐场景优化的前馈4D语言接地框架，通过冻结StreamVGGT编码器提供时空几何表示，SBD将几何Token映射到语言对齐空间，在HyperNeRF和Neu3D上达到SOTA，单场景训练提升2%，多场景训练提升1%。适用时还应保留论文给出的边界：语义精度可能略低于逐场景优化方法；对训练数据中未见过的动态模式可能泛化有限；推理时需要完整的StreamVGGT前向传播。
% METHOD: 2512.20927
\paragraph{Quantile Rendering（2025.12）\cite{2512_20927}}
解决高维特征渲染代价。普通alpha混合需要沿光线处理大量高斯，并为每个高斯混合高维通道；Quantile Rendering把透射率分布均匀分成若干段，只选代表性分位高斯执行特征混合，将复杂度从与所有命中高斯和通道数相关，降到与少量分位样本和通道数相关。GS-Net把该渲染器接入Point Transformer式3D网络，使2D语义损失能高效反传到3D高斯特征预测。相对LangSplatV2的稀疏系数，Q-Render从光线采样端减少高维混合量。
\textbf{实现细节：} Quantile Rendering：分位数渲染加速。Quantile Rendering（2025.12） 提出分位数渲染策略，仅稀疏采样沿每条射线有主导影响的高斯而非密集采样所有交叉高斯，集成到泛化三维神经网络 GS-Net 中实现 512 维特征图的约 43.7 倍渲染加速。
\textbf{实现补充：} 实现上还要核对输入、输出、参数、特征、初始化、训练损失、正则、渲染或推理接口，以及优化中的剪枝、分裂、采样、编码或解码步骤；这些抓手决定该方法到底改了 3DGS 管线的哪一环。
\textbf{方法信息：} NVIDIA 参与了这篇论文。
\textbf{设计动机与边界：} 相对原始 3DGS 以独立高斯和图像重建误差为主的设定，Quantile Rendering 把问题落在以下可实现的改动上：方法包含两个核心组件：(1)Quantile Rendering (Q-Render)，不同于传统体渲染密集采样所有与光线相交的高斯，Q-Render根据超参数K将透射率T∈[0,1]均匀分为K+1段，选择K个分位数高斯——即沿光线主导透射率分布的代表性高斯。仅对这些高斯执行alpha混合和特征渲染，将计算复杂度从O(N×C)降低到O(K×C)，其中K<<N。包含三个子步骤：分位数高斯采样、alpha混合和特征归一化。(2)GS-Net，将Q-Render集成到基于Point Transformer v3的3D神经网络中，通过体素化处理3D高斯，预测高斯特征而非逐场景优化。Q-Render作为2D监督与3D神经网络之间的高效桥梁，允许梯度从图像空间损失反向传播到3D神经网络预测。其设计目的不是单独增加一个模块，而是用上述表示、监督或渲染约束消除原流程在目标任务中的具体失配。提出Quantile Rendering(Q-Render)，一种基于分位数的稀疏采样渲染策略，仅选择沿光线具有主导影响力的少量高斯进行渲染，在512维CLIP特征图上实现约43.7倍加速。适用时还应保留论文给出的边界：分位数数量K需要手动设置；对高斯分布质量有一定要求；仅在室内场景验证。
% METHOD: 2606.28840
\paragraph{DLGStream（2026.06）\cite{2606_28840}}
DLGStream把语言嵌入放进自由视点视频流媒体约束下考虑。离线语言高斯可以为每个场景存大量CLIP特征，但FVV系统要求低帧大小、高FPS和按GOP连续传输；如果直接把高维语言特征附到每帧动态高斯上，带宽和解码开销会迅速失控。DLGStream先按GOP组织多视角动态视频，并参考Swift4D做静态-动态高斯分解：静态背景不用每帧重复变形，动态部分才随时间更新。

语言侧则把CLIP等开放词汇特征压缩成可随Gaussian attributes一起流式传输的动态语言表示。查询时，客户端既能按自由视点渲染视频，也能用文本在4D时空中选中对象或区域。论文在N3DV等动态数据上报告，相对4DLangSplat，DLGStream提升开放词汇查询指标，同时把帧大小降到约十分之一、FPS提高约5倍。相对4-LEGS关注事件定位，DLGStream的增量是把语言高斯放进真实FVV系统的码率、帧率和动态更新预算里。
\textbf{实现细节：} DLGStream（2026.06） 延续 Language-embedded 3DGS 在"Language-embedded dynamic Gaussian video"上的基本设定，主要改变是：面向流式 FVV 的低帧尺寸高 FPS 语言 Gaussian 表示。3DGS 的几何-渲染统一表示保持为底层地图。其设计动机是：针对现有 language Gaussian 表示难以满足 FVV 低帧尺寸和高帧率要求的问题，提出可流式动态语言嵌入表示。具体而言，在 3DGS/4DGS 自由视点视频表示中蒸馏或压缩 CLIP 等语言特征，使动态场景能够进行开放词表查询，同时满足 streaming 约束。适用边界是：语言特征压缩可能牺牲细粒度语义；动态场景中的时序一致查询和边界精度仍需关注。
\textbf{实现补充：} 实现上还要核对输入、输出、参数、特征、初始化、训练损失、正则、渲染或推理接口，以及优化中的剪枝、分裂、采样、编码或解码步骤；这些抓手决定该方法到底改了 3DGS 管线的哪一环。
\textbf{技术实现补充：} DLGStream。DLGStream（2026.06） 延续 Language-embedded 3DGS 在"Language-embedded dynamic Gaussian video"上的基本设定，主要改变是：面向流式 FVV 的低帧尺寸高 FPS 语言 Gaussian 表示。3DGS 的几何-渲染统一表示保持为底层地图。其设计动机是：针对现有 language Gaussian 表示难以满足 FVV 低帧尺寸和高帧率要求的问题，提出可流式动态语言嵌入表示。具体而言，在 3DGS/4DGS 自由视点视频表示中蒸馏或压缩 CLIP 等语言特征，使动态场景能够进行开放词表查询，同时满足 streaming 约束。适用边界是：语言特征压缩可能牺牲细粒度语义；动态场景中的时序一致查询和边界精度仍需关注。
\textbf{方法信息：} 代码：\url{https://github.com/kkkzh/dlgstream}。
\textbf{设计动机与边界：} 相对原始 3DGS 以独立高斯和图像重建误差为主的设定，DLGStream 把问题落在以下可实现的改动上：在 3DGS/4DGS 自由视点视频表示中蒸馏或压缩 CLIP 等语言特征，使动态场景能够进行开放词表查询，同时满足 streaming 约束。其设计目的不是单独增加一个模块，而是用上述表示、监督或渲染约束消除原流程在目标任务中的具体失配。针对现有 language Gaussian 表示难以满足 FVV 低帧尺寸和高帧率要求的问题，提出可流式动态语言嵌入表示。适用时还应保留论文给出的边界：语言特征压缩可能牺牲细粒度语义；动态场景中的时序一致查询和边界精度仍需关注。

\section{6D位姿估计}

6D 位姿方向把 3DGS 从场景表示变成几何匹配工具。高斯模型提供可渲染、带视角相关颜色和连续表面线索的目标几何，位姿估计则反过来寻找相机或物体变换，使渲染线索与观测图像对齐。与传统 PnP、CAD 匹配或纯学习回归相比，这条路线的关键是把高斯的可渲染性、稠密外观和显式几何同时用于姿态求解、合成标注或跟踪。
% METHOD: 2407.15484
\paragraph{6DGS（2024.07）\cite{2407_15484}}
6DGS把相机位姿估计写成“反过来读3DGS”的几何问题。常见做法是先找2D-3D匹配再跑PnP，或者给定初值后反复渲染优化；这些路线要么依赖稳定特征匹配，要么对初始位姿敏感。6DGS从已训练好的3DGS出发，在每个高斯椭球表面构造Radiant Ellicell：把椭球按近似等面积切成小单元，每个单元从椭球中心向外发出一条带颜色和位置属性的光线。目标图像的像素特征与这些光线特征通过注意力建立绑定，等价于判断“这个像素最可能看到了哪条3D高斯光线”。

一旦获得一组像素-光线绑定，位姿就不再需要迭代搜索。多条被选中光线在空间中应当共同指向未知相机中心，6DGS用加权最小二乘求这些光线的最佳交点作为平移；再根据相机中心与光线方向的对应关系恢复旋转。这里的关键增量是把3DGS中连续椭球的表面显式离散成可匹配的光线集合，使“渲染图像像素”变成“像素反查光线”的闭式问题。相对渲染优化式定位，它不需要初始姿态，速度可达到近实时，适合作为3DGS地图上的重定位锚点。
\textbf{实现细节：} 单图像与闭合形式位姿估计。6DGS（2024.07） 从完全不同的角度解决位姿估计——不是优化对象位姿，而是从单个 RGB 图像和 3DGS 场景模型估计相机位姿。核心创新是定义 Ellicell（Ellipsoid Cell）：将每个 3D 高斯椭球体的表面等面积划分为若干单元，在每个单元中心沿从椭球中心到单元中心的方向生成射线，构建射线集 $V = \{\langle v_o, v_d, v_c \rangle\}_{i=1}^N$，其中 $v_o \in \mathbb{R}^{}3$ 为椭球中心（射线原点），$v_d \in \mathbb{R}^{}3$ 为从中心到单元中心的方向，$v_c \in \mathbb{R}^{}3$ 为通过 3DGS 渲染函数 $\mathcal{R}$ 合成的射线颜色。射线-像素绑定通过注意力图 $A$ 实现：计算目标图像像素特征 $\phi(I_t)$ 和射线特征 $\phi(V)$ 之间的注意力权重，注意力值越高表示该射线-像素对越可能对应有效交点。最终通过加权最小二乘求最优射线束的交集直接计算相机中心，旋转则由射线方向与像素方向的对应关系确定。SplatPose（2025.03） 提出 Dual-Attention Ray Scoring Network（DARS-Net），通过几何域注意力机制解耦位置对齐和角度对齐，显式建模方向依赖性以缓解旋转歧义，配合粗到细优化管线通过密集 2D 特征对齐精化位姿估计。
\textbf{实现补充：} 实现上还要核对输入、输出、参数、特征、初始化、训练损失、正则、渲染或推理接口，以及优化中的剪枝、分裂、采样、编码或解码步骤；这些抓手决定该方法到底改了 3DGS 管线的哪一环。
\textbf{设计动机与边界：} 相对原始 3DGS 以独立高斯和图像重建误差为主的设定，6DGS 把问题落在以下可实现的改动上：6DGS流程：(1) 对3DGS模型中每个椭球体，定义Radiant Ellicell将椭球体表面均匀分割为等面积单元格（使用Ramanujan近似计算椭球表面积，沿主轴切片为ribbon再细分为cell）；(2) 每个cell生成一条从椭球中心经过cell中心的光线，光线携带渲染颜色信息；(3) 提取目标图像的像素特征和光线颜色特征；(4) 通过注意力图计算像素-光线绑定分数，选择最佳绑定的光线束；(5) 加权最小二乘法求光线交点确定相机中心位置；(6) 根据相机中心和光线方向确定相机旋转。整个过程是确定性的闭式解，不需要迭代优化或初始姿态估计。由于基于3DGS的快速渲染能力，实现近实时15fps。其设计目的不是单独增加一个模块，而是用上述表示、监督或渲染约束消除原流程在目标任务中的具体失配。提出6DGS，首次通过反转3DGS渲染过程实现无需初始化的6DoF姿态估计。闭式解无需迭代分析合成循环，在真实场景上旋转精度提升12%、平移精度提升22%，达到近实时15fps。适用时还应保留论文给出的边界：依赖3DGS模型的质量和完整性；对严重遮挡或大基线变化的场景可能表现下降。
\textbf{方法信息：} 代码：\url{https://github.com/mbortolon97/6dgs}；GitHub 关注度较高，约 180 stars。



% METHOD: 2503.11692
\paragraph{FloPE（2025.03）\cite{2503_11692}}
FloPE把3DGS用作农业机器人姿态学习的数据工厂，而不是直接把3DGS作为在线定位器。花卉授粉需要花心位置和朝向，真实场景标注成本高，机器人端模型又不能太重。FloPE先用多视角图像重建植物的3DGS，再在高斯点云或渲染视图中快速标注花朵6D姿态；由于相机位姿和3DGS几何已知，同一个3D姿态标注可以自动投到大量新视角中，生成带精确姿态标签的照片级合成训练图。

训练流程是教师-学生式的。大教师模型负责在合成和真实图像上学习花卉检测与姿态估计，姿态可用位置、朝向和9D旋转表示约束；随后把教师预测蒸馏到轻量学生模型，例如面向机器人部署的检测与PoseNet组合。系统还可用RGB-D相机或近距离深度信息对末端位姿做二次修正。相对6DGS从已建高斯场反推出相机位姿，FloPE的增量是利用3DGS批量制造“有真实几何标签的训练视角”，把困难的真实花姿态标注转化为可控的合成数据生成问题。
\textbf{实现细节：} 特殊应用。光场位姿跟踪 利用光场图像提取语义和几何特征转化为视角依赖高斯 splats 作为统一对象表示，支持可微渲染和位姿优化，在包含反射等复杂视觉行为的挑战性物体上与基于模型的跟踪器竞争。LineGS 将 3DGS 扩展到 3D 线段表示，利用边缘处高密度的 Gaussian 点分布精化和优化几何方法生成的初始线段，在几何精度和模型紧凑性上显著提升。
\textbf{方法信息：} 代码：\url{https://github.com/wvu-irl/flope}。
\textbf{设计动机与边界：} 相对原始 3DGS 以独立高斯和图像重建误差为主的设定，FloPE 把问题落在以下可实现的改动上：FloPE分为训练和实施两个流程。训练流程：(1) 使用3DGS重建植物模型，利用其高保真照片级渲染能力；(2) 通过COLMAP获取相机姿态和3D点用于初始化高斯中心；(3) 开发3D姿态标注工具，在3DGS点云中标注花卉姿态（20-25朵花几分钟完成）；(4) 从多视角渲染带标注的合成图像作为训练数据；(5) 知识蒸馏：从大教师模型（花卉检测+姿态估计）到轻量学生模型。实施流程：训练好的花卉姿态估计器部署到机器人上引导授粉操作。其设计目的不是单独增加一个模块，而是用上述表示、监督或渲染约束消除原流程在目标任务中的具体失配。首个利用3DGS生成合成数据用于农业机器人花卉姿态估计的框架，通过知识蒸馏实现轻量级模型（0.6cm位置误差、19.14度角度误差），授粉成功率达78.75%。适用时还应保留论文给出的边界：3DGS重建质量受拍摄条件影响；标注依赖人工判断引入误差；仅限于特定花卉类型。
% METHOD: 2508.15972
\paragraph{UnPose（2025.08）\cite{2508_15972}}
处理零样本无模型物体位姿估计。它从单帧RGB-D图像出发，用Wonder3D生成初始多视角图像，并用不确定性估计判断哪些像素可靠；VGGT对齐相对位姿后，通过位姿图优化重建初始3DGS物体场。后续位姿估计由3DGS渲染和查询图像特征匹配完成，新观测再按不确定性引导融合进地图。相对FloPE需要目标域训练数据，UnPose用扩散先验和不确定性补全较早见到的物体。
UnPose的初始化先把单帧RGB-D目标扩展成一个可优化的小型物体场。Wonder3D生成若干先验视图，LLLA估计这些生成视图的像素不确定性，VGGT提供相对位姿；位姿图优化把真实观测和扩散先验对齐后，得到初始3DGS目标模型。后续查询图像进入系统时，3DGS渲染出的视图与查询图像做特征匹配，给出物体位姿；新观测再按不确定性写回地图，高不确定区域优先被真实图像修正。这个流程的重点是把扩散补全当作可被不确定性约束的初始假设，而不是直接相信生成视图。
\textbf{方法信息：} Huawei 参与了这篇论文。
\textbf{设计动机与边界：} 相对原始 3DGS 以独立高斯和图像重建误差为主的设定，UnPose 把问题落在以下可实现的改动上：方法包含四个模块：(1)初始化：从单帧RGB-D图像出发，使用Wonder3D多视图扩散模型生成初始6个视角图像，采用最后层拉普拉斯近似(LLLA)估计像素级认知不确定性，通过VGGT对齐相对位姿并经位姿图优化重建初始3DGS物体场；(2)位姿估计：利用3DGS模型渲染与查询图像的特征匹配估计初始位姿；(3)后端优化：将扩散先验视图和后续观测整合到位姿图中联合优化；(4)3DGS建图：不确定性引导融合新视图，在不确定区域优先更新，确保全局一致性。其设计目的不是单独增加一个模块，而是用上述表示、监督或渲染约束消除原流程在目标任务中的具体失配。提出UnPose框架，首次实现零样本无CAD模型的6D物体位姿估计。利用预训练扩散模型的像素级认知不确定性引导多视图融合，通过3DGS表示实现高效的增量式3D重建。核心创新是将扩散先验的不确定性估计与位姿图优化结合，在不确定区域优先使用新观测数据，在可靠区域保留先验知识，持续提升位姿估计精度和重建质量。适用时还应保留论文给出的边界：依赖预训练扩散模型的质量；初始化阶段的不确定性估计可能不够精确；对透明或反光物体的处理有待改进。
% METHOD: 2512.13007
\paragraph{Light Field Based 6DoF Tracking（2025.12）\cite{2512_13007}}
Light Field Based 6DoF Tracking用光场相机代替预建CAD或参考视图库。对于未见物体，传统6D跟踪通常需要已知模型、类别先验或前期扫描；这在反光、透明或纹理复杂物体上很难。该方法先用Grounding DINO和SAM2从光场子视图中分割目标，再融合光场视差和Depth Anything深度，把多子视图观测反投影成目标点云。

接着，系统为目标拟合带视角相关外观的高斯泼溅，球谐系数来自多视点颜色观测。连续帧跟踪时，先用彩色ICP做粗配准，再通过光度优化细化两帧高斯泼溅之间的相对变换，得到物体6DoF运动。相对6DGS从场景高斯反推相机位姿，这篇工作把高斯作为“由光场即时构建的物体模型”，用于跨帧物体跟踪；关键优势是无需预建模型，关键依赖则是光场视差和分割质量。
\textbf{实现细节：} 光场位姿跟踪：多方向光线表示。光场位姿跟踪（2025.12） 利用光场图像的多视角特性提取语义和几何特征，转化为视角依赖高斯 splats 作为统一对象表示。与标准 3DGS 位姿方法依赖 RGB 或 RGB-D 输入不同，光场输入天然编码了多方向的光线信息，使高斯 splats 能够捕获反射、折射等复杂视觉行为，在包含此类挑战性外观的物体上与基于精确 CAD 模型的跟踪器竞争。
\textbf{技术实现补充：} 光场位姿跟踪：多方向光线表示。光场位姿跟踪（2025.12） 利用光场图像的多视角特性提取语义和几何特征，转化为视角依赖高斯 splats 作为统一对象表示。与标准 3DGS 位姿方法依赖 RGB 或 RGB-D 输入不同，光场输入天然编码了多方向的光线信息，使高斯 splats 能够捕获反射、折射等复杂视觉行为，在包含此类挑战性外观的物体上与基于精确 CAD 模型的跟踪器竞争。
\textbf{方法信息：} 代码：\url{https://github.com/nagonch/lift-6dof}。
\textbf{设计动机与边界：} 相对原始 3DGS 以独立高斯和图像重建误差为主的设定，Light Field Based 6DoF Tracking of Previously Unobserved Objects 把问题落在以下可实现的改动上：方法流程：（1）光场处理：使用Grounding DINO+SAM2分割物体，融合光场视差和Depth Anything的深度估计，反投影获得3D点云。（2）高斯泼溅初始化：通过拟合球谐系数到多视点颜色观测创建视点相关高斯泼溅。使用彩色ICP粗配准，光度优化精配准。（3）位姿优化：优化最佳对齐连续帧高斯泼溅的变换矩阵估计物体相对运动。关键洞察是光场图像天然充当了先前方法需要预建的物体模型。其设计目的不是单独增加一个模块，而是用上述表示、监督或渲染约束消除原流程在目标任务中的具体失配。首次提出基于光场图像的物体跟踪方法，将视觉基础模型的语义几何特征转换为视点相关高斯泼溅作为在线跟踪参考模型，无需预录制参考视图，在反射物体上达到SOTA模型跟踪器的水平。适用时还应保留论文给出的边界：需要光场相机硬件；计算开销较大；对快速运动的处理有限

\section{3D目标检测}

目标检测方法通常不把3DGS作为最终渲染资产，而是把高斯当作更密集、更结构化的3D特征载体。相对点云或雷达点，Gaussian具有位置、尺度、旋转、不透明度和可投影特征，适合在BEV或检测网络中聚合。
% METHOD: 2410.01647
\paragraph{3DGS-DET（2024.10）\cite{2410_01647}}
较早把3DGS显式用于室内3D目标检测。它先用2D边界检测结果引导3DGS训练，让高斯分布在目标边界处更清楚；随后把2D框投影成3D视锥体，构造Box-Focused采样概率，让目标内部高斯更容易被采样，背景高斯被减少。采样后的高斯属性作为检测器输入。相对直接从点云检测，它用可渲染高斯把多视角图像信息浓缩成带边界引导的3D原语。
3DGS-DET的检测输入来自经过目标边界强化的高斯集合。训练3DGS时，它先把2D边界检测结果叠加到图像监督中，让高斯在前景目标和背景交界处形成更清楚的空间分布。随后，2D检测框被投影成3D视锥体，并转化为Box-Focused采样概率；位于视锥内、可能属于目标的高斯更容易被保留，背景高斯被下采样。检测网络读入的不是渲染图，而是采样后的高斯位置、尺度、旋转、颜色和不透明度等属性。相对点云检测，它多利用了多视图图像重建中已经整合好的边界和外观信息。
\textbf{方法信息：} 代码仓库：\url{https://github.com/yangcaoai/3dgs-det}，GitHub 关注度较高，约 165 stars；核验时完整实现暂未确认。
\textbf{设计动机与边界：} 相对原始 3DGS 以独立高斯和图像重建误差为主的设定，3DGS-DET 把问题落在以下可实现的改动上：3DGS-DET的流水线：(1)使用2D边界引导训练3DGS：在输入图像上执行目标边界检测，将边界叠加到图像上再训练3DGS，使高斯斑点的空间分布更清晰地区分前景目标和背景；(2)Box-Focused采样：将2D框投影到3D空间形成视锥体，3D高斯在视锥内的概率更高，基于该概率分布进行采样以保留更多目标斑点并减少背景噪声；(3)将采样后的3DGS输入检测框架进行训练和推理。方法直接利用原始3DGS的属性（位置、缩放、旋转、颜色、不透明度）作为检测器输入。其设计目的不是单独增加一个模块，而是用上述表示、监督或渲染约束消除原流程在目标任务中的具体失配。首次将3DGS引入室内3D目标检测，提出2D边界引导和Box-Focused采样两大策略，在ScanNet上mAP@0.25提升+6.0、mAP@0.5提升+7.8，在ARKITScenes上mAP@0.25提升+14.9。适用时还应保留论文给出的边界：依赖2D边界检测的质量；采样策略可能遗漏部分目标区域；仅针对室内场景验证
% METHOD: 2507.19856
\paragraph{RaGS（2025.07）\cite{2507_19856}}
RaGS进一步把4D雷达和单目图像融合到同一个高斯场中。RadarGaussianDet3D主要从雷达点出发，解决稀疏雷达到BEV特征的问题；RaGS认为单目图像提供前景语义和纹理边界，4D雷达提供深度、速度和恶劣天气鲁棒性，二者应该共同决定高斯的位置和属性。Frustum-based Localization Initiation先把图像前景点素沿视锥反投影，初始化一批粗高斯中心。

随后，Iterative Multimodal Aggregation在这些高斯附近反复融合图像语义、雷达几何和速度线索，让高斯逐步集中到真实目标区域；Multi-level Gaussian Fusion再把不同层级的高斯场渲染为BEV特征，送入3D检测头。相对固定BEV网格，RaGS的高斯可以把表示容量动态分配给前景和运动目标；相对只用图像proposal，它又利用雷达速度帮助区分动态目标和背景。
\textbf{方法信息：} 代码：\url{https://github.com/shawnnnkb/rags}。
\textbf{设计动机与边界：} 相对原始 3DGS 以独立高斯和图像重建误差为主的设定，RaGS 把问题落在以下可实现的改动上：RaGS采用级联管线构建和渐进优化高斯场：(1)基于视锥的定位初始化(FLI)，将前景点素反投影初始化粗糙高斯中心；(2)迭代多模态聚合(IMA)，显式利用图像语义并隐式整合4D雷达速度几何来精细化感兴趣区域内的高斯；(3)多级高斯融合(MGF)，将高斯场渲染为层次化BEV特征进行3D检测。4D雷达速度维度提供运动线索帮助锚定和精细化高斯空间分布，动态聚焦稀疏和信息丰富区域。其设计目的不是单独增加一个模块，而是用上述表示、监督或渲染约束消除原流程在目标任务中的具体失配。提出首个利用3DGS融合4D雷达和单目线索的3D目标检测框架RaGS，采用级联管线（FLI+IMA+MGF）渐进优化高斯场，利用雷达速度维度锚定和精细化高斯空间分布。适用时还应保留论文给出的边界：依赖4D雷达的精度和可靠性；复杂城区密集场景的检测性能待提升；高斯场的初始化质量对结果影响较大。
% METHOD: 2509.16119
\paragraph{RadarGaussianDet3D（2025.09）\cite{2509_16119}}
RadarGaussianDet3D把4D汽车雷达点云转成高斯BEV特征。雷达点稀疏但有速度信息，传统pillar编码让每个点只贡献到单个BEV格子，容易得到稀疏、不连续的特征图。Point Gaussian Encoder先用轻量PointNet聚合局部邻域，再用自注意力做全局交互，最后为每个雷达点预测高斯属性；这些高斯被splat到BEV平面，一个点可以按尺度和方向影响邻近区域，从而补足雷达稀疏性。

检测损失也改成高斯形式。Box Gaussian Loss把预测框和真值框都转换为3D高斯分布，用分布距离同时约束位置、尺寸和朝向，而不是把框属性拆开独立回归。实现上，论文还设计了Indexing \& Scattering加速BEV splatting，降低遍历式投影的成本。相对普通雷达检测器，RadarGaussianDet3D的核心增量是把“雷达点到BEV格子”的硬映射改成“雷达点到连续高斯分布”的软投影。
\textbf{实现细节：} 雷达与自动驾驶场景。RadarGaussianDet3D（2025.09） 面向 4D 汽车雷达的 3D 检测，与传统室内方法（如 3DGS-DET 针对多视角 RGB 输入）不同，它处理的是雷达点云这种天然稀疏且噪声更大的输入模态。现有方法依赖柱状编码器，每个点仅贡献单个 BEV 网格导致稀疏特征图。RadarGaussianDet3D 设计了 Point Gaussian Encoder（PGE）来弥补这一缺陷，PGE 包含四个步骤：（1）局部特征聚合（LFA）：对每个雷达点的球形邻域 $\mathcal{N}_i$ 应用轻量 PointNet 聚合位置偏移加特征，丰富单点表征；（2）全局特征聚合（GFA）：通过自注意力机制捕获全局上下文 $f_g = \mathrm{FFN}(\mathrm{Attn}(f_l)) + f_l$；（3）高斯属性预测：拼接原始、局部和全局特征后线性投影得到尺度 $s'$、旋转四元数 $q$ 和视角无关特征 $f^g$，其中不透明度固定为 1、位置偏移固定为 0 以避免早期训练中的自由度过大问题，尺度经 sigmoid 激活并限制上界 $s_\mathrm{max}$；（4）BEV 高斯 splatting：将预测的高斯基元 $G = \{(\mu_i, s_i, q_i, o_i, f^g_i)\}_{i=1}^N$ 通过平行投影 splat 到 BEV 平面，用特征 $f^g$ 替代颜色渲染，生成密集 BEV 特征图 $F \in \mathbb{R}^{C \times H \times W}$。相比传统柱状编码器每个点仅影响单个网格，PGE 的每个高斯可自适应地影响多个相邻网格。同时提出 Box Gaussian Loss（BGL）：将真实和预测包围盒 $b = [x, y, z, l, w, h, \theta]$ 转化为 3D 高斯分布 $\mathcal{N}(\mu, \Sigma)$，其中 $\mu = [x,y,z]$，$\Sigma = R_b S_b S_b^T R_b^T$，缩放矩阵 $S_b = \frac{1}{a}\mathrm{diag}(\frac{l}{2}, \frac{w}{2}, \frac{h}{2})$，$a > 0$ 为缩放超参数，通过反向 KL 散度 $D_{\mathrm{KL}}(\hat{\mathcal{N}}, \mathcal{N})$ 统一优化包围盒的全部属性。
\textbf{方法信息：} 代码：\url{https://github.com/xiongweiyi/radargaussiandet3d}。
\textbf{设计动机与边界：} 相对原始 3DGS 以独立高斯和图像重建误差为主的设定，RadarGaussianDet3D 把问题落在以下可实现的改动上：RadarGaussianDet3D包含两个核心创新：(1) Point Gaussian Encoder (PGE)：先通过局部特征聚合（LFA）使用轻量PointNet聚合邻域特征，再通过全局特征聚合（GFA）使用自注意力捕获全局交互，然后预测每个雷达点的高斯属性，最后通过3DGS的BEV splatting将高斯原语投影到BEV平面，产生更密集的特征图。设计了优化的Indexing \& Scattering算法，比遍历法更快，比广播法更省内存。(2) Box Gaussian Loss (BGL)：将真值和预测边界框都转换为3D高斯分布，通过计算分布距离来优化，综合考虑位置、尺寸和朝向等所有属性的关联性。其设计目的不是单独增加一个模块，而是用上述表示、监督或渲染约束消除原流程在目标任务中的具体失配。提出RadarGaussianDet3D，首次将3DGS的splatting技术应用于雷达点云的BEV特征提取。设计了Point Gaussian Encoder (PGE)将雷达点转换为高斯原语后进行BEV rasterization，产生更密集的特征图；同时提出Box Gaussian Loss (BGL)将边界框转换为3D高斯分布进行距离度量，实现更全面一致的优化。适用时还应保留论文给出的边界：依赖雷达点云质量，在极端噪声场景下可能受限；高斯原语参数的预测需要足够的局部特征支撑。

\section{变化检测与场景比较}

变化检测把3DGS用于比较两个时间的场景。关键不是开放词汇，而是把变化前后图像放到同一个3D坐标系中，利用渲染差异、基础模型特征或变化通道定位被移动、消失或新增的物体。
% METHOD: 2411.03706
\paragraph{3DGS-CD（2024.11）\cite{2411_03706}}
3DGS-CD面向“同一场景前后拍了两组图，但相机不对齐、物体可能被移动”的变化检测。它先用变化前图像训练一个标准3DGS，把旧场景变成可从任意新视角渲染的参考；然后利用大部分场景不变这一假设，估计变化后图像的相机位姿。这样，每张变化后图像都能得到一张同视角的“变化前渲染图”，2D变化检测就从跨相机比较变成同视角图像比较。

变化掩码由EfficientSAM特征相似度和Otsu自适应阈值得到，而不是依赖固定类别。随后方法把多视角2D变化区域关联到同一个3D对象，融合成3D变化遮罩，并进一步估计该对象的6D位姿变化。相对Change-3DGS这类在高斯中学习变化属性的路线，3DGS-CD更像一条完整工程流水线：旧场景3DGS提供对齐参考，基础模型提供2D变化线索，多视角几何负责把变化提升为对象级3D结果。
\textbf{方法信息：} 代码：\url{https://github.com/520xyxyzq/3dgs-cd}。
\textbf{设计动机与边界：} 相对原始 3DGS 以独立高斯和图像重建误差为主的设定，3DGS-CD 把问题落在以下可实现的改动上：流水线包含六个模块：(1) 变化前3DGS训练——标准3DGS训练建立场景表示；(2) 变化后相机定位——假设大部分场景不变，使用SfM定位变化后相机；(3) 2D变化检测——在变化后视角渲染变化前图像，使用EfficientSAM图像嵌入计算余弦相似度，Otsu方法动态阈值化得到变化遮罩；(4) 跨视角物体关联——将不同视角的2D变化关联到同一物体；(5) 3D物体分割——融合多视角2D变化得到3D变化遮罩；(6) 物体位姿变化估计——计算6D位姿变化。无需深度输入、用户指令、预定义类别或物体模型。其设计目的不是单独增加一个模块，而是用上述表示、监督或渲染约束消除原流程在目标任务中的具体失配。提出3DGS-CD——首个基于3DGS的3D物体级变化检测方法。创新点包括：(1) 利用3DGS高效渲染能力直接比较变化前后同视角图像；(2) EfficientSAM嵌入进行零样本2D变化检测；(3) 跨视角物体关联和融合获得3D变化遮罩和6D位姿变化。仅需一张变化后图像，18秒完成检测，比NeRF方法快三个数量级。适用时还应保留论文给出的边界：假设大部分场景不变，大范围变化可能影响相机定位；光照变化可能影响2D变化检测精度。
% METHOD: 2412.03911
\paragraph{Change-3DGS（2024.12）\cite{2412_03911}}
用变化前图像训练参考3DGS，再用少量变化后图像学习变化感知表示。系统渲染参考场景，与变化后观测在DINOv2等特征空间中比较，生成候选变化掩码；然后从参考3DGS初始化Change-3DGS，为每个高斯增加独立变化不透明度通道，使被移动或消失的高斯即使RGB不透明度降低，也不会丢失变化信息。相对单纯图像差分，它把变化作为高斯属性优化到3D场景中。
Change-3DGS先把变化前图像训练成参考场景3DGS\_ref，再在变化后图像中渲染同视角参考图。DINOv2等视觉特征比较参考渲染和真实观测，生成候选变化掩码；这些掩码不是最终答案，而是监督一个从3DGS\_ref初始化的变化感知高斯场。每个高斯新增变化通道和独立变化不透明度因子 $\gamma_c$，因此即使某个消失物体的RGB不透明度在新场景优化中降低，它的变化证据仍可保留下来。数据增强用于扩充变化掩码监督，减少少量变化后视角带来的遗漏。相对2D变化检测，Change-3DGS把“哪里变了”显式绑定到3D高斯属性上。
\textbf{实现细节：} Galappaththige et al.（2024.12） 提出了免标签、免位姿的多视角变化检测方法，将变化信息编码为 3DGS 的一个额外通道。该方法首先从参考图像构建 3DGS 参考场景表示，然后在变化后图像视角下渲染参考图像并与之比较，利用 DINOv2 特征得到特征感知变化掩码、通过结构差异得到结构感知变化掩码，两者互补融合为候选变化掩码。接着以参考场景的 3DGS 为初始化，学习一个变化感知 3DGS（Change-3DGS），其中每个高斯增加独立的变化不透明度因子 $\alpha_c$ 和变化幅度 $\beta$——前者允许被移动/移除对象对应的高斯在 RGB 渲染时降低不透明度，同时保留变化掩码中的高变化幅度；后者直接编码变化强度并渲染为变化图。为避免变化高斯被标准不透明度剪枝误删，剪枝条件改为 $\min(\alpha, \alpha_c) < \tau$。仅需 5 张变化后场景的图像，即可学习 3DGS 中的变化通道并生成变化掩码，在复杂多对象场景中 mIoU 和 F1 分别比单视角基线提升 1.7 倍和 1.5 倍，同时支持未见视角的变化掩码生成。
\textbf{技术实现补充：} Galappaththige et al.（2024.12） 提出了免标签、免位姿的多视角变化检测方法，将变化信息编码为 3DGS 的一个额外通道。该方法首先从参考图像构建 3DGS 参考场景表示，然后在变化后图像视角下渲染参考图像并与之比较，利用 DINOv2 特征得到特征感知变化掩码、通过结构差异得到结构感知变化掩码，两者互补融合为候选变化掩码。接着以参考场景的 3DGS 为初始化，学习一个变化感知 3DGS（Change-3DGS），其中每个高斯增加独立的变化不透明度因子 $\alpha_c$ 和变化幅度 $\beta$——前者允许被移动/移除对象对应的高斯在 RGB 渲染时降低不透明度，同时保留变化掩码中的高变化幅度；后者直接编码变化强度并渲染为变化图。为避免变化高斯被标准不透明度剪枝误删，剪枝条件改为 $\min(\alpha, \alpha_c) < \tau$。仅需 5 张变化后场景的图像，即可学习 3DGS 中的变化通道并生成变化掩码，在复杂多对象场景中 mIoU 和 F1 分别比单视角基线提升 1.7 倍和 1.5 倍，同时支持未见视角的变化掩码生成。
\textbf{方法信息：} 代码仓库：\url{https://github.com/gmayday1997/changedet}，GitHub 关注度较高，约 248 stars；核验时完整实现暂未确认。
\textbf{设计动机与边界：} 相对原始 3DGS 以独立高斯和图像重建误差为主的设定，Multi-View Pose-Agnostic Change Localization with Zero Labels 把问题落在以下可实现的改动上：方法流程：(1) 用变化前图像构建参考场景3DGS\_ref。(2) 采集变化后图像，利用DINOv2等视觉模型将渲染参考图与观测推理图进行特征和结构感知比较，生成候选变化掩码。(3) 初始化Change-3DGS（从3DGS\_ref出发），为每个高斯添加变化通道（包含独立变化不透明度因子gamma\_c），在优化推理场景的同时优化变化通道。设计数据增强策略增加训练变化掩码数量。关键设计：独立变化不透明度因子gamma\_c，使得表示消失/移动物体的高斯(其RGB不透明度alpha降低但仍需保持变化信息)不被错误剪枝。其设计目的不是单独增加一个模块，而是用上述表示、监督或渲染约束消除原流程在目标任务中的具体失配。提出首个多视图、无标签、姿态无关的变化检测方法。在3DGS中引入独立变化通道(含变化不透明度因子)构建变化感知3D场景表征。仅用5张变化后图像即可生成变化掩码，在mIoU上提升1.7倍，F1提升1.5倍。还可为未观察视角生成变化掩码。贡献了含光照变化的真实世界数据集。适用时还应保留论文给出的边界：依赖参考场景3DGS的质量；对剧烈外观变化场景(如极暗)COLMAP可能注册失败；变化掩码精度受限于2D特征比较。

\section{场景图与语义建图}

场景图方法把高斯从“可分割物体集合”提升为“带关系和可更新语义的地图”。它们通常先用2D开放词汇检测和分割得到候选对象，再用3DGS提供几何、视角合成和局部更新能力，最终形成可供机器人或空间智能系统查询的结构。
% METHOD: 2502.15309
\paragraph{DynamicGSG（2025.02）\cite{2502_15309}}
DynamicGSG把3DGS从可渲染地图推进到可更新场景图。系统前端用YOLO-World做开放词汇检测，SAM产生实例掩码，CLIP提取对象语义；这些2D观测再通过几何重投影和语义相似性关联到3D地图中的对象节点。与只在图像层面维护场景图不同，DynamicGSG为高斯加入低维实例特征和对象标识，使一组高斯可以被稳定地看作同一个物体。

后端维护两类结构：一类是高斯对象集合，负责外观、几何和局部更新；另一类是层次化场景图，记录对象之间的空间关系和语义关系，并可由视觉语言模型补全关系描述。环境发生变化时，系统不必重新训练整张3DGS，而是把新观测与已有对象匹配，只更新受影响的高斯集合和相关图节点。相对静态语义地图，DynamicGSG的关键增量是把“可渲染、可局部优化”的3DGS性质接入持续场景理解，使机器人或空间智能系统能在动态环境中维护对象级记忆。
\textbf{方法信息：} 代码：\url{https://github.com/geluzhou/dynamic-gsg}。
\textbf{设计动机与边界：} 相对原始 3DGS 以独立高斯和图像重建误差为主的设定，DynamicGSG 把问题落在以下可实现的改动上：DynamicGSG利用3D Gaussian Splatting构建动态高保真开放词汇场景图。系统使用YOLO-World进行开放词汇物体检测，Segment Anything生成实例掩码，CLIP提取语义特征。核心流程包括：(1)物体识别与3D-2D关联，通过联合几何和语义相似性匹配将检测结果与地图物体关联；(2)高斯物体分组，使用联合特征损失监督高斯实例分组并优化高斯地图；(3)层次化场景图构建，使用视觉语言模型分析物体间空间和语义关系；(4)局部动态更新，根据真实环境变化局部更新高斯场景图。每个高斯额外添加低维实例特征和标识符参数。其设计目的不是单独增加一个模块，而是用上述表示、监督或渲染约束消除原流程在目标任务中的具体失配。提出DynamicGSG系统，利用3DGS构建动态高保真开放词汇场景图，支持根据环境变化局部更新高斯场景图，实现长期环境适应。适用时还应保留论文给出的边界：对快速剧烈的环境变化可能无法及时响应；依赖视觉语言模型的检测质量；计算开销需进一步优化以满足实时性要求。
% METHOD: 2603.07660
\paragraph{Holi-Spatial（2026.03）\cite{2603_07660}}
Holi-Spatial把3DGS用作大规模空间智能数据构建引擎。它不是为单个场景做一次分割，而是把原始视频转成多层级3D监督：优化3DGS场景、渲染深度、2D开放集掩码、3D边界框、实例描述和空间问答。几何阶段先用SfM求相机，用Depth-Anything-V3初始化密集点云，再通过3DGS逐场景优化；PGSR、2DGS式几何正则用于强化多视角深度一致性，减少漂浮高斯。

感知阶段，系统均匀采样关键帧，用Gemini生成开放词汇类别并维护动态类别记忆，随后引导SAM3产生掩码；优化后的3DGS渲染深度把2D掩码反投影到3D，形成初始实例。场景级阶段再在3D空间合并、过滤、验证实例，并生成文本描述和空间QA。相对DynamicGSG维护单个环境的在线场景图，Holi-Spatial关注数据飞轮：让3DGS提供几何对齐和视角补全，把2D基础模型输出变成可训练空间智能模型的大规模3D标注。论文发布的Holi-Spatial-4M包含12K优化3DGS场景、1.3M 2D掩码、320K 3D框和1.2M空间QA。
\textbf{设计动机与边界：} 相对原始 3DGS 以独立高斯和图像重建误差为主的设定，Holi-Spatial 把问题落在以下可实现的改动上：Holi-Spatial包含三个阶段：(1)几何优化——使用SfM解析相机参数，利用Depth-Anything-V3初始化密集点云，再通过3DGS进行逐场景优化，集成几何正则化（源自PGSR、2DGS等方法）强制多视角深度一致性，消除浮点伪影；(2)图像级感知——均匀采样关键帧，使用Gemini-3-Pro生成开放词汇类别（维护动态类别标签记忆M\_t确保语义一致性），引导SAM3产生高质量开放集掩码，利用优化后的3DGS渲染深度图将2D掩码反投影到3D空间；(3)场景级精化——在3D中合并、过滤、验证和描述实例，生成密集可靠的空间标注。整个流程实现正向数据飞轮，标注质量可超越人工标注。其设计目的不是单独增加一个模块，而是用上述表示、监督或渲染约束消除原流程在目标任务中的具体失配。提出Holi-Spatial，首个全自动大规模空间感知多模态数据集构建框架，从原始视频输入无需人工干预即可生成包含3DGS重建、深度图、2D掩码、3D边界框、实例描述和空间QA对的多层级空间数据。构建了Holi-Spatial-4M数据集（12K优化3DGS场景、1.3M 2D掩码、320K 3D边界框、1.2M空间QA对），在数据质量和下游任务性能上显著超越现有方法。适用时还应保留论文给出的边界：依赖多个基础模型（Depth-Anything、SAM3、Gemini）的推理质量；自动标注可能引入累积误差；数据集偏向网络视频的分布。
\textbf{方法信息：} 代码：\url{https://github.com/visionary-laboratory/holi-spatial}；GitHub 关注度较高，约 339 stars。



\section{占用网格与语义场景补全}

占用估计把高斯作为从图像到 3D 体素的中间表示。相比直接预测密集体素，高斯更稀疏，也更容易通过可微 splatting 接受 2D 伪标签、深度监督和时序监督。这里的增量不是追求照片级渲染，而是让高斯先承载位置、语义和运动，再投影或离散化为占用网格；因此它位于 3DGS 和自动驾驶语义建图之间。


% METHOD: 2401.01970
\paragraph{FMGS（语义、语言与开放词汇理解）（2024.01）}
FMGS~\cite{2401_01970} 关注的问题是：把 3DGS 接入语义、语言或场景理解任务 FMGS 是宽召回补全中发现的论文，主要根据标题、摘要和 arXiv 类别判定其与 3DGS 的关系。
相对原始 3DGS，FMGS 的变化不是只替换损失函数，而是把表示、初始化、约束或网络预测路径重新组织到任务流程中：该工作被召回是因为论文文本包含 Gaussian/3D Gaussian 相关表述；当前分类先确定其任务域和与 3DGS 的相关性，后续证据流水线再补代码、企业和引用信息。
\textbf{实现细节：} 具体实现可以按输入、核心模块和渲染/优化路径来读。输入端首先取得任务所需的图像、相机、文本、时间序列、语义或几何线索，并把这些线索转换成可以约束高斯中心、尺度、旋转、不透明度、颜色或附加特征的条件。核心结构中，方法使用的主要机制是：从摘要层面看，FMGS 围绕“把 3DGS 接入语义、语言或场景理解任务”展开。若其正文或代码后续可读，应进一步核对它具体改动的是表示、初始化、渲染、优化、损失、动态建模、压缩或应用接口中的哪一环。 训练或推理时，这些模块通常先给出初始化或中间特征，再通过可微 Gaussian rasterization、重建损失、感知/语义约束、几何正则或任务损失回传到高斯参数；如果方法包含前馈网络，则网络输出的是高斯属性、变形场、压缩码、语义特征或控制变量，而不是只输出最终 2D 图像。渲染路径仍以显式高斯投影、深度排序和 alpha 合成为基础，只是在采样、属性预测、局部编辑、动态变形、压缩编码或下游任务接口处加入额外约束。这样写入本章时，它与前后方法的关系是：相对原始 3DGS 和本章已有锚点方法，它可以理解为 3DGS extends\_or\_uses 宽召回分类阶段的关系判断，后续需要按论文方法部分复核。
\textbf{证据与边界：} 本轮召回分类尚未逐项核对实验表格；后续证据补全阶段需要读取论文源码/PDF或项目页确认数据集、指标和对比方法。 该条来自宽召回批次；源码下载受 arXiv e-print 网络限制，部分论文暂时只能依据标题、摘要和类别分类，需要在入可信书前做人工/LLM精读复核。 本书只采用本地 paper pack、数据库分类、源码/引用登记和论文摘要中可核验的信息；没有本地证据支撑的速度、排行榜或开源完整性不在正文中强化。\textbf{方法信息：} 公开链接：\url{https://github.com/xingxingzuo/fmgs}，完整性按本地证据保守处理；引用信号：120 citations (4.0/month, domain-relative cutoff 4.0/month, semantic\_scholar)。
\textbf{设计动机与边界：} 相对原始 3DGS 以独立高斯和图像重建误差为主的设定，FMGS 把问题落在以下可实现的改动上：从摘要层面看，FMGS 围绕“把 3DGS 接入语义、语言或场景理解任务”展开。若其正文或代码后续可读，应进一步核对它具体改动的是表示、初始化、渲染、优化、损失、动态建模、压缩或应用接口中的哪一环。其设计目的不是单独增加一个模块，而是用上述表示、监督或渲染约束消除原流程在目标任务中的具体失配。该工作被召回是因为论文文本包含 Gaussian/3D Gaussian 相关表述；当前分类先确定其任务域和与 3DGS 的相关性，后续证据流水线再补代码、企业和引用信息。适用时还应保留论文给出的边界：该条来自宽召回批次；源码下载受 arXiv e-print 网络限制，部分论文暂时只能依据标题、摘要和类别分类，需要在入可信书前做人工/LLM精读复核。
% METHOD: 2402.03246
\paragraph{SGS-SLAM（语义、语言与开放词汇理解）（2024.02）}
SGS-SLAM~\cite{2402_03246} 关注的问题是：把 3DGS 接入语义、语言或场景理解任务 SGS-SLAM 是宽召回补全中发现的论文，主要根据标题、摘要和 arXiv 类别判定其与 3DGS 的关系。
相对原始 3DGS，SGS-SLAM 的变化不是只替换损失函数，而是把表示、初始化、约束或网络预测路径重新组织到任务流程中：该工作被召回是因为论文文本包含 Gaussian/3D Gaussian 相关表述；当前分类先确定其任务域和与 3DGS 的相关性，后续证据流水线再补代码、企业和引用信息。
\textbf{实现细节：} 具体实现可以按输入、核心模块和渲染/优化路径来读。输入端首先取得任务所需的图像、相机、文本、时间序列、语义或几何线索，并把这些线索转换成可以约束高斯中心、尺度、旋转、不透明度、颜色或附加特征的条件。核心结构中，方法使用的主要机制是：从摘要层面看，SGS-SLAM 围绕“把 3DGS 接入语义、语言或场景理解任务”展开。若其正文或代码后续可读，应进一步核对它具体改动的是表示、初始化、渲染、优化、损失、动态建模、压缩或应用接口中的哪一环。 训练或推理时，这些模块通常先给出初始化或中间特征，再通过可微 Gaussian rasterization、重建损失、感知/语义约束、几何正则或任务损失回传到高斯参数；如果方法包含前馈网络，则网络输出的是高斯属性、变形场、压缩码、语义特征或控制变量，而不是只输出最终 2D 图像。渲染路径仍以显式高斯投影、深度排序和 alpha 合成为基础，只是在采样、属性预测、局部编辑、动态变形、压缩编码或下游任务接口处加入额外约束。这样写入本章时，它与前后方法的关系是：相对原始 3DGS 和本章已有锚点方法，它可以理解为 3DGS extends\_or\_uses 宽召回分类阶段的关系判断，后续需要按论文方法部分复核。
\textbf{证据与边界：} 本轮召回分类尚未逐项核对实验表格；后续证据补全阶段需要读取论文源码/PDF或项目页确认数据集、指标和对比方法。 该条来自宽召回批次；源码下载受 arXiv e-print 网络限制，部分论文暂时只能依据标题、摘要和类别分类，需要在入可信书前做人工/LLM精读复核。 本书只采用本地 paper pack、数据库分类、源码/引用登记和论文摘要中可核验的信息；没有本地证据支撑的速度、排行榜或开源完整性不在正文中强化。\textbf{方法信息：} 代码：\url{https://github.com/shuhongll/sgs-slam}；引用信号：171 citations (5.9/month, domain-relative cutoff 4.0/month, semantic\_scholar)；代码关注：shuhongll/sgs-slam has 516 stars (17.7/month, domain-relative cutoff 21.5/month)。
\textbf{原始方法补充：} SGS-SLAM 针对原始 3DGS 只能离线拟合一组颜色高斯、无法同时定位和维护语义地图的问题，把每个高斯扩展为可联合优化的多通道表示。系统输入连续 RGB-D 帧；跟踪阶段固定地图参数，以当前帧渲染结果估计相机位姿，建图阶段固定已估姿态并更新高斯的几何、外观和语义通道。这个交替流程把位姿优化和地图优化分开，避免两者相互吸收误差；最终由高斯光栅化同时输出 RGB、深度和语义图，以支持后续场景操作。其设计动机是让语义监督与稠密 SLAM 共用显式高斯地图，而不把理解结果附加在一次离线重建之后。
% METHOD: 2403.07494
\paragraph{SemGauss-SLAM（语义、语言与开放词汇理解）（2024.03）}
SemGauss-SLAM~\cite{2403_07494} 关注的问题是：把 3DGS 接入语义、语言或场景理解任务 SemGauss-SLAM 是宽召回补全中发现的论文，主要根据标题、摘要和 arXiv 类别判定其与 3DGS 的关系。
相对原始 3DGS，SemGauss-SLAM 的变化不是只替换损失函数，而是把表示、初始化、约束或网络预测路径重新组织到任务流程中：该工作被召回是因为论文文本包含 Gaussian/3D Gaussian 相关表述；当前分类先确定其任务域和与 3DGS 的相关性，后续证据流水线再补代码、企业和引用信息。
\textbf{实现细节：} 具体实现可以按输入、核心模块和渲染/优化路径来读。输入端首先取得任务所需的图像、相机、文本、时间序列、语义或几何线索，并把这些线索转换成可以约束高斯中心、尺度、旋转、不透明度、颜色或附加特征的条件。核心结构中，方法使用的主要机制是：从摘要层面看，SemGauss-SLAM 围绕“把 3DGS 接入语义、语言或场景理解任务”展开。若其正文或代码后续可读，应进一步核对它具体改动的是表示、初始化、渲染、优化、损失、动态建模、压缩或应用接口中的哪一环。 训练或推理时，这些模块通常先给出初始化或中间特征，再通过可微 Gaussian rasterization、重建损失、感知/语义约束、几何正则或任务损失回传到高斯参数；如果方法包含前馈网络，则网络输出的是高斯属性、变形场、压缩码、语义特征或控制变量，而不是只输出最终 2D 图像。渲染路径仍以显式高斯投影、深度排序和 alpha 合成为基础，只是在采样、属性预测、局部编辑、动态变形、压缩编码或下游任务接口处加入额外约束。这样写入本章时，它与前后方法的关系是：相对原始 3DGS 和本章已有锚点方法，它可以理解为 3DGS extends\_or\_uses 宽召回分类阶段的关系判断，后续需要按论文方法部分复核。
\textbf{证据与边界：} 本轮召回分类尚未逐项核对实验表格；后续证据补全阶段需要读取论文源码/PDF或项目页确认数据集、指标和对比方法。 该条来自宽召回批次；源码下载受 arXiv e-print 网络限制，部分论文暂时只能依据标题、摘要和类别分类，需要在入可信书前做人工/LLM精读复核。 本书只采用本地 paper pack、数据库分类、源码/引用登记和论文摘要中可核验的信息；没有本地证据支撑的速度、排行榜或开源完整性不在正文中强化。\textbf{方法信息：} 代码：\url{https://github.com/irmvlab/semgauss-slam}；引用信号：84 citations (3.0/month, domain-relative cutoff 4.0/month, semantic\_scholar)。
\textbf{设计动机与边界：} 相对原始 3DGS 以独立高斯和图像重建误差为主的设定，SemGauss-SLAM 把问题落在以下可实现的改动上：从摘要层面看，SemGauss-SLAM 围绕“把 3DGS 接入语义、语言或场景理解任务”展开。若其正文或代码后续可读，应进一步核对它具体改动的是表示、初始化、渲染、优化、损失、动态建模、压缩或应用接口中的哪一环。其设计目的不是单独增加一个模块，而是用上述表示、监督或渲染约束消除原流程在目标任务中的具体失配。该工作被召回是因为论文文本包含 Gaussian/3D Gaussian 相关表述；当前分类先确定其任务域和与 3DGS 的相关性，后续证据流水线再补代码、企业和引用信息。适用时还应保留论文给出的边界：该条来自宽召回批次；源码下载受 arXiv e-print 网络限制，部分论文暂时只能依据标题、摘要和类别分类，需要在入可信书前做人工/LLM精读复核。
% METHOD: 2403.09637
\paragraph{GaussianGrasper（语义、语言与开放词汇理解）（2024.03）}
GaussianGrasper~\cite{2403_09637} 关注的问题是：把 3DGS 接入语义、语言或场景理解任务 GaussianGrasper 是宽召回补全中发现的论文，主要根据标题、摘要和 arXiv 类别判定其与 3DGS 的关系。
相对原始 3DGS，GaussianGrasper 的变化不是只替换损失函数，而是把表示、初始化、约束或网络预测路径重新组织到任务流程中：该工作被召回是因为论文文本包含 Gaussian/3D Gaussian 相关表述；当前分类先确定其任务域和与 3DGS 的相关性，后续证据流水线再补代码、企业和引用信息。
\textbf{实现细节：} 具体实现可以按输入、核心模块和渲染/优化路径来读。输入端首先取得任务所需的图像、相机、文本、时间序列、语义或几何线索，并把这些线索转换成可以约束高斯中心、尺度、旋转、不透明度、颜色或附加特征的条件。核心结构中，方法使用的主要机制是：从摘要层面看，GaussianGrasper 围绕“把 3DGS 接入语义、语言或场景理解任务”展开。若其正文或代码后续可读，应进一步核对它具体改动的是表示、初始化、渲染、优化、损失、动态建模、压缩或应用接口中的哪一环。 训练或推理时，这些模块通常先给出初始化或中间特征，再通过可微 Gaussian rasterization、重建损失、感知/语义约束、几何正则或任务损失回传到高斯参数；如果方法包含前馈网络，则网络输出的是高斯属性、变形场、压缩码、语义特征或控制变量，而不是只输出最终 2D 图像。渲染路径仍以显式高斯投影、深度排序和 alpha 合成为基础，只是在采样、属性预测、局部编辑、动态变形、压缩编码或下游任务接口处加入额外约束。这样写入本章时，它与前后方法的关系是：相对原始 3DGS 和本章已有锚点方法，它可以理解为 3DGS extends\_or\_uses 宽召回分类阶段的关系判断，后续需要按论文方法部分复核。
\textbf{证据与边界：} 本轮召回分类尚未逐项核对实验表格；后续证据补全阶段需要读取论文源码/PDF或项目页确认数据集、指标和对比方法。 该条来自宽召回批次；源码下载受 arXiv e-print 网络限制，部分论文暂时只能依据标题、摘要和类别分类，需要在入可信书前做人工/LLM精读复核。 本书只采用本地 paper pack、数据库分类、源码/引用登记和论文摘要中可核验的信息；没有本地证据支撑的速度、排行榜或开源完整性不在正文中强化。\textbf{方法信息：} 代码：\url{https://github.com/mrsecant/gaussiangrasper}；引用信号：112 citations (4.0/month, domain-relative cutoff 4.0/month, semantic\_scholar)；代码关注：mrsecant/gaussiangrasper has 174 stars (6.2/month, domain-relative cutoff 21.5/month)。
\textbf{设计动机与边界：} 相对原始 3DGS 以独立高斯和图像重建误差为主的设定，GaussianGrasper 把问题落在以下可实现的改动上：从摘要层面看，GaussianGrasper 围绕“把 3DGS 接入语义、语言或场景理解任务”展开。若其正文或代码后续可读，应进一步核对它具体改动的是表示、初始化、渲染、优化、损失、动态建模、压缩或应用接口中的哪一环。其设计目的不是单独增加一个模块，而是用上述表示、监督或渲染约束消除原流程在目标任务中的具体失配。该工作被召回是因为论文文本包含 Gaussian/3D Gaussian 相关表述；当前分类先确定其任务域和与 3DGS 的相关性，后续证据流水线再补代码、企业和引用信息。适用时还应保留论文给出的边界：该条来自宽召回批次；源码下载受 arXiv e-print 网络限制，部分论文暂时只能依据标题、摘要和类别分类，需要在入可信书前做人工/LLM精读复核。
% METHOD: 2403.11625
\paragraph{GaussNav（语义、语言与开放词汇理解）（2024.03）}
GaussNav~\cite{2403_11625} 关注的问题是：把 3DGS 接入语义、语言或场景理解任务 GaussNav 是宽召回补全中发现的论文，主要根据标题、摘要和 arXiv 类别判定其与 3DGS 的关系。
相对原始 3DGS，GaussNav 的变化不是只替换损失函数，而是把表示、初始化、约束或网络预测路径重新组织到任务流程中：该工作被召回是因为论文文本包含 Gaussian/3D Gaussian 相关表述；当前分类先确定其任务域和与 3DGS 的相关性，后续证据流水线再补代码、企业和引用信息。
\textbf{实现细节：} 具体实现可以按输入、核心模块和渲染/优化路径来读。输入端首先取得任务所需的图像、相机、文本、时间序列、语义或几何线索，并把这些线索转换成可以约束高斯中心、尺度、旋转、不透明度、颜色或附加特征的条件。核心结构中，方法使用的主要机制是：从摘要层面看，GaussNav 围绕“把 3DGS 接入语义、语言或场景理解任务”展开。若其正文或代码后续可读，应进一步核对它具体改动的是表示、初始化、渲染、优化、损失、动态建模、压缩或应用接口中的哪一环。 训练或推理时，这些模块通常先给出初始化或中间特征，再通过可微 Gaussian rasterization、重建损失、感知/语义约束、几何正则或任务损失回传到高斯参数；如果方法包含前馈网络，则网络输出的是高斯属性、变形场、压缩码、语义特征或控制变量，而不是只输出最终 2D 图像。渲染路径仍以显式高斯投影、深度排序和 alpha 合成为基础，只是在采样、属性预测、局部编辑、动态变形、压缩编码或下游任务接口处加入额外约束。这样写入本章时，它与前后方法的关系是：相对原始 3DGS 和本章已有锚点方法，它可以理解为 3DGS extends\_or\_uses 宽召回分类阶段的关系判断，后续需要按论文方法部分复核。
\textbf{证据与边界：} 本轮召回分类尚未逐项核对实验表格；后续证据补全阶段需要读取论文源码/PDF或项目页确认数据集、指标和对比方法。 该条来自宽召回批次；源码下载受 arXiv e-print 网络限制，部分论文暂时只能依据标题、摘要和类别分类，需要在入可信书前做人工/LLM精读复核。 本书只采用本地 paper pack、数据库分类、源码/引用登记和论文摘要中可核验的信息；没有本地证据支撑的速度、排行榜或开源完整性不在正文中强化。\textbf{方法信息：} 代码：\url{https://github.com/xiaohanlei/gaussnav}；引用信号：67 citations (2.4/month, domain-relative cutoff 4.0/month, semantic\_scholar)；代码关注：xiaohanlei/gaussnav has 218 stars (7.9/month, domain-relative cutoff 21.5/month)。
\textbf{设计动机与边界：} 相对原始 3DGS 以独立高斯和图像重建误差为主的设定，GaussNav 把问题落在以下可实现的改动上：从摘要层面看，GaussNav 围绕“把 3DGS 接入语义、语言或场景理解任务”展开。若其正文或代码后续可读，应进一步核对它具体改动的是表示、初始化、渲染、优化、损失、动态建模、压缩或应用接口中的哪一环。其设计目的不是单独增加一个模块，而是用上述表示、监督或渲染约束消除原流程在目标任务中的具体失配。该工作被召回是因为论文文本包含 Gaussian/3D Gaussian 相关表述；当前分类先确定其任务域和与 3DGS 的相关性，后续证据流水线再补代码、企业和引用信息。适用时还应保留论文给出的边界：该条来自宽召回批次；源码下载受 arXiv e-print 网络限制，部分论文暂时只能依据标题、摘要和类别分类，需要在入可信书前做人工/LLM精读复核。
% METHOD: 2403.15624
\paragraph{Semantic Gaussians（语义、语言与开放词汇理解）（2024.03）}
Semantic Gaussians~\cite{2403_15624} 关注的问题是：把 3DGS 接入语义、语言或场景理解任务 Semantic Gaussians 是宽召回补全中发现的论文，主要根据标题、摘要和 arXiv 类别判定其与 3DGS 的关系。
相对原始 3DGS，Semantic Gaussians 的变化不是只替换损失函数，而是把表示、初始化、约束或网络预测路径重新组织到任务流程中：该工作被召回是因为论文文本包含 Gaussian/3D Gaussian 相关表述；当前分类先确定其任务域和与 3DGS 的相关性，后续证据流水线再补代码、企业和引用信息。
\textbf{实现细节：} 具体实现可以按输入、核心模块和渲染/优化路径来读。输入端首先取得任务所需的图像、相机、文本、时间序列、语义或几何线索，并把这些线索转换成可以约束高斯中心、尺度、旋转、不透明度、颜色或附加特征的条件。核心结构中，方法使用的主要机制是：从摘要层面看，Semantic Gaussians 围绕“把 3DGS 接入语义、语言或场景理解任务”展开。若其正文或代码后续可读，应进一步核对它具体改动的是表示、初始化、渲染、优化、损失、动态建模、压缩或应用接口中的哪一环。 训练或推理时，这些模块通常先给出初始化或中间特征，再通过可微 Gaussian rasterization、重建损失、感知/语义约束、几何正则或任务损失回传到高斯参数；如果方法包含前馈网络，则网络输出的是高斯属性、变形场、压缩码、语义特征或控制变量，而不是只输出最终 2D 图像。渲染路径仍以显式高斯投影、深度排序和 alpha 合成为基础，只是在采样、属性预测、局部编辑、动态变形、压缩编码或下游任务接口处加入额外约束。这样写入本章时，它与前后方法的关系是：相对原始 3DGS 和本章已有锚点方法，它可以理解为 3DGS extends\_or\_uses 宽召回分类阶段的关系判断，后续需要按论文方法部分复核。
\textbf{证据与边界：} 本轮召回分类尚未逐项核对实验表格；后续证据补全阶段需要读取论文源码/PDF或项目页确认数据集、指标和对比方法。 该条来自宽召回批次；源码下载受 arXiv e-print 网络限制，部分论文暂时只能依据标题、摘要和类别分类，需要在入可信书前做人工/LLM精读复核。 本书只采用本地 paper pack、数据库分类、源码/引用登记和论文摘要中可核验的信息；没有本地证据支撑的速度、排行榜或开源完整性不在正文中强化。\textbf{方法信息：} 代码：\url{https://github.com/sharinka0715/semantic-gaussians}；引用信号：93 citations (3.4/month, domain-relative cutoff 4.0/month, semantic\_scholar)；代码关注：sharinka0715/semantic-gaussians has 265 stars (9.6/month, domain-relative cutoff 21.5/month)。
\textbf{设计动机与边界：} 相对原始 3DGS 以独立高斯和图像重建误差为主的设定，Semantic Gaussians 把问题落在以下可实现的改动上：从摘要层面看，Semantic Gaussians 围绕“把 3DGS 接入语义、语言或场景理解任务”展开。若其正文或代码后续可读，应进一步核对它具体改动的是表示、初始化、渲染、优化、损失、动态建模、压缩或应用接口中的哪一环。其设计目的不是单独增加一个模块，而是用上述表示、监督或渲染约束消除原流程在目标任务中的具体失配。该工作被召回是因为论文文本包含 Gaussian/3D Gaussian 相关表述；当前分类先确定其任务域和与 3DGS 的相关性，后续证据流水线再补代码、企业和引用信息。适用时还应保留论文给出的边界：该条来自宽召回批次；源码下载受 arXiv e-print 网络限制，部分论文暂时只能依据标题、摘要和类别分类，需要在入可信书前做人工/LLM精读复核。
% METHOD: 2404.01223
\paragraph{Feature Splatting（语义、语言与开放词汇理解）（2024.04）}
Feature Splatting~\cite{2404_01223} 关注的问题是：把 3DGS 接入语义、语言或场景理解任务 Feature Splatting 是宽召回补全中发现的论文，主要根据标题、摘要和 arXiv 类别判定其与 3DGS 的关系。
相对原始 3DGS，Feature Splatting 的变化不是只替换损失函数，而是把表示、初始化、约束或网络预测路径重新组织到任务流程中：该工作被召回是因为论文文本包含 Gaussian/3D Gaussian 相关表述；当前分类先确定其任务域和与 3DGS 的相关性，后续证据流水线再补代码、企业和引用信息。
\textbf{实现细节：} 具体实现可以按输入、核心模块和渲染/优化路径来读。输入端首先取得任务所需的图像、相机、文本、时间序列、语义或几何线索，并把这些线索转换成可以约束高斯中心、尺度、旋转、不透明度、颜色或附加特征的条件。核心结构中，方法使用的主要机制是：从摘要层面看，Feature Splatting 围绕“把 3DGS 接入语义、语言或场景理解任务”展开。若其正文或代码后续可读，应进一步核对它具体改动的是表示、初始化、渲染、优化、损失、动态建模、压缩或应用接口中的哪一环。 训练或推理时，这些模块通常先给出初始化或中间特征，再通过可微 Gaussian rasterization、重建损失、感知/语义约束、几何正则或任务损失回传到高斯参数；如果方法包含前馈网络，则网络输出的是高斯属性、变形场、压缩码、语义特征或控制变量，而不是只输出最终 2D 图像。渲染路径仍以显式高斯投影、深度排序和 alpha 合成为基础，只是在采样、属性预测、局部编辑、动态变形、压缩编码或下游任务接口处加入额外约束。这样写入本章时，它与前后方法的关系是：相对原始 3DGS 和本章已有锚点方法，它可以理解为 3DGS extends\_or\_uses 宽召回分类阶段的关系判断，后续需要按论文方法部分复核。
\textbf{证据与边界：} 本轮召回分类尚未逐项核对实验表格；后续证据补全阶段需要读取论文源码/PDF或项目页确认数据集、指标和对比方法。 该条来自宽召回批次；源码下载受 arXiv e-print 网络限制，部分论文暂时只能依据标题、摘要和类别分类，需要在入可信书前做人工/LLM精读复核。 本书只采用本地 paper pack、数据库分类、源码/引用登记和论文摘要中可核验的信息；没有本地证据支撑的速度、排行榜或开源完整性不在正文中强化。\textbf{方法信息：} 公开链接：\url{https://github.com/jmiaodong/feature-3d-gaussian}，完整性按本地证据保守处理；引用信号：63 citations (2.3/month, domain-relative cutoff 4.0/month, semantic\_scholar)。
\textbf{设计动机与边界：} 相对原始 3DGS 以独立高斯和图像重建误差为主的设定，Feature Splatting 把问题落在以下可实现的改动上：从摘要层面看，Feature Splatting 围绕“把 3DGS 接入语义、语言或场景理解任务”展开。若其正文或代码后续可读，应进一步核对它具体改动的是表示、初始化、渲染、优化、损失、动态建模、压缩或应用接口中的哪一环。其设计目的不是单独增加一个模块，而是用上述表示、监督或渲染约束消除原流程在目标任务中的具体失配。该工作被召回是因为论文文本包含 Gaussian/3D Gaussian 相关表述；当前分类先确定其任务域和与 3DGS 的相关性，后续证据流水线再补代码、企业和引用信息。适用时还应保留论文给出的边界：该条来自宽召回批次；源码下载受 arXiv e-print 网络限制，部分论文暂时只能依据标题、摘要和类别分类，需要在入可信书前做人工/LLM精读复核。
% METHOD: 2407.20229
\paragraph{Improving 2D Feature Representations by 3D-Aware Fine-Tuning（语义、语言与开放词汇理解）（2024.07）}
Improving 2D Feature Representations by 3D-Aware Fine-Tuning~\cite{2407_20229} 关注的问题是：把 3DGS 接入语义、语言或场景理解任务 Improving 2D Feature Representations by 3D-Aware Fine-Tuning 是宽召回补全中发现的论文，主要根据标题、摘要和 arXiv 类别判定其与 3DGS 的关系。
相对原始 3DGS，Improving 2D Feature Representations by 3D-Aware Fine-Tuning 的变化不是只替换损失函数，而是把表示、初始化、约束或网络预测路径重新组织到任务流程中：该工作被召回是因为论文文本包含 Gaussian/3D Gaussian 相关表述；当前分类先确定其任务域和与 3DGS 的相关性，后续证据流水线再补代码、企业和引用信息。
\textbf{实现细节：} 具体实现可以按输入、核心模块和渲染/优化路径来读。输入端首先取得任务所需的图像、相机、文本、时间序列、语义或几何线索，并把这些线索转换成可以约束高斯中心、尺度、旋转、不透明度、颜色或附加特征的条件。核心结构中，方法使用的主要机制是：从摘要层面看，Improving 2D Feature Representations by 3D-Aware Fine-Tuning 围绕“把 3DGS 接入语义、语言或场景理解任务”展开。若其正文或代码后续可读，应进一步核对它具体改动的是表示、初始化、渲染、优化、损失、动态建模、压缩或应用接口中的哪一环。 训练或推理时，这些模块通常先给出初始化或中间特征，再通过可微 Gaussian rasterization、重建损失、感知/语义约束、几何正则或任务损失回传到高斯参数；如果方法包含前馈网络，则网络输出的是高斯属性、变形场、压缩码、语义特征或控制变量，而不是只输出最终 2D 图像。渲染路径仍以显式高斯投影、深度排序和 alpha 合成为基础，只是在采样、属性预测、局部编辑、动态变形、压缩编码或下游任务接口处加入额外约束。这样写入本章时，它与前后方法的关系是：相对原始 3DGS 和本章已有锚点方法，它可以理解为 3DGS extends\_or\_uses 宽召回分类阶段的关系判断，后续需要按论文方法部分复核。
\textbf{证据与边界：} 本轮召回分类尚未逐项核对实验表格；后续证据补全阶段需要读取论文源码/PDF或项目页确认数据集、指标和对比方法。 该条来自宽召回批次；源码下载受 arXiv e-print 网络限制，部分论文暂时只能依据标题、摘要和类别分类，需要在入可信书前做人工/LLM精读复核。 本书只采用本地 paper pack、数据库分类、源码/引用登记和论文摘要中可核验的信息；没有本地证据支撑的速度、排行榜或开源完整性不在正文中强化。\textbf{方法信息：} 引用信号：75 citations (3.2/month, domain-relative cutoff 4.0/month, semantic\_scholar)。
\textbf{设计动机与边界：} 相对原始 3DGS 以独立高斯和图像重建误差为主的设定，Improving 2D Feature Representations by 3D-Aware Fine-Tuning 把问题落在以下可实现的改动上：从摘要层面看，Improving 2D Feature Representations by 3D-Aware Fine-Tuning 围绕“把 3DGS 接入语义、语言或场景理解任务”展开。若其正文或代码后续可读，应进一步核对它具体改动的是表示、初始化、渲染、优化、损失、动态建模、压缩或应用接口中的哪一环。其设计目的不是单独增加一个模块，而是用上述表示、监督或渲染约束消除原流程在目标任务中的具体失配。该工作被召回是因为论文文本包含 Gaussian/3D Gaussian 相关表述；当前分类先确定其任务域和与 3DGS 的相关性，后续证据流水线再补代码、企业和引用信息。适用时还应保留论文给出的边界：该条来自宽召回批次；源码下载受 arXiv e-print 网络限制，部分论文暂时只能依据标题、摘要和类别分类，需要在入可信书前做人工/LLM精读复核。
% METHOD: 2502.17288
\paragraph{GaussianFlowOcc（2025.02）\cite{2502_17288}}
GaussianFlowOcc把自动驾驶占用预测中的密集体素中间层换成稀疏3D高斯。传统占用网络直接从多相机图像预测体素语义，监督通常稀疏且昂贵；GaussianFlowOcc先让Gaussian Transformer从图像特征中生成一组带位置、不透明度、尺度、旋转和语义logits的高斯，再用可微splatting把这些高斯投回相机视图。这样，Grounded-SAM产生的2D语义伪标签和Metric3D产生的深度伪标签都可以直接监督高斯，而不必先获得完整3D人工标注。

模型内部的Gaussian Transformer包含图像到高斯的交互和高斯之间的信息传播：GICA等跨注意力模块从多视角图像特征中采样与每个高斯相关的信息，induced attention让高斯集合内部交换上下文，多个预测头输出几何、语义和不透明度。为了处理动态物体，Temporal Gaussian Splatting再为高斯估计到相邻帧的3D flow，把当前高斯移动到邻近时刻后重新渲染，用时序监督减少单帧伪标签噪声。训练后，这些稀疏高斯可以离散化成占用网格进行评测。相对直接体素化的占用网络，它的关键增量是把3DGS变成弱监督桥梁：2D基础模型负责提供廉价监督，高斯负责在3D中融合和传播监督。
\textbf{实现细节：} GaussianFlowOcc：稀疏高斯时序占用估计。GaussianFlowOcc（2025.02） 从表示和训练两个层面同时革新占用估计流程。在表示层面，它将传统密集体素网格替换为 $N$ 个稀疏 3D 高斯基元（每个基元参数化为位置、不透明度、尺度、旋转四元数和语义 logits），消除对 3D 体素标签的依赖，转而通过高斯泼溅将估计结果渲染回 2D 相机空间进行监督。在模型架构层面，GaussianFlowOcc 设计了高斯变换器（Gaussian Transformer），由三个迭代块组成，每个块包含三类注意力：（1）与图像特征的交叉注意力（使用 $M$ 个诱导点实现线性复杂度）；（2）高斯基元之间的自注意力；（3）与前帧高斯的时序注意力。高斯头部（Gaussian Heads）为每个基元估计尺度、旋转、不透明度和语义 logits。在时序建模方面，时序模块（Temporal Module）为每个高斯计算 3D 时序光流偏移，将其与前一帧的时序邻居关联，并在高斯泼溅过程中补偿动态对象的运动，从而实现弱监督下的时序一致性。 %=====================================================================。
\textbf{设计动机与边界：} 相对原始 3DGS 以独立高斯和图像重建误差为主的设定，GaussianFlowOcc 把问题落在以下可实现的改动上：GaussianFlowOcc用稀疏3D高斯分布替代传统密集体素网格表示场景。核心架构为Gaussian Transformer，利用集合注意力机制将图像特征直接转换为3D高斯属性（位置、不透明度、尺度、旋转、语义logits）。每个高斯初始化为可学习的位置和特征。高斯通过泼溅管线渲染回输入相机空间，使用2D标签（Grounded-SAM语义伪标签和Metric3D深度）计算损失。关键创新是Temporal Gaussian Splatting：渲染相邻帧增强监督信号。时序模块为每个高斯估计到各时间帧的3D流，建模动态物体运动，减轻训练期间物体运动引起的时序不一致。支持弱监督训练，无需昂贵的密集3D体素标注。结果可轻松转换为体素体积以便与先前方法比较。其设计目的不是单独增加一个模块，而是用上述表示、监督或渲染约束消除原流程在目标任务中的具体失配。提出用稀疏3D高斯替代密集体素网格进行占用估计，引入时间高斯泼溅和时序模块处理动态物体，在nuScenes弱监督设置下显著超越所有先前方法，推理速度比SOTA快50倍。适用时还应保留论文给出的边界：稀疏高斯表示可能遗漏小物体；弱监督的精度上限受伪标签质量限制；时序流的长期预测稳定性需改进。
\textbf{方法信息：} Bosch Research 参与了这篇论文；代码：\url{https://github.com/boschresearch/gaussianflowocc}。

% METHOD: 2605.26500
\paragraph{3D Gaussian Map with Open-Set Semantic Grouping for Vision-Language Navigation（语义、语言与开放词汇理解）（2026.05）}
3D Gaussian Map with Open-Set Semantic Grouping for Vision-Language Navigation~\cite{2605_26500} 关注的问题是：把 3DGS 接入语义、语言或场景理解任务 3D Gaussian Map with Open-Set Semantic Grouping for Vision-Language Navigation 是宽召回补全中发现的论文，主要根据标题、摘要和 arXiv 类别判定其与 3DGS 的关系。
相对原始 3DGS，3D Gaussian Map with Open-Set Semantic Grouping for Vision-Language Navigation 的变化不是只替换损失函数，而是把表示、初始化、约束或网络预测路径重新组织到任务流程中：该工作被召回是因为论文文本包含 Gaussian/3D Gaussian 相关表述；当前分类先确定其任务域和与 3DGS 的相关性，后续证据流水线再补代码、企业和引用信息。
\textbf{实现细节：} 具体实现可以按输入、核心模块和渲染/优化路径来读。输入端首先取得任务所需的图像、相机、文本、时间序列、语义或几何线索，并把这些线索转换成可以约束高斯中心、尺度、旋转、不透明度、颜色或附加特征的条件。核心结构中，方法使用的主要机制是：从摘要层面看，3D Gaussian Map with Open-Set Semantic Grouping for Vision-Language Navigation 围绕“把 3DGS 接入语义、语言或场景理解任务”展开。若其正文或代码后续可读，应进一步核对它具体改动的是表示、初始化、渲染、优化、损失、动态建模、压缩或应用接口中的哪一环。 训练或推理时，这些模块通常先给出初始化或中间特征，再通过可微 Gaussian rasterization、重建损失、感知/语义约束、几何正则或任务损失回传到高斯参数；如果方法包含前馈网络，则网络输出的是高斯属性、变形场、压缩码、语义特征或控制变量，而不是只输出最终 2D 图像。渲染路径仍以显式高斯投影、深度排序和 alpha 合成为基础，只是在采样、属性预测、局部编辑、动态变形、压缩编码或下游任务接口处加入额外约束。这样写入本章时，它与前后方法的关系是：相对原始 3DGS 和本章已有锚点方法，它可以理解为 3DGS extends\_or\_uses 宽召回分类阶段的关系判断，后续需要按论文方法部分复核。
\textbf{证据与边界：} 本轮召回分类尚未逐项核对实验表格；后续证据补全阶段需要读取论文源码/PDF或项目页确认数据集、指标和对比方法。 该条来自宽召回批次；源码下载受 arXiv e-print 网络限制，部分论文暂时只能依据标题、摘要和类别分类，需要在入可信书前做人工/LLM精读复核。 本书只采用本地 paper pack、数据库分类、源码/引用登记和论文摘要中可核验的信息；没有本地证据支撑的速度、排行榜或开源完整性不在正文中强化。\textbf{方法信息：} 公开链接：\url{https://github.com/cvpr-org/author-kit}，完整性按本地证据保守处理；引用信号：13 citations (9.1/month, domain-relative cutoff 4.0/month, semantic\_scholar)。
\textbf{设计动机与边界：} 相对原始 3DGS 以独立高斯和图像重建误差为主的设定，3D Gaussian Map with Open-Set Semantic Grouping for Vision-Language Navigation 把问题落在以下可实现的改动上：从摘要层面看，3D Gaussian Map with Open-Set Semantic Grouping for Vision-Language Navigation 围绕“把 3DGS 接入语义、语言或场景理解任务”展开。若其正文或代码后续可读，应进一步核对它具体改动的是表示、初始化、渲染、优化、损失、动态建模、压缩或应用接口中的哪一环。其设计目的不是单独增加一个模块，而是用上述表示、监督或渲染约束消除原流程在目标任务中的具体失配。该工作被召回是因为论文文本包含 Gaussian/3D Gaussian 相关表述；当前分类先确定其任务域和与 3DGS 的相关性，后续证据流水线再补代码、企业和引用信息。适用时还应保留论文给出的边界：该条来自宽召回批次；源码下载受 arXiv e-print 网络限制，部分论文暂时只能依据标题、摘要和类别分类，需要在入可信书前做人工/LLM精读复核。
% METHOD: 2607.03872
\paragraph{SharpSplat（语义、语言与开放词汇理解）（2026.07）}
SharpSplat~\cite{2607_03872} 关注的问题是：把 3DGS 接入语义、语言或场景理解任务 SharpSplat 是宽召回补全中发现的论文，主要根据标题、摘要和 arXiv 类别判定其与 3DGS 的关系。
相对原始 3DGS，SharpSplat 的变化不是只替换损失函数，而是把表示、初始化、约束或网络预测路径重新组织到任务流程中：该工作被召回是因为论文文本包含 Gaussian/3D Gaussian 相关表述；当前分类先确定其任务域和与 3DGS 的相关性，后续证据流水线再补代码、企业和引用信息。
\textbf{实现细节：} 具体实现可以按输入、核心模块和渲染/优化路径来读。输入端首先取得任务所需的图像、相机、文本、时间序列、语义或几何线索，并把这些线索转换成可以约束高斯中心、尺度、旋转、不透明度、颜色或附加特征的条件。核心结构中，方法使用的主要机制是：从摘要层面看，SharpSplat 围绕“把 3DGS 接入语义、语言或场景理解任务”展开。若其正文或代码后续可读，应进一步核对它具体改动的是表示、初始化、渲染、优化、损失、动态建模、压缩或应用接口中的哪一环。 训练或推理时，这些模块通常先给出初始化或中间特征，再通过可微 Gaussian rasterization、重建损失、感知/语义约束、几何正则或任务损失回传到高斯参数；如果方法包含前馈网络，则网络输出的是高斯属性、变形场、压缩码、语义特征或控制变量，而不是只输出最终 2D 图像。渲染路径仍以显式高斯投影、深度排序和 alpha 合成为基础，只是在采样、属性预测、局部编辑、动态变形、压缩编码或下游任务接口处加入额外约束。这样写入本章时，它与前后方法的关系是：相对原始 3DGS 和本章已有锚点方法，它可以理解为 3DGS extends\_or\_uses 宽召回分类阶段的关系判断，后续需要按论文方法部分复核。
\textbf{证据与边界：} 本轮召回分类尚未逐项核对实验表格；后续证据补全阶段需要读取论文源码/PDF或项目页确认数据集、指标和对比方法。 该条来自宽召回批次；源码下载受 arXiv e-print 网络限制，部分论文暂时只能依据标题、摘要和类别分类，需要在入可信书前做人工/LLM精读复核。 本书只采用本地 paper pack、数据库分类、源码/引用登记和论文摘要中可核验的信息；没有本地证据支撑的速度、排行榜或开源完整性不在正文中强化。\textbf{方法信息：} 代码：\url{https://github.com/donalleniii/sharpsplat}。


本章的方法表明，3DGS场景理解的关键不在于给每个高斯简单贴一个标签，而在于语义属性如何被可靠地提升、压缩、渲染和查询。Feature 3DGS提供了可渲染特征场的底座；Gaussian Grouping、SAGA和FlashSplat解决2D掩码到3D实例的提升；OpenGaussian、LangSplatV2、GOI和GAGS则围绕开放词汇特征的压缩、粒度和查询稳定性展开。往后看，真正可靠的系统往往会混合这些路线：用基础模型提供2D语义，用3DGS保证跨视角几何一致，用码本或稀疏表示解决高维语言成本，再用任务特定模块接入检测、位姿、变化检测或场景图。

\textbf{设计动机与边界：} 相对原始 3DGS 以独立高斯和图像重建误差为主的设定，SharpSplat 把问题落在以下可实现的改动上：从摘要层面看，SharpSplat 围绕“把 3DGS 接入语义、语言或场景理解任务”展开。若其正文或代码后续可读，应进一步核对它具体改动的是表示、初始化、渲染、优化、损失、动态建模、压缩或应用接口中的哪一环。其设计目的不是单独增加一个模块，而是用上述表示、监督或渲染约束消除原流程在目标任务中的具体失配。该工作被召回是因为论文文本包含 Gaussian/3D Gaussian 相关表述；当前分类先确定其任务域和与 3DGS 的相关性，后续证据流水线再补代码、企业和引用信息。适用时还应保留论文给出的边界：该条来自宽召回批次；源码下载受 arXiv e-print 网络限制，部分论文暂时只能依据标题、摘要和类别分类，需要在入可信书前做人工/LLM精读复核。
\textbf{尾项实现澄清：} SharpSplat 的出发点是原始 3DGS 在无人机建筑立面上容易把锐利结构转为平滑漂浮几何。它保持标准高斯渲染管线，用 SAM 3 建筑掩码提取建筑学显著边缘，并在训练时将渲染图像梯度同这些边缘对齐；该边缘正则迫使高斯收敛到清晰的屋顶、立面和轮廓过渡，而非改动表示或光栅化器。
